\chapter{C1: ordinary and small-source charts}
\label{ch:cert-3}
Formula tags are local to their certificate section. All source and testing fields remain separately labelled.

\section{\texorpdfstring{The $q$-infinitesimal C1.1 source at both small testing angles}{The q-infinitesimal C1.1 source at both small testing angles}}
\label{component:pc376284fc7ae}

\label{pc376284fc7ae:sec:r7-c1-q-inf-local}
\begin{theorem}
On the original C1.1 source with $q=k^2\eta$ infinitesimal, the entire
original parent is tensor SOS whenever both independent testing
half-angles $u,v$ are infinitesimal. All original opening rates and
$\beta=0$ are retained. The boundary $D=0$ is empty here because
$D_b$ has residue $1/2$.
\end{theorem}
This statement includes every state; it does not yet assert the
ordinary or unbounded $v$ testing charts.

Use the preceding source notation and put $c=\sqrt q$, $e=c/g$,
$h=H/R$, and $E_*=h-v$. Work first with $Z=1$ and $u,v,h>0$.
Define
\[
 A=ch,\quad B=cL/R,\quad C=uv/b,\quad D=qg,\quad
 \Gamma=(ABCD)^{1/4},\quad d_i=A_i-\Gamma.
\]
Then
\[
 \Gamma^2=qhv,\qquad BC=e v^2h,\qquad
 u h=uE_*+bC.
 \tag{QI1}
\]
Here $h=u(W+uY)/(Rb^2)$ is separated into a positive pure-source
factor and a positive pure-testing factor. Every radical below is
therefore legal; no mixed inverse or order comparison is used.
The exact original inequality $\beta\ge0$ gives $t=O(e^2)$, while
$t$ is bounded. Hence $t/e$ and $\sqrt{t/e}$ are bounded.

\paragraph{Exact scalar, with the $D$ packet restored.}
Write $U=u^2,V=v^2$, $D_u=1+U,D_v=1+V$, and
$a_t=1-\rho_0+(t/R)V$. At both state normals zero the entire original
scalar is
\[
 a h^2+b_*hv+c_*v^2+P_L L^2+2D_uC^2+4\alpha_b g^2D_u,
 \tag{QI2}
\]
where
\begin{align*}
 a={}&4D_vR^2\{c_b(1-\rho_0)^2+a_b\rho_0^2/D_u\}
 +2b^2R^2\{a_t^2+V\rho_0^2/D_u\}\\
 &+4\alpha_b g^2b^4VR^2\rho_0^2/D_u
 +4D_bb^4R^2a_t^2,\\
 b_*={}&-4R\{D_u(1+tV/R)-U\rho_0\}
 -8\alpha_b g^2b^2R\rho_0-8D_bb^2Ra_tD_u,\\
 c_*={}&4D_bD_u^2.
\end{align*}
All these rational functions have bounded source coefficient arrays;
only $D_u$ and positive pure-source units occur in their denominators.
The $r$ openings cancel exactly as in (SC12).
In particular
\[
 D_b=\frac{N_D}{2JP_d},\qquad
 N_D=\frac{K(1+qg^2)(1+b^2r)-JP_d}{b^2}
\]
has a polynomial numerator: the numerator before division vanishes
identically at $b^2=0$, where $P_d=1+qg^2$. This is a finite monic
division, so it introduces no unbounded source coefficient array.
Also $\alpha_b/q^2$, $h_2/q$, and $h_b/q$ have explicit bounded
arrays from their source formulas. These observations justify the
$q$ coquotients below uniformly through both source openings.

Use formal specialization at $q=0$ only to construct finite polynomial
coquotients; it is not a new physical source. Put $z=b^2$,
$J_0=1-z(1-r)$, and $S=1-r+tV$. The three specialized coefficients are
\begin{align*}
 a_0={}&2D_v\{(1+zr)(1-r)+r/D_u\}
       +2z\{S^2+Vr^2/D_u\}+2z^2S^2/J_0,\\
 b_0={}&-4\{D_u(1+tV)-Ur\}-4zSD_u/J_0,\\
 c_0={}&2D_u^2/J_0.
\end{align*}
Define $\widetilde b_0=b_0+4tD_uV$, and the bounded finite
coquotients
\[
 X=a_0+\widetilde b_0+c_0=U X_U+V X_V,
 \qquad Y=-\widetilde b_0-2c_0=U Y_U+V Y_V.
 \tag{QI3}
\]
For definiteness, first subtract the value at $U=0$, divide by $U$,
and then divide that value by $V$. Their two $V$ coquotients are
\[
 X_V=2+2zr(1-2t)+\frac{2z(1+zr)t^2V}{J_0},
 \qquad Y_V=\frac{4zt}{J_0}.
 \tag{QI4}
\]
Thus $X_V=2+\varepsilon_V$, $\varepsilon_V\in B_{\rm Inf}$;
$X_U,Y_U$ have bounded arrays. Likewise write
$a=a_0+qa_1$, $b_*=b_0+qb_1$, $c_*=c_0+qc_1$ by finite numerator
division. All $a_1,b_1,c_1$ have bounded arrays. Direct source
specialization gives
\[
 \overline{a_1-c_1}=2,\qquad
 \overline{b_1+2c_1}=-8,\qquad
 \overline{R^2P_L/q}=2,\qquad
 \overline{4\alpha_bD_u/q^2}=2.
 \tag{QI5}
\]

The following is the exact scalar identity to be allocated:
\begin{align*}
 (\mathrm{QI2})={}&(c_0+qc_1)E_*^2+2D_uC^2
 +X_U(uE_*+bC)^2+Y_Uu(uE_*+bC)E_*\\
 &+X_V(vh)^2+Y_Vv(vh)E_*
 -4(t v/e)D_uBC\\
 &+(a_1-c_1)A^2+(b_1+2c_1)\Gamma^2
 +(R^2P_L/q)B^2+(4\alpha_bD_u/q^2)D^2.
 \tag{QI6}
\end{align*}
It uses only (QI1) and $qv^2=2\Gamma^2-A^2+qE_*^2$.
In particular the apparently dangerous negative $t$ term is already
an infinitesimal-coefficient $BC$ term, since $v$ is infinitesimal.

\paragraph{Two extra contact squares and their finite first jets.}
Set
\[
 F_*=vh-gD,\quad
 R_F=\frac{\sqrt{g\Gamma}}{c}(\Gamma-D),\quad
 V_0=cv,\quad R_E=\sqrt\Gamma(\sqrt A-\sqrt{V_0}).
\]
All four roots vanish at actual contact. Exact multiplication gives
\begin{align*}
 (vh)^2&=F_*^2+2R_F^2+2gF_*d_4+g^2D^2-4g^2\Gamma d_4,\\
 c\Gamma E_*&=2\Gamma d_1-R_E^2,\\
 \Gamma F_*&=-2g\Gamma d_4+R_F^2/g.
 \tag{QI7}
\end{align*}
The root $R_E$ receives a strict real square pivot from $E_*^2$:
\[
 qE_*^2=(A-V_0)^2
       =4R_E^2+(\sqrt A-\sqrt{V_0})^4.
 \tag{QI8}
\]
Reserve this fraction $q$ of the strict $E_*$ pivot; the remaining
coefficient still has positive real residue. The leading $2(vh)^2$
in (QI6), by (QI7), gives strict $F_*$ and $R_F$ pivots, with only
infinitesimal quadratic corrections and the printed first-jet term.
The remaining leading angular form is $2(\sum A_i^2-4\Gamma^2)$,
with the strict finite toric square list from (SC16).

For precision, the only additional bare-$\Gamma$ normal couplings
left by (QI6), apart from ordinary toric first jets, have the form
$\Gamma E_*a_*+\Gamma F_*b_*'$ with
\[
 a_*=bu(2X_U+Y_U)+Y_Vvg,
 \qquad b_*'=2\varepsilon_Vg.
\]
At contact $u=bc,v=cg$ these satisfy
\[
 a_{*0}/c=b^2(2X_{U0}+Y_{U0})+Y_{V0}g^2\in B_{\rm Inf},
 \qquad b_{*0}'/g=2\varepsilon_{V0}\in B_{\rm Inf}.
\]
Their differences from contact have finite bounded monic Taylor
arrays. The conductors (SC18) allocate those differences against
$E_*,F_*$ with infinitesimal coefficients. Apply (QI7) to their
constant terms. What remains is a bounded infinitesimal toric
quadratic plus $\Gamma\sum\mu_i d_i+\Gamma^2\sigma$.
The exact scalar value and three original first jets imply its
contact value and first-jet balances. The same two finite divisions
as in (SC17)--(SC20), or the Hadamard contact-lift identity, allocate
it on the strict toric root list. No infinite series or unsupported
ideal-membership assertion is used: (QI3), the $q$ coquotients,
and the two displayed contact coefficient ratios are the finite
certificates needed for the additional roots.
Every coefficient function in these Taylor steps depends only on
$u,v$. Thus only the first two conductors in (SC18)--(SC20) are
used; the third state-coordinate jet supplies the inactive Hadamard
balance. No bound on $(W+uY)/R$ after $Z$ normalization is assumed.

\paragraph{Both complete state normals.}
Use $\widehat N=\sqrt{H_0}N$, $H_0=4h_2/R+4h_b/t$, as in the
preceding uniform-normal lemma. Add back the centered $D$ packet
from (SC10), including its normals. Its scalar root is exactly
\[
 J_*=(T+b^2tV)h-D_uv
 =E_*+b^2c(1-tg^2)A-buC+b^2tvF_*+b^2tvgD.
 \tag{QI9}
\]
Its full root is $J_*+ubvgN-b^2N_y/\sqrt s$.
All additional hard coefficients are infinitesimal after the
$\widehat N$ normalization.

For the first $N$ normal in (SC14), its $H$ term is a bounded
infinitesimal multiple of $u h=uE_*+bC$; its other terms are
infinitesimal multiples of $B,C$. For the harmonic normal, multiply
its $H$ coefficient by $\sqrt s$. After the exact opening
cancellations, its value at $q=U=V=0$ is zero. Finite numerator
coquotients therefore give
\[
 \sqrt s\,[H\text{-coefficient in }\mathcal L_y]
       =qK_q+UK_U+VK_V
 \tag{QI10}
\]
with bounded arrays. Use
\[
 qH=RcA,\quad u^2H=Ru(uE_*+bC),\quad
 v^2H=Rv(F_*+gD).
\]
The two remaining terms of $\mathcal L_y$ have infinitesimal
coefficients on $C,D$ since $\alpha_b/q^2$ is bounded. Thus both
complete first normals, including (QI9), are infinitesimal
combinations of $E_*,F_*,A,B,C,D$. Their exact contact zeros make
the coefficient sum of $A,B,C,D$ zero at contact. One finite monic
division and (SC18) allocate these on the same strict square list.
The hard block has residue $\operatorname{diag}(1,4)$ and all its
other entries have zero residue. This proves a bounded positive
residue Gram for the entire original parent, after restoring all
retained positive packets.

\paragraph{States and signs.}
The restriction $Z=0$ is positive by the original pure-testing state
restriction. For $Z\ne0$, divide the state by the pure-testing $Z$;
its size is immaterial, since every coefficient array above depends
only on $u,v$. On the negative-core chart, exact sign gauges give
$u,v>0$ and the inherited pure-testing Schur identity forces
$W+uY>0$. Hence $h>0$, as required. The other core-sign chart and
literal testing axes are direct original exits. Complex states
split into real and imaginary parts. This completes the asserted
small-$u$, small-$v$ theorem.


\section{\texorpdfstring{A uniform hard normal for the $q$-infinitesimal branch}{A uniform hard normal for the q-infinitesimal branch}}
\label{component:p7d43f3fccb7d}

This intermediate lemma by itself is not an admission of the full
$q$-infinitesimal C1.1 source. The following two completion sections
supply the scalar allocation and remote joins. Keep (SC1)--(SC3), but do
not discard the whole second paired packet. Its addition to the
$N^2$ coefficient in (SC3) gives exactly
\[
 \frac4R(h_2+h_bD_v)+\frac{4h_bM}{tR}D_v
       =H_0+\frac{4h_b}{t}v^2,
 \qquad H_0=\frac{4h_2}{R}+\frac{4h_b}{t}.
\]
Retain the last summand as the positive square
$4h_b(vN)^2/t$, and use the pure-source normalization
\[
                         \widehat N=\sqrt{H_0}\,N.
\]
The exact original inequality $\beta\ge0$ gives
$h_b/t\ge g^2/2$. The explicit formula for $h_2$ gives
$h_2/q$ a positive bounded unit, uniformly through both openings.
Consequently $H_0\ge c_0(q+g^2)$ for some positive real $c_0$.
This argument also permits $\beta=0$ literally. It does not require
$e=\sqrt q/g$ or its reciprocal to be bounded.

The coefficient of $\widehat N^2$ from this source part is exactly
one. Its remaining contribution $2g^2u^2v^2/H_0$ is bounded and
infinitesimal on the small-angle testing chart. Moreover
\[
 \frac{g}{\sqrt{H_0}},\quad
 \frac{\sqrt q}{\sqrt{H_0}}
\]
are bounded. Since $h_b-h_2=O(qg^2)$, the complete first normal
(SC14), after this normalization, is an infinitesimal-coefficient
linear combination of $B,H,V$ for $u,v$ infinitesimal. The first
coefficient is bounded because
\[
 \frac{h_bD_v-h_2}{\sqrt{qH_0}}
            =O(g^2+v^2),
\]
and the other two coefficients are bounded by constant multiples of
$buv$ and $v$. The hard cross coefficient with $N_y$ is likewise
infinitesimal.

The issue left by this lemma is scalar and angular: $H$ is unweighted in this
statement, whereas the four-term circuit controls $\sqrt q H$.
One must retain a strict $D$-based contact normal and allocate every
unweighted $H$ coefficient on that normal and the circuit. The
naive full extraction (SC10) fails in this branch: its raw $P_H$
satisfies the exact formal specialization
\[
       P_H\big|_{q=0,\,v=0}=-\frac{2ru^2}{1+u^2},
\]
so $P_H/q$ need not be bounded. Thus the corrected hard-normal
normalization removes the former $g/\sqrt q$ obstruction. The later
protected-square lift establishes the missing scalar allocation.


\section{\texorpdfstring{The other Bob testing charts for $q$ infinitesimal}{The other Bob testing charts for q infinitesimal}}
\label{component:pd2ecc6ac1891}

\label{pd2ecc6ac1891:sec:r7-c1-q-inf-remote}
\begin{theorem}
The original C1.1 source with $q=k^2\eta$ infinitesimal is tensor SOS
at every independent test. Together with the $q$-unit theorem this
closes the entire original C1.1 source.
\end{theorem}
The inherited $C,A$-large admission already handles $u$ not
infinitesimal. The preceding local theorem handles $u,v$
infinitesimal. Here are the two remaining $v$ charts. Every positive
opening rate and the literal permitted coefficient boundary is kept.

\paragraph{Ordinary $v$: the full normal residue is strict.}
Let $v$ be a positive bounded unit, use $Z=1$ on the negative-core
chart, and retain the roots of (QI1). Normalize $N$ by $H_0$ as
before, but now restore the positive packet $4h_b(vN)^2/t$.
Write
\[
 \tau=\operatorname{res}t,\qquad
 \rho=\operatorname{res}(t/e^2),\qquad
 a^2=\frac{\tau}{2(1+\tau)},\quad
 \gamma^2=\frac{\rho}{2(1+\tau)}.
\]
The original $\beta\ge0$ equation implies $0\le\rho\le1$;
$0\le\tau\le1$ follows from $M=R-t>0$ and $R\sim1$.
The entire residue on $(\widehat N,B,C)$ has diagonal entries
\[
 1+\frac{v^2}{1+\tau},\qquad 2(1+\tau v^2),\qquad 2,
\]
and polynomial cross coefficients
\[
 4av^2,\qquad -4\gamma v,\qquad -4\sqrt{\tau\rho}\,v.
\]
Its exact square decomposition is
\begin{align*}
 &2(C-\gamma v\widehat N-\sqrt{\tau\rho}\,vB)^2
       +\widehat N^2+2B^2\\
 &\qquad +(1-\rho)v^2
    \left(\frac{\widehat N}{\sqrt{1+\tau}}+
                         \sqrt{2\tau}\,B\right)^2.
 \tag{QR1}
\end{align*}
Thus it is strictly positive definite, also at $\rho=1$, the
possible $\beta=0$ endpoint, and at $\tau=0$ or $1$.
The other residue blocks are exactly
\[
                       2E_*^2+2(vh)^2+4D_vN_y^2.
\]
There are no nonzero leading cross entries with these blocks.

For completeness, the boundedness check uses the entire raw scalar
(QI2) and both full normal formulas (SC14), with (QI9) restored.
Terms containing $u h$ use $u h=uE_*+bC$; all their coefficients
are infinitesimal. The $N$--$B$ and $N$--$C$ coefficients with
nonzero residue are precisely those retained in (QR1). The harmonic
normal has infinitesimal coefficient on $h$ for bounded $v$:
after the opening cancellation its coefficient lies in
$(b^2,q,u^2)$ with bounded arrays on this chart. Its other
coefficients are infinitesimal as well.
All remaining source-constant terms and weak monomials enter the
same strict Gram through the exact bounded pure-testing identities
\[
                   h=(vh)/v,\qquad
                   1=((vh)-vE_*)/v^2.
\]
For example $A=ch$ and $D=qg$ are infinitesimal-coefficient
combinations of these two roots. Consequently the actual full Gram
has the positive residue just displayed. Restore all retained
packets. The same sign, $Z=0$, and complex-state joins as in the
local theorem apply; no restriction on the size of the normalized
state was used.

\paragraph{Unbounded $v$: a five-root endpoint frame.}
Put $w=1/v$, so $w$ is infinitesimal, and multiply the original
parent by $w^2$. The exact source formulas give positive real
constants $\mu,d_0$ with
\[
 C/d\ge\mu/b^2,\qquad A/d\ge\mu/b^2,\qquad
 D/d\ge d_0 b^2.
\]
For instance sufficiently small fixed constants below $1/4$ work,
since $b^2C/d,b^2A/d$ are uniformly bounded below and
$D/(db^2)\sim1/2$. Retain all excesses and the entire $m,\alpha,
\beta$ packets. No assumption on $e$ is needed.
Use the five roots
\[
              P=uW/b,\quad Q=Y/b,\quad X,\quad J_b'=bZ,
                     \quad V_b=uwZ.
\]
The homogenized signed core is exactly
\begin{align*}
 K={}&V_b^2-2buV_bQ-2wPJ_b'+2wV_bX+b^2w^2P^2
       +2buwPX-2b^2w^2PQ\\
 &+(D_u+w^2)(X^2+b^2Q^2)-2buwXQ.
 \tag{QR2}
\end{align*}
The three selected event packets are
\[
 4\mu D_w(P-Q)^2+4\mu D_uD_wQ^2,
 \qquad
 4d_0(b^2wP+buX-b^2wQ-J_b')^2.
 \tag{QR3}
\]
Combining these with $2K$, the bounded Gram has strictly positive
residue
\[
 4\mu\{(P-Q)^2+Q^2\}+4d_0(J_b')^2+2X^2+2V_b^2.
\]
All other displayed coefficients have zero residue. This is a full
state-quadratic certificate, not a restriction to a contact graph.
It includes the homogeneous endpoint $w=0$ directly and uses no
inverse involving independent source and testing variables.

The infinitesimal, unit, and unbounded $v$ charts exhaust all tests
remaining after the inherited $C,A$ admission. Hence the entire
$q$-infinitesimal C1.1 source, and consequently C1.1 itself, is closed.


\section{\texorpdfstring{C1.1 with $k^2\eta$ a unit: a complete contact certificate}{C1.1 with k2eta a unit: a complete contact certificate}}
\label{component:p95350f69a353}

\label{p95350f69a353:sec:r7-c1-q-unit}
\begin{theorem}
On the original C1.1 source, suppose $q=k^2\eta$ is a positive bounded
unit. The entire original parent is tensor SOS at every independent
projective test and every real or complex state. All positive opening
rates and the permitted literal boundary $D=0$ are included.
\end{theorem}
The proof uses the exact full quadratic (SC3), not its leading scalar.
The $q$-infinitesimal part of C1.1 is not included in this theorem.

\paragraph{One more centered packet.}
Use the variables of (SC1)--(SC5), and put $s_D=J/R$. Retain the
additional nonnegative square
\[
                  4D_b(J_b+s_DH)^2.                 \tag{SC10}
\]
It vanishes at actual contact, where $H=RvZ$ and $J_b=-JvZ$.
Its removal replaces $4D_bJ_b^2$ in (SC3) exactly by
\[
                  -8D_bs_DHJ_b-4D_bs_D^2H^2.
\]
In particular it removes the entire standalone $(vZ)^2$ term.
This step allows $D_b=0$ literally and does not divide by $D_b$.

Set $Z=1$ on the remaining state chart. Put $c=\sqrt q$ and
$w=(W+uY)/R$, and use the four separated positive monomials
\[
 A=\frac{cuw}{b^2},\quad B=\frac{cvw}{bg},\quad
 C=\frac{uv}{b},\quad D=c^2g,\quad
 \Gamma=(ABCD)^{1/4}.
 \tag{SC11}
\]
Their exact contact is $(u,v,w)=(k,cg,bg)$, with $k=bc$, and
$A=B=C=D$ there. Notice that $H=RA/c$, $L=RB/c$, $V=C$,
$g=D/q$, and $Hv=R\Gamma^2/q$.

Here are explicit formulas for the entire scalar at $N=N_y=0$,
after (SC10). Write $a_t=1-\rho_0+pv^2$. Define
\begin{align*}
 P_H={}&4D_v\{c_b(1-\rho_0)^2+a_b\rho_0^2/D_u\}
       +2b^2\{a_t^2+v^2\rho_0^2/D_u\}\\
 &+4\alpha_bg^2b^4v^2\rho_0^2/D_u
       -8D_bs_Db^2a_t-4D_bs_D^2,\\
 P_L={}&4[h_2(1-p)+h_bpD_v]/R,\\
 P_{Hv}={}&-4[D_u(1+pv^2)-u^2\rho_0]
             -8\alpha_bg^2b^2\rho_0+8D_bs_DD_u.
\end{align*}
The exact scalar is
\[
 f=f_AA^2+f_BB^2+f_CC^2+f_DD^2-f_\Gamma\Gamma^2,
 \tag{SC12}
\]
where
\[
 f_A=R^2P_H/q,\quad f_B=R^2P_L/q,\quad
 f_C=2D_u,\quad f_D=4\alpha_bD_u/q^2,
 \quad f_\Gamma=-RP_{Hv}/q.
\]
These are rational functions of $u,v$ over bounded source coefficients,
with only positive-unit denominators. For example,
\[
 c_b(1-\rho_0)^2=
   \frac{K(1+b^2r)E(1-r)}{2P_dR^2},\qquad
 a_b\rho_0^2=\frac{P_4K^2r}{2P_dR^2}.
\]
Thus neither $r$ opening survives in their denominators.
For $u,v,b,g$ infinitesimal their residues are
\[
               \overline f_A=\overline f_B=
 \overline f_C=\overline f_D=2,\qquad
               \overline f_\Gamma=8.               \tag{SC13}
\]
This uses the exact source residues in (SC2), including
$\overline D_b=(1-2\overline q)/2$. It does not assume a positive
lower bound on $D_b$.

Here is an explicit finite denominator certificate used below.
Every function just displayed is $P(u,v)/(1+u^2)$ for a polynomial
$P$ with bounded pure-source coefficient array. For a monomial use
\[
 x^m-x_0^m=(x-x_0)\sum_{j=0}^{m-1}x^{m-1-j}x_0^j.
\]
Applying this identity to each monomial supplies finite first and
second difference arrays with bounded source coefficients. For the
denominator use only the exact identity
\[
 \frac1{1+u^2}-\frac1{1+k^2}
 =-\frac{(u-k)(u+k)}{(1+u^2)(1+k^2)}.
\]
Thus the first and second finite Taylor arrays about $(k,cg)$
are bounded and use only pure-testing $1+u^2$ and pure-source
$1+k^2$ denominators, both units. The same construction after
$u=U^2,v=V^2$ gives the half-variable arrays below. This is a
finite polynomial identity construction, not an analytic series.

\paragraph{Every first normal and the complete hard block.}
Let $\delta_\rho=\rho_0-\rho$. The full first-normal coefficients
of the remaining quadratic, with no normal suppressed, are
\begin{align*}
 \mathcal L_N={}&\frac8R(h_bD_v-h_2)L
       +4buvg(a_t-2D_bs_D)H-4D_uvgV,\\
 \mathcal L_y={}&8D_v\sqrt{s}\,\delta_\rho H\\
 &+\frac H{\sqrt s}\left[
 -\frac{8D_va_bu^2\rho_0}{D_u}+8D_bs_Db^2
 +\frac{8\alpha_bg^2b^4v^2\rho_0}{D_u}
 -4b^2a_t+\frac{4b^2v^2\rho_0}{D_u}\right]\\
 &+\frac{4bu}{\sqrt{s}}V
       -\frac{8\alpha_bgb^2v}{\sqrt{s}}(gZ).
                                                        \tag{SC14}
\end{align*}
Each is a bounded infinitesimal-coefficient linear combination of
$A,B,C,D$. At actual contact it is zero exactly, since the source
is a stationary zero and every removed packet vanishes there.
The coefficients of $N^2,N_y^2,NN_y$ are respectively
\[
\begin{aligned}
 H_{NN}&=4(h_2+h_bD_v)/R+2g^2u^2v^2,\\
 H_{yy}&=4D_v\left(1-\frac{a_bu^2}{D_us}\right)
       +\frac{4\alpha_bg^2b^4v^2}{D_us}
       +\frac{2b^2}{s}\left(1+\frac{v^2}{D_u}\right),\\
 H_{Ny}&=-4buvg/\sqrt{s}.
\end{aligned}                                                     \tag{SC15}
\]
Thus the hard Gram has residue $\operatorname{diag}(4q,4)$.
In (SC15), $H_{Ny}$ denotes the polynomial coefficient of $NN_y$,
not the off-diagonal matrix entry.

\paragraph{A finite geometric-mean square lift.}
Here is a detailed lift for the particular raw module (SC12).
It is not a claim about arbitrary contact-vanishing polynomials.
Define $\alpha_i=A_i-\Gamma$, with
$(A_1,A_2,A_3,A_4)=(A,B,C,D)$, and put
\[
 T_4=\sum_iA_i^2-4\Gamma^2.
\]
A fixed positive Gram on the following finite roots represents $T_4$:
\[
 \alpha_i,\qquad
 Q_{ij|kl}=\sqrt{A_iA_j}-\sqrt{A_kA_l},\qquad
 \sqrt\Gamma(\sqrt{A_i}-\sqrt{A_j}),\qquad
 \sqrt\Gamma((A_iA_j)^{1/4}-(A_kA_l)^{1/4}).          \tag{SC16}
\]
Use all four $\alpha_i$, all three pair partitions, and all six
unordered differences in the third group. To verify strictness, average
positive multiples of the symmetric circuit decomposition in (SC9) and
\[
 T_4=\sum_i\alpha_i^2+2\Gamma\left(\sum_iA_i-4\Gamma\right).
\]

over the three decompositions
\[
 \sum_iA_i-4\Gamma=(\sqrt{A_i}-\sqrt{A_j})^2
 +(\sqrt{A_k}-\sqrt{A_l})^2
 +2((A_iA_j)^{1/4}-(A_kA_l)^{1/4})^2.
\]
The additional differences $A_i-A_j$ in the first decomposition are
$\alpha_i-\alpha_j$ and give nonnegative corrections to the already
strict $\alpha$ block.

Write $\epsilon_i=f_i-2$, $\epsilon_\Gamma=f_\Gamma-8$, and
$\sigma=\sum_i\epsilon_i-\epsilon_\Gamma$. Then
\[
 f-2T_4=\sum_i\epsilon_i\alpha_i^2
       +2\Gamma\sum_i\epsilon_i\alpha_i+\Gamma^2\sigma.
 \tag{SC17}
\]
All $\epsilon$ have zero residue. The exact contact value implies
$\sigma(k,cg)=0$. Clear the pure-testing denominator $D_u$ if needed;
monic division gives
$\sigma=(u-k)\sigma_u+(v-cg)\sigma_v$ with bounded coefficients.
For the monomials (SC11), the exact conductors are
\[
\begin{aligned}
 \Gamma(u-k)&=\sqrt{uk}\,Q_{13|24},\\
 \Gamma(v-cg)&=\sqrt{vcg}\,Q_{23|14},\\
 \Gamma(w-bg)&=\sqrt{wbg}\,Q_{12|34}.
\end{aligned}                                                     \tag{SC18}
\]
Thus the last term in (SC17) becomes $\Gamma$ times a linear
combination of the $Q$ roots with infinitesimal bounded coefficients.
The exact identity
\[
 Q_{ij|kl}-\tfrac12(A_i+A_j-A_k-A_l)
 =-\tfrac12(\sqrt{A_i}-\sqrt{A_j})^2
  +\tfrac12(\sqrt{A_k}-\sqrt{A_l})^2               \tag{SC19}
\]
changes these terms into infinitesimal quadratic corrections on (SC16)
and $\Gamma\sum_i\mu_i\alpha_i$, where each $\mu_i$ is infinitesimal
and is rational in $\sqrt u,\sqrt v$ with bounded coefficient arrays
and positive-unit denominators.

The exact three contact jets imply that the four source values
$\mu_i(k,cg)$ coincide. Indeed the differential vectors of
$\alpha_i/\Gamma$ at contact are the four rows
\[
 (1/2,-1/2,1/2),\quad(-1/2,1/2,1/2),\quad
 (1/2,1/2,-1/2),\quad(-1/2,-1/2,-1/2)
\]
in logarithmic coordinates $(u/k,v/(cg),w/(bg))$. Their sole linear
relation is the sum of all four rows. Cancelling the nonzero source
factors in this equality establishes coincidence; no inverse of a
small number enters a Gram coefficient.
Call the common infinitesimal value $\mu_*$. The common term is
\[
 \mu_*\Gamma\sum_i\alpha_i
   =\mu_*\Gamma\left(\sum_iA_i-4\Gamma\right),
\]
an infinitesimal correction on the explicitly displayed variance roots.
For each remaining $\mu_i-\mu_*$, monic division in $\sqrt u,\sqrt v$
has bounded coquotients. Apply the half-variable conductors
\[
\begin{aligned}
 \Gamma(\sqrt u-\sqrt k)
   &=(uk)^{1/4}\sqrt\Gamma
          ((AC)^{1/4}-(BD)^{1/4}),\\
 \Gamma(\sqrt v-\sqrt{cg})
   &=(vcg)^{1/4}\sqrt\Gamma
          ((BC)^{1/4}-(AD)^{1/4}).
\end{aligned}                                                     \tag{SC20}
\]
Their coefficient factors are infinitesimal. Multiplying by the
remaining $\alpha_i$ proves that the entire scalar correction in
(SC17), not merely its low-order jet, is a bounded zero-residue
quadratic form on (SC16). This uses only two finite monic divisions.
The boundedness of their coefficient arrays follows from the explicit
rational functions (SC12), whose test denominators are powers of
$D_u$, and from the displayed square-root substitutions; bounded
function values alone are not being used as a substitute for this fact.

The same argument for a first normal is shorter. If
$\mathcal L=\sum_i a_iA_i$, write
$\mathcal L=\sum_i a_i\alpha_i+\Gamma\sum_i a_i$.
Its exact contact zero gives $\sum_i a_i(k,cg)=0$; one monic division
and (SC18) allocate it on (SC16), with infinitesimal coefficients.
Apply this to both complete expressions (SC14).
Together with (SC15), this yields a bounded strictly positive-residue
Gram for the entire original remainder. Restore (SC10), the paired
packets, and all positive pure-source and pure-testing normalizations.

\paragraph{Exhaustion and legal radicals.}
For $q$ a unit, $e^2=q/g^2$ is unbounded. The exact adjacent-event
rotations and the $W$-unit state argument (SC8) therefore reduce every
unadmitted negative-core test to $u,v,W$ infinitesimal and $Z$ a unit.
Literal axes are direct original-core exits. Use exact sign gauges and
normalize $Z=1$. The inherited pure-testing Schur identity forces
$W+uY>0$ on the remaining negative-core chart: with
$w_t=(W+uY)/D_u$, the eliminated scalar is
\[
 \frac{uvD_uD_v}{\mathcal D}
       (D_uuv w_t^2-2D_uD_vw_t+uv),\qquad
 \mathcal D=u^2v^4+u^2v^2+u^2+v^4+2v^2+1>0,
\]
which is nonnegative when $w_t\le0$; its eliminated two-square
matrix is strictly positive. Hence $w=(W+uY)/R>0$.
Every monomial in (SC11), and every product under a radical in
(SC16)--(SC20), factors into a positive pure-source quantity and a
positive pure-testing quantity. The occurrence of $R$ in $w$ is a
positive pure-source unit, not a mixed denominator. All test
coefficients remain bounded because $u,v,w$ are infinitesimal.
If the general toric contact lemma is used, its third coordinate
is literally the pure-testing $W+uY$, with pure-source contact
$Rbg$; this direction is inactive in every coefficient function.
Complex state evaluations split into real and imaginary parts.
This supplies every independent test and proves the theorem.



\section{\texorpdfstring{C1.1: an opening-uniform full normal and a new angular circuit}{C1.1: an opening-uniform full normal and a new angular circuit}}
\label{component:p9afbe73899af}

\label{p9afbe73899af:sec:r7-c1-small}
This section concerns the original C1.1 source, not an enlarged family.
All coefficients are divided by $d>0$. Put
\[
 b=\eta^{-1/2},\quad g=\sqrt\zeta,\quad q=k^2\eta,
 \qquad x=b^2q,\quad a=b^2(1-r).
 \tag{SC1}
\]
Here $b,g$ are positive infinitesimals, $q>0$ is bounded, and
$\overline q\le1/2$. The new bounded two-normal reduction below is
asserted for $q$ a positive unit. The exact polynomial identities
remain valid more generally, but their indicated boundedness argument
would not control $g/\sqrt q$ when $q$ is infinitesimal.
No all-tests closure is inferred from the reduction.

Keep $K,R,T,M,J,E,Q,P_4,P_d,P_\alpha$ from (NC1)--(NC2), and set
\[
\begin{aligned}
 c_b&=b^2C/d=\frac{K(1+b^2r)E}{2(1-r)P_d},&
 a_b&=b^2A/d=\frac{P_4}{2rP_d},\\
 D_b&=D/(db^2),&\alpha_b&=\alpha/(db^2g^2),\\
 h_2&=\frac{Kq(1+b^2r)(1-tg^2)(1+qg^2)}{2P_d},&
 h_b&=\frac{Kq(1+b^2r)[1+(1-t)g^2]}{2P_d}.
\end{aligned}
\]
Thus $m/d=h_2/(g^2M)$ and $\beta/d+1/2=h_b/(g^2t)$.
The exact relation
\[
 h_b-h_2=\frac{Kq(1+b^2r)E g^2}{2P_d}
 \tag{SC2}
\]
removes the apparent opening dependence from the first normal of the
paired $m,\beta$ currents. The quantities $D_b,\alpha_b,h_2,h_b$ are
bounded; $\alpha_b,h_2,h_b$ are units in the stated $q$-unit domain.
Their residues are respectively $(1-2q)/2,q^2/2,q/2,q/2$, with
bounded source letters on the right understood at residue.

\paragraph{Exact full-parent reduction.}
For bounded affine tests define
\[
 Q_b=(W+uY)/b,\quad Y_b=Y/b^2,\quad p=t/R,\quad
 N=\frac{(X+uZ)/b-pvQ_b}{g},
\]
\[
 H=uQ_b/b,\quad L=vQ_b/g,\quad V=uvZ/b,\quad
 B_b=gN+pvQ_b,\quad D_u=1+u^2,\quad D_v=1+v^2.
\]
All roots use only pure-source inverses. Retain the two nonnegative
original packets
\[
 \frac{4h_2t}{MR}N^2,
 \qquad \frac{4h_bM}{tR}D_vN^2.
\]
The complete original parent divided by $db^2$, minus these packets,
is exactly
\begin{align*}
 G={}&\frac4R\{(h_2+h_bD_v)N^2+2(h_bD_v-h_2)LN
              +[h_2(1-p)+h_bpD_v]L^2\}\\
 &+4D_v\{c_b(H-D_uY_b)^2+a_bD_uY_b^2\}
   +4D_bJ_b^2+4\alpha_b g^2D_u(vb^2Y_b-Z)^2+2K_b,
                                                        \tag{SC3}\\
 J_b={}&b^2(H-D_uY_b)+ubvB_b-D_uvZ,\\
 K_b={}&b^2(H-D_uY_b)^2+2buv(H-D_uY_b)B_b+u^2v^2B_b^2\\
 &+D_uv^2b^2Y_b^2+D_uV^2
       -2D_uV(Q_b+vB_b)+2D_uu^2vZY_b.
\end{align*}
These are entire quadratic polynomials in the state; no restriction
of a state normal has been made. In particular, (SC3) is a direct
input for collecting both complete first normals and their full hard
quadratic, rather than an identity solely on a contact graph.

\paragraph{A harmonic normal removes both $r$ openings.}
Put
\[
 s=c_b+a_b,\quad \rho=c_b/s,\quad h=c_ba_b/s,
 \quad \rho_0=Kr/R,\qquad
 N_y=\sqrt{s}\,(D_uY_b-\rho_0H).
\]
Then the complete $C,A$ line in (SC3), divided by $4D_v$, equals
\[
 \left[N_y+\sqrt{s}(\rho_0-\rho)H\right]^2+hH^2
 -\frac{a_bu^2}{D_u}
           \left(\frac{N_y}{\sqrt{s}}+\rho_0H\right)^2.
 \tag{SC4}
\]
The inverse $1/D_u$ is pure testing, and every coefficient in (SC4)
is bounded uniformly through either $r\to0$ or $r\to1$.
Indeed $s$ is bounded below, $h$ is a positive unit, and
$a_b/s$, $a_b\rho_0/\sqrt{s}$ and $a_b\rho_0^2$ are bounded.
There is no mixed harmonic denominator.

Here is the exact source cancellation that also centers the normal
without an opening loss. Define
\[
 S_0=Kr(1+b^2r)E+(1-r)P_4,
\]
\[
 T_0=q^2(1-tg^2)^2+qrg^2+2qtg^2-qg^2-2q+r.
\]
Then
\[
 \rho=\frac{Kr(1+b^2r)E}{S_0},\qquad
 Kr-\rho R=-\frac{b^2r(1-r)K T_0}{S_0},\qquad
 h=\frac{K(1+b^2r)E P_4}{2P_dS_0}.             \tag{SC5}
\]
Both $S_0$ and the other denominators are positive units. Also
$s\asymp[r(1-r)]^{-1}$, so $\sqrt{s}(\rho_0-\rho)$ is
infinitesimal without a lower bound on either opening.
At the actual contact,
\[
 (u,v)=(b\sqrt q,\sqrt q g),\quad
 Q_b=RgZ,\quad Y_b=r\sqrt q gZ,
\]
and both $N,N_y$ vanish exactly. Thus (SC3)--(SC5) remove every
unbounded source weight while retaining the actual contact and every
state normal.

\paragraph{The leading angular circuit.}
For $u,v$ infinitesimal, the leading scalar after the harmonic state
normal has been removed is twice
\[
 (1-q)H^2+qL^2+V^2+(1-2q)(vZ)^2+q^2(gZ)^2-2VQ_b.
\]
Here the leading relation $H(vZ)=VQ_b$ is exact for the displayed
roots. It splits as
\[
 (1-2q)(H-vZ)^2+
 qH^2+qL^2+V^2+q^2(gZ)^2-4qVQ_b.              \tag{SC6}
\]
The final four positive terms form one AM--GM circuit, since
\[
 (\sqrt qH)(\sqrt qL)V(qgZ)=(qVQ_b)^2.
\]
On the positive negative-core chart its radical factors separate
between the source and testing fields. This proves positivity of the
leading scalar, including $q=1/2$, where its separate binomial drops
out. The leading hard normal coefficients are $4q$ on $N$ and $4$
on $N_y$. This organizing form alone is not yet a certificate for
all the higher terms in (SC3).

A useful exact conductor can be stated without comparing any source
and testing rate. Normalize $Z=1$ on its pure-testing state chart,
put $c=\sqrt q$, $w=(W+uY)/R$, and now use the contact-exact roots
\[
 H_0=uw/b^2,\ L_0=vw/(bg),\ V_0=uv/b,\ D_0=c^2g,
 \quad A_0=cH_0,\ B_0=cL_0,\ C_0=V_0.
\]
Their contact is $(u,v,w)=(bc,cg,bg)$. Set
\[
 E_0=H_0-v,\quad
 S_v=A_0-D_0-cE_0,\quad
 S_u=C_0-D_0-S_v,\quad S_w=B_0-D_0-S_v.
\]
Direct multiplication gives
\[
\begin{aligned}
 H_0(u-bc)&=(u-bc)E_0+bS_u,\\
 H_0(v-cg)&=(v-cg)E_0+(v/c)S_v,\\
 H_0(w-bg)&=(w-bg)E_0+(bg/c)S_w.
\end{aligned}                                                     \tag{SC7}
\]
Only the positive-unit inverse $1/c$ occurs. These identities lift
contact-vanishing coefficients times $H_0$ to the displayed root
module with bounded coefficients on the small pure-testing chart.
They do not alone prove a full scalar square lift: a first-jet
syzygy among the redundant angular roots still has to be allocated,
particularly when the coefficient of $E_0^2$ vanishes at $q=1/2$.

\paragraph{Exact small-cost angular pivots.}
The four-term circuit has the exact symmetric decomposition
\[
 \sum_{i=1}^4A_i^2-4\sqrt{A_1A_2A_3A_4}
 =\frac13\sum_{i<j}(A_i-A_j)^2
  +\frac23\sum_{\{ij\}\mid\{kl\}}
       (\sqrt{A_iA_j}-\sqrt{A_kA_l})^2,
\]
where the second sum is over the three pair partitions. Use
$(A_1,A_2,A_3,A_4)=(A_0,B_0,C_0,D_0)$ above and assume the positive
chart $u,v,w>0$. Then the following identities are exact:
\[
\begin{aligned}
 c\sqrt{uw}\,E_0
 &=\sqrt{uw}(A_0-D_0)
      +b\sqrt{cg}(\sqrt{A_0D_0}-\sqrt{B_0C_0}),\\
 \sqrt{vw}(u-bc)
 &=\frac{b^{3/2}}{\sqrt c}
             (\sqrt{A_0C_0}-\sqrt{B_0D_0}),\\
 \sqrt{uw}(v-cg)
 &=b\sqrt{g/c}
             (\sqrt{B_0C_0}-\sqrt{A_0D_0}),\\
 \sqrt{uv}(w-bg)
 &=\frac{b^{3/2}\sqrt g}{c}
             (\sqrt{A_0B_0}-\sqrt{C_0D_0}).
\end{aligned}                                                     \tag{SC9}
\]
The coefficient factors on the right are infinitesimal; all radicals
separate between the source and testing fields. Consequently any fixed
positive real multiples of the squares of the four left-hand roots can
be inserted as additional positive Gram pivots at an infinitesimal cost
in the displayed strict circuit Gram. To see this without a limiting
assertion, write the four identities as $P=TZ$, with $Z$ the displayed
circuit roots and every entry of $T$ in the bounded maximal ideal. Then
$Z^TGZ=\tau P^TP+Z^T(G-\tau T^TT)Z$, and the second matrix retains
the strictly positive real residue of $G$ for every fixed real $\tau>0$.
This is a genuine full-rates Gram mechanism. It still requires showing
that the entire first-normal and scalar remainders of (SC3) allocate on
these weighted roots; the identities do not make that last membership
automatic.

\paragraph{A quantified remote improvement when $e$ is unbounded.}
The source identity $e^2=q/g^2$ and (NI3) imply that $m,\beta$ are
both unbounded when $e$ is unbounded. Three inverse-cycle rotations
send this adjacent pair to $C,A$ and transform tests as
$(u,v)\mapsto(v,1/u)$. Therefore the adjacent-event theorem covers
all tests with $v$ not infinitesimal. Combined with the previous
$C,A$ admission, only $u,v$ infinitesimal remain in this $e$ type.
This includes every $q$-unit source considered above.

There is also a complete state-chart admission there. Normalize the
pure-testing state by a coordinate of maximal absolute value. If
$W$ is a unit, cap $C,A,m,\beta$ at a sufficiently large positive
real $\mu$, retain all excesses, and leave every other event as its
original nonnegative square. At $u=v=X=Y=0$, the quadratic part in
$(u,v,X,Y)$ is
\[
 2(uW-Y)^2+2(X+uZ)^2-4uvWZ+
 4\mu\{(uW-Y)^2+Y^2+(vW-X-uZ)^2+(X+uZ)^2\}.
 \tag{SC8}
\]
The event part is strictly positive for every real $W\ne0$ and bounded
real $Z$: its zero equations give successively $Y=u=X=v=0$. A sufficiently
large finite $\mu$ therefore absorbs the entire signed quadratic core.
The exact remainder has degree at least three in these four small
roots, with bounded coefficients in $W,Z$, so it allocates with zero
residue. This proves the entire original parent on this chart.
If the core has negative sign and $W$ is infinitesimal, then $Z$ is
a unit, since otherwise the core residue forces $X=Y=0$, contrary
to maximal-coordinate normalization. Thus the remaining state chart
has $Z=1$, $W$ infinitesimal, and both testing half-angles infinitesimal.
This is a full quantified narrowing, not closure of C1.1.

\paragraph{Completion and exact outstanding obligation.}
The following section supplies a finite scalar square lift and both
complete first normals for the $q$-unit C1.1 subdomain, including
$D=0$. It uses a contact-centered $D$ packet followed by the geometric
mean conductor; an unweighted $E_0^2$ perturbation by itself is not
controlled by the circuit uniformly at arbitrary relative rates.
For $q$ infinitesimal, the uniform normal and protected contact lift
in Sections~\ref{pc376284fc7ae:sec:r7-c1-q-inf-local}--\ref{pd2ecc6ac1891:sec:r7-c1-q-inf-remote}
resolve the additional $g/\sqrt q$ issue and all remote tests.
No assumption of bounded $e$ or $1/e$ is imposed. The original
checker verifies thirty-nine whole-polynomial and source identities;
the later modules include their own exact checks and scoped review.


\section{\texorpdfstring{C1 UH: a homogeneous biquadratic conductor}{C1 UH: a homogeneous biquadratic conductor}}
\label{component:pf7ffa921793c}

\label{pf7ffa921793c:sec:r7-c1-uh}
All coefficients below are divided by the original positive pure-source
$d$. We treat precisely (NI9) and (NI10), with no alteration of the native
source (NC1)--(NC7). The revision-6 common-event theorem already covers
all tests except positive unit $u$ and $w=1/v$ infinitesimal.
Put $x=k^2$, $K=1+x$, $\sigma=1/\eta$, $z=1/\zeta=q^2$,
$b=r/z$, $H=h/\sigma$, and retain $\sigma>0$ (it need not be a unit).
On (NI9), $\sigma$ is infinitesimal and $z$ is a unit; on (NI10),
$z$ is infinitesimal and $\sigma$ is bounded. Throughout, $k,H$ are
positive units, $b,z,\sigma$ are bounded, $Q=1-xb>0$, and
\[
 r=zb,\quad R=1+xzb,\quad M=R-z+\sigma zH,
 \quad J=1+\sigma(xH-1+zb),\quad\theta=1-\sigma H,
 \quad E_s=1-xH.
\]
The exact source polynomials, all with bounded arguments, are
\begin{align*}
 P={}&xK+\sigma x+\sigma z+\sigma zb x^2
 +x^2\sigma zH(1-\sigma H)+2x\sigma^2zH
 +\sigma^2z^2bx^2H,\\
 V={}&z(1+x\sigma H)E_s+x(1-zb)(1+\sigma zH),\\
 B_\alpha={}&R\sigma zH+RxH+ME_s.
\end{align*}
Writing $m=m_{\rm old}/z$, $c=C/z$, $a=Az$, $\ell=\alpha z$,
\[
 c={K(1+\sigma zb)E_s\over2(1-zb)P},\quad
 m={Kx(1+\sigma zb)H(\sigma z+x)\over2MP},\quad
 a={V\over2bP},\quad
 \ell={x^2(1+\sigma zb)B_\alpha\over2QP}.                 \tag{UH1}
\]
Define the polynomial quotients
\[
 N_D={K(\sigma z+x)(1+\sigma zb)-JP\over\sigma},\qquad
 N_\beta={Kx(1+\sigma zb)(1+\sigma zH)-\theta P\over\sigma}.
\]
Then $D=\sigma N_D/(2JP)$ and
$\beta=\sigma N_\beta/(2\theta P)$, exactly.

\paragraph{The bounded full parent.}
Let $D_u=1+u^2$, $D_w=1+w^2$. Its six event amplitudes are
\[
 M_w=W-wX+uY-uwZ,\ R_w=uW-Y,\ Y,\
 J_w=uwW+uX-wY-Z,\ Y-wZ,\ B_w=X+uZ.
\]
Thus the exact parent is
\[
 F=4mzM_w^2+4czD_wR_w^2+4(a/z)D_uD_wY^2
 +4DJ_w^2+4(\ell/z)D_u(Y-wZ)^2+4\beta D_wB_w^2+2K_w, \tag{UH2}
\]
where, without any restriction to a state ray,
\begin{align*}
 K_w={}&u^2w^2Z^2-2u^2wYZ-2uwWZ+2uw^2XZ
       +u^2w^2W^2+2u^2wWX-2uw^2WY\\
      &+(D_u+w^2)(X^2+Y^2)-2uwXY.
\end{align*}
For a potentially unbounded weight choose a finite cap, keeping bounded
weights unchanged. Subtract and retain the following nonnegative
centered excesses:
\begin{align*}
 &4(m-m_*)[qM_w-kM(wZ/q)]^2,\\
 &4(c-c_*)D_w[qR_w-x(1-r)(wZ/q)]^2,\\
 &4(a-a_*)D_uD_w[(Y-xbwZ)/q]^2,
 \quad4(\ell-\ell_*)D_u[(Y-xbwZ)/q]^2.          \tag{UH3}
\end{align*}
The first line uses $qM_w$ and the second $qR_w$; thus these are
exactly excesses of the scaled weights in (UH1). Their remainders
are bounded because $mM,c(1-r),ab,\ell Q$ are bounded.
When $\sigma$ is a unit, additionally retain
\[
 4(D-D_*)(J_w+JZ)^2,
 \qquad4(\beta-\beta_*)D_w(X+uhZ)^2.               \tag{UH4}
\]
Here $DJ$ and $\beta\theta$ are bounded. When $\sigma$ is
infinitesimal, no (UH4) cap is needed: $D,\beta=O(\sigma^2)$,
as proved below. Denote the complete remainder by $G$.
Every retained row vanishes at the actual contact
\[
 u=k,\quad w={q\sqrt\sigma\over k},\qquad
 (W,X,Y,Z)=({\sqrt\sigma\over q},-kh,kbq\sqrt\sigma,1).
                                                               \tag{UH5}
\]
No comparison between independent source and test rates occurs.

\paragraph{Exact homogeneous conductor, including every first normal.}
Put $h=\sigma H$, $\delta=u-k$, and use the six roots
\[
\begin{gathered}
 N_x=X+khZ,\quad N_y=(Y-kbq^2W)/q,\\
 U=wW-\sigma Z/k,\quad N=(wZ-q^2W/k)/q,\\
 S=q\delta W,\quad T=\sigma\delta Z.
\end{gathered}                                     \tag{UH6}
\]
Substitute $X=N_x-khZ$, $Y=qN_y+kbq^2W$ in the entire
bounded remainder $G$. Its complete quadratic normal block is the
coefficient collection in $(N_x,N_y)$. Its scalar and first-normal
coefficients have exactly the following supports:
\begin{align*}
 G_{00}&={A\over q^2}w^2Z^2+B w^2W^2+C\sigma^2 Z^2
                  +D_1 wWZ+E q^2W^2,\\
 G_{10}&=A_xwW+B_xZ+C_xw^2Z,\\
 G_{01}&=q(A_y+B_yw^2)W+C_ywZ/q.                    \tag{UH7}
\end{align*}
All eleven displayed coefficient families and the quadratic normal
coefficients are bounded polynomials in $u$, with bounded source
coefficients; $B_x$ is divisible by $\sigma$. In the unit-$\sigma$
case divisibility means a bounded quotient, not a new polynomial
identity. Formulas below specify every coefficient before caps, and
(UH3)--(UH4) specify their exact changes. No state normal is omitted.

Set $s=\sigma/k$, $a_0=q^2/k$ and
\[
 L=2A/k+2B\sigma/k+D_1,\quad
 J_0=C-B/k^2,\quad K_0=A/k^2+2B\sigma/k^2+D_1/k+E.
\]
The scalar is identically
\[
 A N^2+B U^2+qLWN+\sigma^2J_0Z^2+q^2K_0W^2.       \tag{UH8}
\]
Exact source contact and all stationary jets give
\[
 L(k)=J_0(k)=K_0(k)=0,\qquad K_0'(k)=-\sigma J_0'(k),
\]
\[
 B_x(k)+s[A_x(k)+a_0C_x(k)]=0,\qquad
 A_y(k)+C_y(k)/k+s a_0B_y(k)=0.                    \tag{UH9}
\]
These hold separately for every centered excess, hence for every cap.
Write the monic divisions
\[
 L=\delta L_1,\quad J_0=\delta j+\delta^2J_2,\quad
 K_0=-\sigma\delta j+\delta^2K_2,
\]
\[
 B_x+s(A_x+a_0C_x)=\sigma\delta H_x,\qquad
 A_y+C_y/k+s a_0B_y=\delta H_y.
\]
All quotients are bounded. The scalar and full first normals become
\begin{align*}
 G_{00}={}&A N^2+B U^2+(L_1+k\sigma j)SN-kjTU
                          +J_2T^2+K_2S^2,\\
 G_{10}={}&(A_x+a_0C_x)U+C_xqwN+H_xT,\\
 G_{01}={}&(C_y+s q^2B_y)N+qB_ywU+H_yS.            \tag{UH10}
\end{align*}
Together with the unchanged full quadratic normal block, this is an
exact bounded six-root Gram for $G$, not a scalar restriction.
The only root relation needed for strict completion is
\[
 T^2=\sigma S^2+k\sigma\delta SN-k\delta TU.        \tag{UH11}
\]
Its verification is one multiplication using (UH6). Adding any bounded
multiple of its zero changes no parent. No mixed inverse, root, order,
or source/test identification occurs in (UH6)--(UH11).

For the cap argument it is useful to print every centered row directly
in this same frame. In the order $m,c,(a,\ell),D,\beta$, they are
\begin{align*}
 &uzN_y-(kR+z\delta)N-qwN_x-(zQ/k)S,\\
 &S-zN_y-x(1-r)N,\qquad N_y-xbN,\\
 &(u-kbz)U+uN_x-qwN_y+(E_s/k)T,\qquad N_x+HT.
                                                               \tag{UH20}
\end{align*}
These are exact identities before taking any residue.

\paragraph{Explicit coefficient input.}
For (UH2) before caps put $j_u=u-kbz$, $p_u=1+ukbz$,
$t_u=kh-u$. The full scalar coefficients in (UH7) are
\begin{align*}
 A={}&4mz^2t_u^2+4\ell D_u+(4\beta+2)zt_u^2,\\
 B={}&(4cz+4D+2)j_u^2+4aD_uxb^2z,\\
 C={}&[4D(1+ukh)^2+4\beta(u-kh)^2+2D_uxh^2]/\sigma^2,\\
 D_1={}&8mzp_ut_u-8Dj_u(1+ukh)-8\ell D_ukb
          -4u^2k(h+bz)-4u+4uxhbz,\\
 E={}&4mp_u^2+4cj_u^2+4(a+\ell)D_uxb^2+2D_uxb^2z.
\end{align*}
The first normals are
\begin{align*}
 A_x={}&-8mzp_u+(8D+4)uj_u,\\
 B_x={}&-8Du(1+ukh)+8\beta(u-kh)-4D_ukh,\\
 C_x={}&-8mzt_u+(8\beta+4)(u-kh),\\
 A_y={}&8mzup_u-8czj_u+8(a+\ell)D_ukb+4D_ukbz,\\
 B_y={}&-8czj_u+8aD_ukb-(8D+4)j_u,\\
 C_y={}&8mz^2ut_u-8\ell D_u+8Dz(1+ukh)+4uzt_u.
\end{align*}
The coefficients of $N_x^2,N_y^2,N_xN_y$ are respectively
\[
 4mzw^2+4Du^2+4\beta D_w+2(D_u+w^2),
\]
\[
 4mz^2u^2+4cz^2D_w+4aD_uD_w+4\ell D_u+4Dzw^2
                            +2z(D_u+w^2),
 \qquad-8mzuqw-8Duqw-4uqw.                         \tag{UH12}
\]
Here products such as $Du$ mean $D\,u$, whereas $D_u=1+u^2$.

\paragraph{Infinitesimal $\sigma$: source collapse and strict residue.}
The exact polynomial $P_m=(N_D+xN_\beta)/\sigma$ is bounded.
At $\sigma=0$, $N_D=-xN_\beta$ and
\[
 N_\beta=x(K+z)H+bxz-x-z.
\]
Nonnegativity of $D,\beta$ forces
\[
 H=H_0+O(\sigma),\quad
 H_0={x+zQ\over x(K+z)},\qquad E_0=1-xH_0={R\over K+z}.
                                                               \tag{UH13}
\]
Indeed $JD+x\theta\beta=\sigma^2P_m/(2P)$ and $J,\theta$
are units. Thus $D,\beta=O(\sigma^2)$, including either literal zero.
Moreover
\[
 (D+x\beta)/\sigma^2\longrightarrow KH_0E_0/2>0.
\]
The other limiting coefficients are
\[
 c_0={R\over2x(1-zb)(K+z)},\quad m_0={xH_0\over2(1-zQ)},
 \quad a_0^{\rm wt}={H_0\over2b},\quad
 \ell_0={xR\over2Q(K+z)}.                           \tag{UH14}
\]
At a matched angle $u=k$, the residue Gram (UH10) splits into
$(N_y,N,S)$ and $(N_x,U,T)$, because their only off-block coefficients
contain $w$. Before caps, the $(N_y,N)$ block is positive definite;
its angular Schur complement is
\[
 \Gamma=8RH_0/K>0.                                \tag{UH15}
\]
The $(N_x,U)$ block is precisely
\[
 2K(N_x+xH_0U)^2+2KxH_0E_0U^2.                   \tag{UH16}
\]
The remaining $T$ couplings are bounded. As $\sigma,\delta$ have
zero residue, (UH11) permits an arbitrarily large positive $T^2$
pivot without changing any other residue coefficient. This proves a
strict full six-root residue Gram. For $u^2\ne x$, use the full raw
frame $(E_t,F_t,G_t,R_t,X,T_0)=(qW,Y/q,wZ/q,wW,X,\sigma Z)$.
Its leading first block is
\begin{align*}
 &4m_0[E_t+uz(F_t-G_t)]^2+4c_0(uE_t-zF_t)^2
       +4a_0^{\rm wt}D_uF_t^2+4\ell_0D_u(F_t-G_t)^2\\
 &\qquad+2z\{u^2(G_t-F_t)^2+F_t^2\}-4uE_tG_t.
\end{align*}
Put $v_0=u^2$ and $M_0=1-zQ$. In $(F_t,F_t-G_t)$ the hard block
is diagonal, with weights
\[
 B_h={2(x+zQ)[v_0(1-r)+R]\over bx(1-r)(K+z)},\qquad
 L_h={2R[v_0(x+zQ)+xM_0]\over Q(K+z)M_0}.
\]
The exact $E_t$ Schur complement is
\[
 {2(v_0-x)^2(1+v_0)RH_0\over
 [v_0(1-r)+R][v_0(x+zQ)+xM_0]}.                  \tag{UH17}
\]
It has a positive-unit limit through every opening when $v_0\ne x$.
The hard minimizers are
\[
 {F_t\over E_t}={ubK\over v_0(1-r)+R},\qquad
 {F_t-G_t\over E_t}={-uKQ\over v_0(x+zQ)+xM_0},
\]
also bounded through every opening. The other raw blocks are
$2\{u^2R_t^2+2u^2R_tX+D_uX^2\}$ and
$4(D+u^2\beta)T_0^2/\sigma^2$, both strictly positive.
This gives the full bounded away-angle lift, without a determinant-only
claim or any omitted first state normal.

\paragraph{Unit $\sigma$: angular Schur sign.}
Here (NI10) gives $z=0$ at residue. Put $h=\sigma H$,
$J=1+\sigma(xH-1)$, $\theta=1-\sigma H$. The source signs imply
\[
 0<\sigma<K,\qquad {1\over K+\sigma}\le H
       \le{x+\sigma\over x(K+\sigma)}.             \tag{UH18}
\]
These inequalities are stated on the interior; its closure includes
$\sigma=K$ when $J=\theta=0$.
The limits are
\[
 m={KxH\over2(K+\sigma)},\quad c={K(1-xH)\over2x(K+\sigma)},
 \quad a={1\over2b(K+\sigma)},\quad
 \ell={x\over2Q(K+\sigma)},
\]
\[
 D={K\over2J(K+\sigma)}-\tfrac12,
 \quad\beta={K\over2\theta(K+\sigma)}-\tfrac12.
\]
The two hard blocks in (UH10) are positive definite. Their angular
Schur coefficients $\gamma_S,\gamma_T$ have no $ST$ term and satisfy
\[
 \gamma_S+\sigma\gamma_T
 ={ -2p(H)\over\sigma(K+\sigma)}>0,\qquad
 p(H)=Kx(K+4\sigma)H^2-2K(x+2\sigma)H+x.           \tag{UH19}
\]
Convexity of $p$ and (UH18) prove strict negativity: its endpoint
values are respectively
\[
 {-\sigma(3\sigma x+4\sigma+4x+4)\over(K+\sigma)^2},\quad
 {-\sigma(4\sigma x+3\sigma+4x^2+4x)\over x(K+\sigma)^2}.
\]
Use (UH11) to distribute this positive total between $S^2,T^2$;
for example choose its multiple so both resulting Schur coefficients
are positive. This yields the required strict six-root Gram.

For completeness, angles $u^2\ne x$ are strictly positive as well.
Eliminating the two hard state variables from the leading form,
multiplying by the positive source factor $(K+\sigma)/2$, and
writing $U_0=u^2$, $P_0=(w/q)(Z/q)$, $T_0=(Z/q)^2$ gives a positive
Laurent function
\[
 A+B P_0^2+C P_0^2/T_0+D_0T_0-2E P_0,
\]
where
\[
 A=KxH+K(1-xH)U_0/x,\ B=x(1+U_0),\
 C={KU_0\over J+\theta U_0},
\]
\[
 D_0=(1+U_0)\left[{K(1+U_0)\over J+\theta U_0}-(K+\sigma)\right],
 \quad E={K\sqrt{U_0}(1+U_0)\over J+\theta U_0}.
\]
The strict criterion is $E<\sqrt{AB}+\sqrt{CD_0}$.
Squaring only to analyze these positive real residue quantities gives
\[
 (E^2-AB-CD_0)^2-4ABCD_0
 ={K^2(1+U_0)^2\over(J+\theta U_0)^2}(U_0-x)^2Q_2(U_0).
\]
The quadratic $Q_2$ has positive constant $H^2x^2J^2$, positive
leading coefficient $\theta^2(1-xH)^2$, and $Q_2(x)=xp(H)<0$.
Explicitly, with $E_s=1-xH$,
\[
 Q_2(U_0)=E_s^2\theta^2U_0^2
 +2E_s\{x\theta^2+E_s\theta J+(K+\sigma)\theta-2K\}U_0
 +x^2H^2J^2.
\]
The concave quadratic
$U_0(K+\sigma)-[x^2H+(1-xH)U_0](J+\theta U_0)$
is negative at zero and infinity and positive at $x$; at each of its
two zeros the preceding discriminant is negative. Consequently its
positive interval lies inside $Q_2<0$. This proves the criterion for
all $U_0\ne x$. Sign charts with $P_0\le0$ are immediate. These
square roots are an analysis of real residues only and are not tensor
roots or mixed localizations.

This inequality is directly a bounded Gram allocation, not a
comparison of mixed ratios. Before choosing projective coordinates,
put $E_t=qW$, $G_t=wZ/q$, $R_t=wW$. The exact relation
$E_tG_t=R_tZ$ permits distributing the cross term between the two
quadratic blocks $(E_t,G_t)$ and $(R_t,Z)$. The strict inequality
above is exactly the condition that both $2\times2$ blocks can be
made positive definite by a finite real allocation. It therefore
handles every independent rate of $q,w$.
At $J=0$ or $\theta=0$, the root-interval argument permits one endpoint
at zero or infinity. At the double face $J=\theta=0$, necessarily
$\sigma=K$, $H=1/K$. As $J+\theta U_0\to0$,
\[
 E-\sqrt{CD_0}\longrightarrow K\sqrt{U_0},\qquad
 AB-K^2U_0=(U_0-x)^2.
\]
Thus the constrained allocation remains strictly positive away from
$U_0=x$, including this double face without dividing by a vanishing
residue denominator.

\paragraph{Opening faces, tensor conclusion, and exact scope.}
The finite-cap argument is applied to the full Gram above. For
$\sigma$ infinitesimal the possible diverging rows have openings
$1-zb$, $1-zQ$, $b$, $Q$; at most two can diverge together. The pairs
$(c,m)$, $(m,a)$, $(c,\ell)$ have independent hard rows; other pairs
are excluded by $0<zb,xb<1$ and $1-zQ>0$. In the $(N_y,N)$ columns
of (UH20), their determinants have absolute values $kzK,k,x$,
respectively; each is a unit on its compatible face. The hard block
is uniformly coercive, since it contains positive-unit multiples
of $N_y^2$ and $(N_y-N)^2$. Its minimizing hard coordinates for
fixed angular $S$ are explicitly
\[
 {N_y\over S}={b(2bxz-x+1)\over K},\qquad
 {N\over S}={4bxz-x-2z+1\over xK}.
\]
They are bounded through every opening. Thus the limit of the
uncapped Schur complement is the Schur complement restricted to the
diverging centered rows' kernel. The quotients (UH15)--(UH16) stay
positive units there. The elementary finite-dimensional penalty
lemma now applies: a bounded symmetric form strictly positive on
the common kernel becomes strictly positive after sufficiently large
finite positive multiples of these rows' squares are added. This
proves the finite-cap claim, with the entire angular right-hand sides
in (UH20) retained. Away from $u^2=x$, (UH17) and its printed bounded
hard minimizers give the same constrained strictness directly.

For unit $\sigma$, only $b,Q,J,\theta$ can be openings.
The $b,Q$ rows coincide in the hard direction $N_y-xbN$ and cannot
both diverge. The $J,\theta$ rows on $(N_x,U)$ are independent even
if both vanish; their determinant is the unit $u$. Formula (UH19)
remains a positive unit through these faces. More explicitly, the
two hard blocks dominate positive-unit multiples of
$N_y^2+(N_y-N)^2$ and $N_x^2+(N_x+U)^2$. Their angular minimizers are
\[
 {N_y\over S}={b(1-HxK)\over K},\quad
 {N\over S}={1-HxK\over xK},
\]
\[
 {N_x\over T}={\sigma H^2xK-HxK-\sigma H(2x+1)+x\over\sigma K},
 \quad {U\over T}={HxK(K+\sigma)-xK-\sigma\over x\sigma K}.
\]
Their $D$ and $\beta$ row defects divided by $T$ are respectively
$k(HK-1)J/(\sigma K)$ and $-x\theta(HK-1)/(\sigma K)$.
Hence every relevant constrained limit exists with bounded
minimizers; the same penalty argument applies. On the angular
relation $T^2=\sigma S^2$, its strictly positive Schur value is
(UH19), and after taking finite caps the bounded zero (UH11)
distributes that value into two strict angular pivots. The away-angle
allocation and its single/double-face limits above give the same
conclusion there. All removed excesses remain exact positive
source/test squares. No positive lower bound on an original opening
has been imposed.

A strict bounded residue Gram lifts by the inherited tensor Gram
lemma. Every root in (UH6) is a finite sum of separated source and
testing factors. After restoring the retained squares this certifies
the complete original parent, including $D=0$ or $\beta=0$.
Together with revision-6 remote tests, this closes both (NI9) and
(NI10). Those are the entire actual UH portion of C1.5. The already
proved UU and UI portions and exact source stabilizers therefore close
the entire C1.5 row, subject to the exact checks and mathematical review
recorded alongside this section.


\section{\texorpdfstring{C1 UI: completion of the ordinary/small strip}{C1 UI: completion of the ordinary/small strip}}
\label{component:pa0b854a04f49}

\label{pa0b854a04f49:sec:r6-c1-ui}
Use the unchanged native C1 source in the representative of (NI7):
$k,\eta$ are positive bounded units and $\zeta=g^2$ is positive
infinitesimal. All original openings and $D=0$ are retained.
The exact stabilizer recovers the other representative; no condition
$k\le1$ is imposed after that change. This section proves the missing
ordinary/small testing strip, and joins the remote theorem (CR1)--(CR13).

\paragraph{Exact bounded remainder.}
Normalize by the positive pure-source $a$. Write $x=k^2$,
$r=1-a\eta$, $K=1+x$, $R=1+xr$, $M=R-t$, and
\[
 J_s=K-xtg^2-a,\quad E=1-x\eta+tx\eta g^2,\quad
 Q_c=1-xg^2r,\quad y=x\eta g^2.
\]
Here $a,t,r,1-r,M,J_s$ are positive bounded openings; $E,Q_c,R,K,r+\eta$
are positive units. The exact polynomials $P_4,P_d,P_\alpha$ are (NC2).
The coefficients after division by $a$ are
\[
\begin{aligned}
 c&=2/a,& A_0&=\frac{2\eta^2P_4}{Kr(r+\eta)E},&
 d_0&=\frac{4\eta P_d}{K(r+\eta)E},\\
 D_0&=\frac{2(1+y)}{J_sE}-d_0/2,&
 \alpha_0&=\frac{2x^2\eta P_\alpha}{KQ_cE},&
 \alpha/a&=g^2\alpha_0.
\end{aligned}                                                     \tag{CU1}
\]
In particular $c,A_0$ are strong; $d_0,\alpha_0$ are units; and
$ca,A_0r,D_0J_s$ are bounded with continuous rational expressions
through every opening. Put
\[
 h_2=\frac{2x\eta(1-tg^2)(1+y)}E,\qquad
 h_b=\frac{2x\eta[1+(1-t)g^2]}E,
 \qquad h_b-h_2=2x\eta g^2.
\]
Both $h_2,h_b$ are positive units. The exact coefficients are
$m/a=h_2/(g^2M)$ and $\beta/a+d_0/2=h_b/(g^2t)$.

For the test put $D_u=1+u^2,D_v=1+v^2$ and
\[
 Q=W+uY,\quad B=X+uZ,\quad p=t/R,
 \quad N=(B-pvQ)/g,\quad L=vQ/g.
\]
The inverse $W=Q-uY$, $X=gN+pvQ-uZ$ is polynomial.
Retain the two positive squares
\[
 \frac{4h_2t}{MR}N^2,\qquad
 \frac{4h_bM}{tR}D_vN^2.                            \tag{CU2}
\]
Their remainder is exactly (CL7). Choose bounded caps
$c_*\le c$, $A_*\le A_0$, $D_*\le D_0$, with $c_*,A_*>0$,
$D_*\ge0$; bounded original weights can be kept unchanged.
Retain the three positive centered excesses
\[
\begin{aligned}
 &4(c-c_*)D_v(r_t-avZ)^2,\quad
 4(A_0-A_*)D_uD_v(Y-rvZ/\eta)^2,\\
 &4(D_0-D_*)(J_t+J_svZ)^2,
\end{aligned}                                                     \tag{CU3}
\]
where
\[
 r_t=uQ-D_uY,\quad J_t=r_t-D_uvZ+ugvN+upv^2Q.
\]
All five retained rows vanish at the actual contact
\[
 (u,v)=(k,k\sqrt\eta g),\quad
 (W,X,Y,Z)=(g/\sqrt\eta,k(tg^2-1),krg/\sqrt\eta,1).
                                                               \tag{CU4}
\]
Set $c_a=2-c_*a$, $A_r=A_0r-A_*r$, $D_J=D_0J_s-D_*J_s$.
After removing (CU2)--(CU3), the entire parent divided by $a$ is
\begin{align*}
G={}&\frac4R\{(h_2+h_bD_v)N^2+2(h_bD_v-h_2)LN
              +[h_2(1-p)+h_bpD_v]L^2\}\\
 &+4c_*D_vr_t^2+8c_aD_vvZr_t-4c_aaD_vv^2Z^2\\
 &+4A_*D_uD_vY^2+\frac{8A_r}{\eta}D_uD_vvZY
                   -\frac{4A_rr}{\eta^2}D_uD_vv^2Z^2\\
 &+4D_*J_t^2-8D_JvZJ_t-4D_JJ_sv^2Z^2\\
 &+4\alpha_0D_u(gvY-gZ)^2
       +2d_0\{J_t^2+D_u(v^2Y^2-2vYZ-v^2Z^2)\}.
                                                               \tag{CU5}
\end{align*}
Thus $G$ is an exact bounded Gram on $(N,Q,Y,vZ,gZ,L)$ whenever
$u$ is a bounded unit and $v$ infinitesimal. Every source denominator
in its coefficients is a unit. Formula (CU5), rather than a scalar
restriction, supplies all normal terms below.

\paragraph{Local contact lifting with all first normals.}
Put $\kappa=k\sqrt\eta$, $\varepsilon=\kappa g$,
$s=R/(k\eta)$, $\delta=u-k$, and define
\[
 E_Q=Q-svZ,\quad E_Y=Y-rvZ/\eta,\quad
 R_h=(vQ-s\varepsilon^2Z)/\varepsilon,
 \quad S=v\delta Z,\quad T=\varepsilon\delta Z.
\]
Let $c_L=4[h_2(1-p)+h_bpD_v]/R$, a positive bounded unit.
The sole potentially unbounded $L$ row can be replaced exactly:
\[
 L=\kappa R_h+kR gZ,\qquad
 L^2-\kappa^2R_h^2=2kR vQZ-k^2R^2g^2Z^2.           \tag{CU6}
\]
Consequently (CU5) equals $c_L\kappa^2R_h^2$
plus $8\kappa(h_bD_v-h_2)NR_h/R$ and a polynomial $G_0$.
Explicitly, in (CU5) replace its first line by
\[
 \frac4R(h_2+h_bD_v)N^2+8k(h_bD_v-h_2)gNZ
                 +c_L(2kRvQZ-k^2R^2g^2Z^2).        \tag{CU7}
\]
All coefficients of $G_0$ are bounded polynomials in $u,v,g$ with
pure-source unit denominators. The coefficient of $NR_h$ is
infinitesimal, since $h_bD_v-h_2=2x\eta g^2+h_bv^2$.

Here are exact prescriptions for every normal coefficient, avoiding
any implicit omission. Substitute $Q=E_Q+svZ$ and
$Y=E_Y+rvZ/\eta$ in $G_0$ and collect the quadratic polynomial
in $(N,E_Q,E_Y)$. The constant coefficient is $\phi(u,\xi)Z^2$,
where $\xi=v^2$ and $\phi\in(g^2,\xi)$. The linear coefficients
of $E_Q,E_Y$ are $vL_Q(u,\xi)Z,vL_Y(u,\xi)Z$; that of $N$ is
$gL_N(u,\xi)Z$, where $L_N\in(g^2,\xi)$. These parity and ideal
assertions follow term by term in (CU5),(CU7). The remaining quadratic
normal block is the full coefficient collection of the same expression.
For explicit pen-and-paper formulas, set
\[
 b=r/\eta,\quad F_0=us-D_ub,\quad
 G_0^{\rm sc}=F_0-D_u+up\xi s,\quad
 H_0=(8D_*+4d_0)G_0^{\rm sc}-8D_J.
\]
Then, writing $D_v=1+\xi$ in these polynomials,
\begin{align*}
 L_Q={}&2c_LkR+8D_vu(c_*F_0+c_a)+u(1+p\xi)H_0,\\
 L_Y={}&-8D_vD_u(c_*F_0+c_a)
       +8D_uD_v(A_*b+A_r/\eta)-D_uH_0\\
      &+8\alpha_0D_ug^2(b\xi-1)+4d_0D_u(b\xi-1),\\
 L_N={}&8k(h_bD_v-h_2)+u\xi H_0.
\end{align*}
The complete scalar coefficient is
\begin{align*}
 \phi={}&c_L(2kRs\xi-k^2R^2g^2)
   +4c_*D_vF_0^2\xi+8c_aD_vF_0\xi-4c_aaD_v\xi\\
 &+4A_*D_uD_vb^2\xi+8(A_r/\eta)D_uD_vb\xi
                    -4A_rrD_uD_v\xi/\eta^2\\
 &+4D_*(G_0^{\rm sc})^2\xi-8D_JG_0^{\rm sc}\xi-4D_JJ_s\xi\\
 &+4\alpha_0D_ug^2(b\xi-1)^2
   +2d_0\{(G_0^{\rm sc})^2\xi+D_u(b^2\xi^2-2b\xi-\xi)\}.
\end{align*}
These formulas display the two asserted ideals directly. They also
give a short explicit input for the monic divisions, without an
unprinted state-normal computation.

The source is an exact stationary zero and every removed square
vanishes there. Therefore
\[
 \phi=\phi_u=\phi_\xi=0,\qquad L_Q=L_Y=L_N=0
                       \quad\text{at }(u,\xi)=(k,\varepsilon^2).
\]
Successive monic polynomial division gives, with bounded coefficients,
\[
 \phi=\delta^2(\varepsilon^2B_0+\xi B_1)
       +\delta(\xi-\varepsilon^2)B_2
       +(\xi-\varepsilon^2)^2B_3,                  \tag{CU8}
\]
\[
 L_i=\delta L_{i0}+(\xi-\varepsilon^2)L_{i1} (i=Q,Y),\quad
 L_N=\delta(g^2L_{N0}+\xi L_{N1})
                         +(\xi-\varepsilon^2)L_{N2}.
\]
For clarity, each quotient is obtained by subtracting the value at
$u=k$ and dividing by the monic $u-k$, then by the monic
$\xi-\varepsilon^2$. The second $u$ quotient of $\phi$ stays in
$(g^2,\xi)$; its constant and first $u$ terms have the stated double
and simple zeros in $\xi$. The same argument gives the $L_N$ ideal.
No small source number is inverted in these divisions.

Put $R_s=(v^2-\varepsilon^2)Z/\varepsilon$. The exact root identities
\[
 \varepsilon R_s=(\varepsilon R_h-vE_Q)/s,\quad
 R_sT=(R_hT-SE_Q)/s,\quad
 S^2-T^2=\delta(R_hT-SE_Q)/s                       \tag{CU9}
\]
allocate (CU8) and every first normal on the six roots
$(N,E_Q,E_Y,R_h,S,T)$. For example $v\delta Z=S$,
$\delta(\xi-\varepsilon^2)Z^2=(R_hT-SE_Q)/s$, and
$(\xi-\varepsilon^2)^2Z^2=(\varepsilon R_h-vE_Q)^2/s^2$.
The $N$ normal has only infinitesimal couplings. All quadratic normal
coefficients are already bounded. This is the required exact full lift.

\paragraph{Strict matched Hessian and source boundary caps.}
First take finite interior source residues. The leading form is
(CL2). Put $V=k\sqrt\eta V_1$, $W=W_1/\sqrt\eta$,
$Y=kY_1/\sqrt\eta$, $Z=1$. At the contact
$(W_1,Y_1,V_1,u)=(1,r,1,k)$, the Schur complement of the two state
coordinates in half the Hessian, in coordinates $(V_1,u)$, is
\[
\begin{aligned}
 G_{11}&=\frac{32\eta x^2P_4}{KE_0},&
 G_{12}&=\frac{16\eta k^3(R^2-2KE_0)}{KE_0R},\\
 G_{22}&=\frac{32\eta x(P_4+xrP_d)}{R^2P_4},&
 \det G&=\frac{256\eta^2x^3(4K^2E_0-R^2)}{K^2E_0^2}.
\end{aligned}                                                     \tag{CU10}
\]
Here $E_0=1-x\eta$, $P_4=KE_0+x(1-r)$, $P_d=1+2x+rx^2$.
There are no opening denominators. The leading $D\ge0$ condition is
\[
 K(2+xr)E_0-R^2\ge0,
\]
so $4K^2E_0-R^2\ge K(4K-2-xr)E_0>0$.
The state block is uniformly bounded below since $c,A_0$ are strong.
Thus the matched Hessian is strict, uniformly through every permitted
source boundary after the appropriate hard directions are retained.

We give the boundary justification explicitly. Divergences can occur
only at $a=0$, $r=0$, or $J_s=0$, for $c,A_0,D_0$, respectively.
Since $a+J_s\sim K$, $c$ and $D_0$ cannot both diverge. If $a=0$,
then $r=1$, so $c$ is the only divergent weight. The only possible
divergent pair is $A_0,D_0$. Their hard-state rows are $E_Y$ and
$kE_Q-KE_Y$, with nonzero determinant $k$. All angular parts of
these centered rows are bounded. Solving these at most two real
linear equations therefore supplies a bounded competitor for the
minimizing state vector. The bounded finite part and (CU10) imply
that the state minimizer is bounded as well, by completing its
positive state square. The full matched Hessian hence has a uniform
positive lower bound on the Euclidean norm in a nearby finite-source
family. On a vector in the limiting hard kernel, adjust only the two
state coordinates to annihilate the nearby exact hard rows. This
adjustment tends to zero, and the divergent penalties vanish exactly.
The bounded finite part is consequently strictly positive on that
kernel. The elementary real finite-penalty lemma supplies finite caps.

Every allowed residue is approachable by such finite sources. For
fixed $x>0$, the real leading inequalities are
\[
 0\le r\le1,\qquad
 \eta_{\min}(r)=\frac{1-r}{K}\le\eta\le
 \eta_{\max}(r)=\frac{1-R^2/[K(2+xr)]}{x},
\]
with positive interval width $R/[xK(2+xr)]$. Perturb $r$ into
$(0,1)$ and its interval coordinate into $(0,1)$.
This makes $a,r,J$ and the $D$ numerator strictly positive while
converging to the given residue. Taking $t=R/2$ and then sufficiently
small real $g>0$ restores all original source signs, including
$\beta>0$. This construction, rather than substitution into a
singular leading denominator, justifies the boundary limit.

In the six-root residue Gram, $N$ is an isolated positive pivot.
The $(E_Q,E_Y)$ block is strict, $R_h$ has positive pivot
$c_L\kappa^2$, $S$ couples only to $E_Q,E_Y$, and $T$ only to $R_h$.
Setting $S=T$ is the matched Hessian under an invertible real change
of coordinates: $R_h/g$ has relative-angle derivative $2s\kappa\ne0$,
and $T/g=\kappa\delta$ supplies the ordinary-angle derivative.
After the hard Schur complements, the sum of the $S,T$ diagonals is
therefore strictly positive. Add a real multiple of the last zero
in (CU9) to redistribute this sum into two positive diagonals.
The extra coefficients contain $\delta$ and vanish in the residue.
The bounded strict-residue Gram lemma proves the full parent at
$u-k$ and $v$ infinitesimal, including every source opening.

\paragraph{Ordinary angles away from contact.}
Let now $\bar u^2\ne\bar k^2$, with $u$ a unit and $v$ infinitesimal.
Use the bounded frame $(N,r_t,Y,vZ,gZ,L)$ of (CU5), and the exact zero
\[
 L(gZ)-Q(vZ)=0,\qquad Q=(r_t+D_uY)/u.              \tag{CU11}
\]
For finite leading source weights, $N$ is an isolated positive pivot.
The remaining residue is
\[
 q^TPq+2Ub^Tq+c_2U^2+4hL^2+c_0T^2,
 \quad q=(r_t,Y)^T,\ U=vZ,\ T=gZ,
\]
where the source coefficients are (CL1) and
\[
\begin{gathered}
 s_D=4D+2d,\quad s_C=4C+s_D,\quad
 P=\operatorname{diag}(s_C,4AD_u),\\
 \ell=(1/u,D_u/u)^T,\quad b=(-s_DD_u,-2dD_u)^T,\\
 c_0=4\alpha D_u,\qquad c_2=s_DD_u^2-2dD_u.
\end{gathered}                                                     \tag{CU12}
\]
Apply the real rank-one allocation lemma (UA1)--(UA3), using (CU11).
We verify its hypotheses uniformly through all source boundaries.
Writing $d_a=\ell^TP^{-1}\ell$, $e_a=\ell^TP^{-1}b$,
$h_a=c_2-b^TP^{-1}b$, direct inversion yields
\[
\begin{aligned}
 d_a&=\frac{aJE_0}{8u^2(JE_0+a)}+\frac{D_u}{4Au^2},\\
 e_a&=-\frac{aD_u}{u(JE_0+a)}-\frac{dD_u}{2Au},\\
 h_a&=\frac{8D_u^2}{JE_0+a}-2dD_u-\frac{d^2D_u}{A},\\
 P^{-1}b&=(-aD_u/(JE_0+a),-d/(2A))^T.
\end{aligned}                                                     \tag{CU13}
\]
All are bounded continuous, since $JE_0+a$ is a positive unit and
$1/A$ has a continuous bounded formula proportional to $r$.
The same holds for $P^{-1}\ell$.

Normalize the Laurent polynomial (CL4) by $As_C$.
Its $U$ and $T^{-1}$ coefficients are, respectively,
$1/[2(JE_0+a)]$ and $\alpha/4$, positive units.
Its $T$ coefficient is
\[
 h\{(D+d/2)/s_C+2CD/(As_C)+d/(8A)\},
\]
also a positive unit: when $1/A=0$, $r=0$ forces $a$ a unit,
so $(D+d/2)/s_C=a/[4(JE_0+a)]$ is a unit; otherwise the last term
is positive. The $U^{-1}$ coefficient is
$h\alpha(1/s_C+1/(4A))$. If it vanishes, $r=0$ and $J=0$,
and the $T/U$ coefficient has the positive term $hD/s_C=h/4$.
All remaining normalized nonconstant coefficients are bounded and
nonnegative. The constant is bounded by the exact contact equation.
Thus strict convexity and coercivity in logarithmic coordinates
persist at every boundary. The normalized Laurent function retains
its unique stationary zero $(U,T)=(x,x\eta)$.

For fixed $u^2\ne x$, the rank-one determinant polynomial is
\[
 p(T)=c_0+(h_a+4hd_ac_0)T+4h(d_ah_a+e_a^2)T^2
      =\frac{16D_u}{As_C}T\mathcal L(u^2,T)>0.
\]
Its constant and leading coefficients are positive units by the
preceding coefficient bounds. The strict allocation interval
therefore survives at every source boundary. The two real Gram
blocks have a uniform positive lower bound; (CU13) controls their
state minimizers even if $P$ is unbounded. The centered divergent
rows have hard parts $r_t$, $Y$, $r_t$; the only simultaneous pair
is $Y,r_t$, whose determinant is one. Exact small state corrections
annihilate the large rows in the nearby finite family. The same
finite-penalty argument as above supplies caps in (CU3). The exact
bounded remainder (CU5), after (CU11), has strict residue.
This proves every ordinary/small test away from contact.

\paragraph{Exhaustive admission.}
The two arguments cover the whole ordinary/small testing strip;
the exact gauges cover negative contact signs. The remote theorem
(CR1)--(CR13), including CR5a, covers its complement.
All transformations are homogeneous in the state; no state coordinate
was divided out in the tensor proof. Real and imaginary state parts
are treated separately. Hence the entire actual UI subdomain of C1.5,
and its exact stabilizer images, has an all-tests tensor-SOS proof.
The actual UH subdomain of C1.5 remains a separate obligation.


\section{\texorpdfstring{C1 UI: the positive Laurent model}{C1 UI: the positive Laurent model}}
\label{component:p050c7ae6d9b8}
% Edited scoped mathematical packet, 22 October 2026.

\label{p050c7ae6d9b8:sec:c1-ui-leading}

This is the positive Laurent-model lemma used in the complete UI proof.
Write $x=k^2$, $r=1-a\eta$, $K=1+x$, $R=1+xr$ and
\[
 E_0=1-x\eta,\quad J=K-a,\quad
 P_4=KE_0+x(1-r),\quad P_d=1+2x+rx^2.
\]
At a real interior source point the leading coefficients of the
parent divided by $a$ are
\[
\begin{aligned}
 h&=\frac{2x\eta}{RE_0},& C&=\frac2a,&
 A&=\frac{2\eta^2P_4}{Kr(r+\eta)E_0},\\
 d&=\frac{4\eta P_d}{K(r+\eta)E_0},&
 D&=\frac2{JE_0}-\frac d2,&
 \alpha&=\frac{2x^2\eta R}{KE_0}.
\end{aligned}                                                     \tag{CL1}
\]
All are positive except $D\ge0$, including its literal zero. Here
$\alpha$ is the coefficient after extracting the weak $g^2$ scale.
For a test $(u,V)$ put $D_u=1+u^2$ and define
\begin{align*}
f={}&4hV^2(W+uY)^2+4C(uW-Y)^2+4AD_uY^2\\
 &+(4D+2d)(uW-Y-VD_uZ)^2+4\alpha D_uZ^2\\
 &-4dD_uVYZ-2dD_uV^2Z^2.                           \tag{CL2}
\end{align*}
Let $M_f$ be its symmetric three-state matrix and set $U=u^2$, $T=V^2$.
Direct collection gives
\[
 \det M_f=64D_u^2UT\,\mathcal L(U,T),                 \tag{CL3}
\]
where
\begin{align*}
\mathcal L={}&AC(D+d/2)U+
 \frac{h\alpha(A+C+D+d/2)}U+\frac{A\alpha(C+D+d/2)}T\\
 &+hD(A+C+d/2)\frac TU+hC(D+d/2)UT\\
 &+h(AD+2CD+Ad/2+Cd/2+Dd/2+d^2/4)T+L_0,\\
L_0={}&ACD-ADd/2+h\alpha(C+D+d/2)
                         -d^2(A+C+D)/4-d^3/8.       \tag{CL4}
\end{align*}
Every nonconstant coefficient is nonnegative. The coefficients of
$U,U^{-1},T,T^{-1}$ are strictly positive, also when $D=0$.
Thus in coordinates $(\log U,\log T)$ the Hessian is strictly
positive and the function is coercive. Substitution of (CL1) gives
\[
 \mathcal L(x,x\eta)=\partial_U\mathcal L(x,x\eta)
                   =\partial_T\mathcal L(x,x\eta)=0. \tag{CL5}
\]
Strict convexity proves that $\mathcal L$ is positive away from this
single contact. The $(W,Y)$ principal block of (CL2) is strictly
positive for $U,T>0$, because its $C$ and $A$ terms already have
trivial common kernel. Hence $f$ is positive semidefinite everywhere
and positive definite off the contact locus. Literal endpoints follow
by continuity. This proof uses no numerical determinant table.

There is a useful corrected formula at the large-$D$ boundary.
Impose its hard equation $uW-Y-VD_uZ=0$, eliminating $W$.
The resulting two-state determinant is $16D_u^2T\mathcal L_D$, where
\begin{align*}
\mathcal L_D={}&hT\{(C+A+d/2)/U+(2C+A+d/2)+CU\}\\
 &+h\alpha(1+1/U)+AC(U+1)+A\alpha/T-Ad/2-d^2/4.
                                                               \tag{CL6}
\end{align*}
In particular the coefficients of $T/U$ and $UT$ are
$h(C+A+d/2)$ and $hC$, respectively. Interchanging these coefficients
is incorrect. All three identities follow by polynomial collection in the displayed model.


\section{\texorpdfstring{C1 UI: a full-parent remote complement}{C1 UI: a full-parent remote complement}}
\label{component:p013c226a79b3}

\label{p013c226a79b3:sec:c1-ui-remote}
The source has actual ordinary Alice angle and infinitesimal Bob
angle. Apply the exact native stabilizer (AS5) as in (NI7) and take
$\eta$ a positive bounded unit, $\zeta=g^2$ infinitesimal and $k$ a
bounded unit. The reflected orientation need not satisfy $k\le1$.
Its $E$ remains a positive unit by the exact source entry in (NI7).
Normalize by the positive pure-source $d$.
Equations (NC2)--(NC4) give
\[
 C/d,\ A/d,\quad g^2m/d,\quad g^2\beta/d,\quad
 \alpha/(dg^2)\quad\hbox{strong},\qquad D/d\ge0.       \tag{CR1}
\]
For example $P_d,P_\alpha,Q,r+\eta$ are units, whereas the
positive openings $M,t,r,1-r$ are bounded and may be infinitesimal.
The $-1/2$ in $\beta/d$ disappears after multiplication by $g^2$.
In this particular source subdomain $\beta=0$ is impossible; $D=0$
is allowed throughout. Choose positive real constants
$\mu,\nu,c,a,b$ below the five quantities in (CR1), respectively
in the order $g^2m/d,g^2\beta/d,C/d,A/d,\alpha/(dg^2)$.
Whenever needed we cap these coefficients at
\[
 m=\mu/g^2,\quad\beta=\nu/g^2,\quad C=c,\quad A=a,
 \quad\alpha=bg^2.                                  \tag{CR2}
\]
Every unused packet, every positive excess, and the complete $D$
packet is retained as a square. The signed core, divided by two, is
\[
 K=J^2+D_u(vY-Z)^2+D_v(X+uZ)^2-D_uD_vZ^2,
 \quad J=uW+uvX-Y-vZ.                              \tag{CR3}
\]
Only the core and the explicitly selected caps need a bounded Gram.
This observation keeps all upper coefficient scales out of the proof.

\begin{theorem}[C1 UI remote tests]
For every source satisfying (CR1), the complete original C1 parent
is tensor SOS at every independent testing pair except possibly
$u$ a bounded nonzero unit and $v$ infinitesimal. The conclusion
includes literal projective endpoints, both signs, all states,
and $D=0$.
\end{theorem}
\begin{proof}
All constants in residue forms below are real residues. Roots may be
unbounded, but all coefficients in their exact Gram are bounded.
First put $Q=W+uY$, $B=X+uZ$; then
\[
 M=vQ-B,\quad uW-Y=uQ-D_uY,\quad
 J=uQ-D_uY+uvB-vD_uZ.                               \tag{CR4}
\]
These identities account for every state normal.

If $u$ is a bounded unit and $v$ projectively infinite, put $w=1/v$
and homogenize by $w^2$. Use roots
$(E_Q,E_B,Y,Z)=(Q/g,B/g,Y,Z)$ and cap only $m,\beta$.
The exact homogenized core is
\[
 [w(ugE_Q-D_uY)+ugE_B-D_uZ]^2
 +D_u(Y-wZ)^2+D_wg^2E_B^2-D_uD_wZ^2.
\]
Its residue is $D_u(Y^2+u^2Z^2)$. Including the two caps gives
\[
 4\mu E_Q^2+4\nu E_B^2+2D_u(Y^2+u^2Z^2)>0.         \tag{CR5}
\]

If $u$ is infinitesimal and $v$ projectively infinite, the preceding
four-root frame loses the $Z$ pivot. Keep $w=1/v$ and instead use
$(E_Q,E_B,Y,U,T)$, where $U=uZ$ and $T=gwZ$. Put
$A_0=w(ugE_Q-D_uY)+ugE_B$. The exact homogenized core is
\[
 K_w=A_0^2+D_uY^2+D_wg^2E_B^2+D_uU^2
       -2D_ug(wE_Q+E_B)U+2wD_uuYU.
\]
All its coefficients are bounded, with residue $Y^2+U^2$.
Cap $m,\beta,\alpha$; their amplitudes after normalization are
$E_Q-wE_B$, $E_B$, and $gY-T$, with the positive factors
$D_w,D_u$ in the last two packets. The full selected residue is
\[
 4\mu E_Q^2+4\nu E_B^2+2Y^2+2U^2+4bT^2>0.       \tag{CR5a}
\]
This includes literal $u=0$ or $w=0$; no division by either was used.

If $u$ is infinitesimal and $v$ a bounded unit, use
$E_Q=Q/g,E_B=B/g,Y,U=uZ,Z_g=gZ$, and cap $m,\beta,\alpha$.
Substitute $W=gE_Q-uY$, $X=gE_B-uZ$ in (CR3).
All terms involving $Z$ collect into $uZ$; the resulting exact
core is a polynomial in $(u,v,g;E_Q,E_B,Y,U)$.
Its residue is $D_vY^2+v^2U^2$. The full selected residue is
\[
 4\mu(vE_Q-E_B)^2+4\nu D_vE_B^2
       +2D_vY^2+2v^2U^2+4bZ_g^2>0.                 \tag{CR6}
\]
The $\alpha$ amplitude is $gvY-Z_g$, so no inverse of $u$ was used
in constructing this Gram.

If both $u,v$ are infinitesimal, set
\[
 U=uQ,\quad V=vQ/g,\quad E_B=B/g,\quad Z_g=gZ,
 \quad S=uvZ,
\]
and include $Y$ as the sixth root. The exact core is
\begin{align*}
K={}&(U-D_uY+uvgE_B)^2+D_uv^2Y^2+D_vg^2E_B^2+D_uS^2\\
 &+2D_uuYS-2D_uuVZ_g-2D_uvgE_BS.                 \tag{CR7}
\end{align*}
In particular the apparently problematic term $UvZ$ is exactly
$uVZ_g$, with no mixed inverse. Cap all five weights in (CR2).
The entire selected residue is
\[
 4\mu(V-E_B)^2+4\nu E_B^2+(4c+2)(U-Y)^2
          +4aY^2+4bZ_g^2+2S^2>0.                  \tag{CR8}
\]

It remains to cover $u$ projectively infinite. Put $p=1/u$,
homogenize by $p^2$, and use $D_p=1+p^2$,
$Q=pW+Y$, $B=pX+Z$. Then
\[
 M=vQ-B,\quad r=W-pY,\quad
 J=D_p(W+vX)-p(Q+vB).                              \tag{CR9}
\]
For $v$ a bounded unit, use $(Q/g,B/g,W,X)$ and cap $m,\beta,C$.
The core after $Y=gE_Q-pW$, $Z=gE_B-pX$ is a bounded polynomial
with residue $(W+vX)^2$. The selected residue is
\[
 4\mu(vE_Q-E_B)^2+4\nu D_vE_B^2
           +4cD_vW^2+2(W+vX)^2>0.                 \tag{CR10}
\]

For infinitesimal $v$, use $(W,R,Y,V,E_B)$ with
$R=vX$, $V=vQ/g$, $E_B=B/g$. Cap $m,\beta,C,A$, retaining
$\alpha$ entirely. The exact core is
\begin{align*}
K_p={}&(W+D_pR-pY-pvgE_B)^2+D_vg^2E_B^2\\
 &+D_p\{v^2Y^2-2vgYE_B+2pYR-(vgE_B-pR)^2\}.
                                                               \tag{CR11}
\end{align*}
Every coefficient is bounded. The selected residue is
\[
 4\mu(V-E_B)^2+4\nu E_B^2+4cW^2+4aY^2
                    +2(W+R)^2>0.                 \tag{CR12}
\]

Finally, when both testing angles are infinite, also put $w=1/v$
and homogenize by $p^2w^2$. Use $(Q/g,B/g,W,X)$ and cap
$m,\beta,C$. Substitution $Y=gE_Q-pW$, $Z=gE_B-pX$ in the
homogenized core is polynomial and has residue $X^2$.
The selected residue is
\[
 4\mu E_Q^2+4\nu E_B^2+4cW^2+2X^2>0.               \tag{CR13}
\]
Each displayed positive form is strict on the complete root space.
The bounded positive-residue Gram lemma therefore supplies tensor
SOS. Restoring the retained packets proves the original parent.
The bounded unit/unit testing chart is already covered by the
inherited noncommuting-residue theorem. This leaves exactly the
ordinary/small testing strip stated in the theorem. The exact gauges
$(u,W,X)\mapsto(-u,-W,-X)$ and
$(v,W,Y)\mapsto(-v,-W,-Y)$ preserve every event square and the core.
All homogeneous formulas allow $p=0$ or $w=0$ literally.
Real and imaginary state components are treated separately.
\end{proof}

The ordinary/small testing strip is completed in Section~\ref{component:pa0b854a04f49}. Together with the remote charts proved here, it covers the entire actual UI source. The actual UH source is treated by the separate complete conductor in Section~\ref{component:pf7ffa921793c}.


\section{\texorpdfstring{The centered-$A$ certificate for every bounded positive $\delta$}{The centered-A certificate for every bounded positive delta}}
\label{component:pccaa793fd0bf}


\begin{theorem}
Use the unchanged physical native C1 source with $k,\zeta,r=1-a\eta$
positive infinitesimals and $\eta$ a positive bounded unit. Suppose
$e=k\sqrt{\eta/\zeta}$ is bounded below by a positive real number;
it may be unbounded. Then the entire original parent is tensor SOS at
every independent projective test and every real or complex state.
Every positive opening rate, literal $\beta=0$, and every rate of
$k^2/r$ is included.
\end{theorem}

Use all exact definitions and the centered positive packet (CA1) in
Section~\ref{component:p3763434080c2}. The sole change is that
$\delta=\Delta R/K$ is now any positive bounded coefficient, so it
must not be inverted in a bounded Gram. The source chart (CA3),
$\ell=Qr=K-aE$, and the exact cancelled coquotients remain unchanged.
Their only denominators are positive units, independent of the rate
of $\ell$ relative to $k$.

\paragraph{Scale the entire remainder and its full normal.}
Let $\mathscr G$ be the exact full remainder (CA5). Put
\[
 \widehat G=\delta^2\mathscr G,\qquad
 \widehat M=\delta M_B,\quad \widehat Y=\delta Y,
 \quad\widehat W=\delta W.
\]
The original testing variables have not been rescaled against a
coefficient rate; these displayed quantities are explicit separated
tensor roots. Their complete formulas are
\[
 q_A=vD_uw,\quad q_B=v(D_u-uN-puD_uv w),
\]
\[
 \widehat W=\delta w+u(\widehat M+\delta hq_B)/D_u,
 \quad\widehat Y=\delta uw-(\widehat M+\delta hq_B)/D_u.
\]
Set
\[
 \widehat B_2=\frac{8a\widetilde c_2K^2}{R\widetilde c_3Q^2},
 \qquad
 \widehat B_b=\frac{8a\widetilde c_bK^2}{R\widetilde c_3Q^2}.
\]
The entire exact scaled remainder is
\begin{align*}
\widehat G={}&\widehat B_2\{N^2-2q_AN+(1-p)q_A^2\}
 +\widehat B_bD_v\{N^2+2q_AN+pq_A^2\}\\
&+B_3D_v\{\widehat M^2+2\delta q_B\widehat M+\delta^2hq_B^2\}
 +B_5\{\widehat M^2-2\delta q_B\widehat M+\delta^2(1-h)q_B^2\}\\
&+4A_*D_uD_v(\widehat Y-r u\widehat W)^2
 +4LD_uD_v(2u\widehat W\widehat Y-r u^2\widehat W^2)\\
&+4\alpha D_u(v\widehat Y-\delta)^2
 +2dD_u\{(v\widehat Y-\delta)^2-\delta^2D_v\}.
 \tag{CS1}
\end{align*}
All coefficients are bounded; every occurrence of $\delta$ has
nonnegative exponent. There is no inverse of $\delta$, $\ell$, $k$,
$M$, $J$ or $t$. The inverse $1/a$ in $A_*$ is a positive unit.
Every first-normal and hard-normal term is retained in (CS1).

\paragraph{Uniform contact grading.}
Keep the exact roots $A_0,B_0,C_0,D_0,x,y,z,h_0,q_1,q_2$ and polynomial
$P$ in (CA6). Use the same weights
$\operatorname{wt}(\delta,k,E,\ell,t,u,w,v)=(0,1,0,1,0,1,1,2)$.
Weight zero means bounded here, not bounded away from zero.
The exact finite coefficient collection from (CS1) is
\[
 \begin{aligned}
 (\widehat G|_{N=\widehat M=0})_{\le6}
 &=\frac{4P}{E^4}x^2-\frac{16(k+u)}{E^2}xz+\frac8E T_4,\\
 (\widehat L_N)_{\le3}&=0,\\
 (\widehat L_M)_{\le3}
 &=\frac{8\{E^2+2E-2\ell(E+1)\}}{E^2}x.
 \end{aligned}                                                    \tag{CS2}
\]
The full hard block on $(N,\widehat M)$ has residue
\[
        \operatorname{diag}(16/E,4(E^2+4)/E).          \tag{CS3}
\]
These are precisely the old identities after the indicated scaling,
but their validity for an infinitesimal $\delta$ requires the
nonnegative-exponent assertion in (CS1), not cancellation of a
purported unit $\delta$.

The scalar is in
$(A_0,B_0,C_0)^2+(D_0^2)$ over the ring with $E$ and the displayed
positive units inverted. Both complete first normals are in
$(A_0,B_0,C_0)$. Direct monomial allocation in (CS1) proves these
memberships with nonnegative $\delta$ exponents. At zero normals,
$q_A=A_0D_u$, $\delta q_B=B_0D_u(1-puv w)$, and both
$\widehat Y,u\widehat W$ belong to $(B_0,C_0)$ with bounded-unit
coefficients. The sole constant from the $\alpha$ packet satisfies
$\alpha\delta^2/D_0^2=2a\widetilde c_\alpha/\widetilde c_3$ exactly.
Thus these are memberships of the entire rational expression, not merely
its finite collection. The unchanged monic
contact conductor lifts them without a $\delta$ inverse. Its contact
substitutions preserve the weights, even if $\delta$ is infinitesimal.
After (CS2), all remaining scalar weights are at least seven and all
remaining normal weights at least four, hence the lifted correction
coefficients have zero residue. The finite checker verifies complete lifts of the displayed low-weight
collection and explicitly rejects every negative exponent other
than the allowed positive-unit $E$ inverse.

The leading soft form is
\[
            \frac{4(E+4)}{E^2}x^2+\frac8E T_4.
\]
Use the strict two-pairing circuit on $T_4$. Its only hard coupling
is the full first normal $8(E+2)x/E$ on $\widehat M$, and
\[
 \det\begin{pmatrix}4(E^2+4)/E&4(E+2)/E\\
                    4(E+2)/E&4(E+4)/E^2\end{pmatrix}
                   =256/E^3>0.                       \tag{CS4}
\]
Together with $16/E$ on $N$, this gives the strict eight-root Gram.

\paragraph{Exhaustion and restoration.}
The independent outer/state lemma Section~\ref{component:p364e14238053}
uses the unit-$e$ hypothesis only to infer a positive lower bound on
$m/d$ in (CAO2). When $e$ is unbounded, $m/d$ is instead unbounded,
so that same lower bound and every clipping argument hold. Its
$c/d,A/d,D/d$ bounds, exact centered root identity, ordinary-$u$
certificate, unit-$w$ certificate, and literal projective endpoint
arguments do not contain $e$ or $\delta$. Thus it reduces every
remaining independent test to exactly the contact chart settled above.
Restore $\delta^2$ times every retained original packet, including
(CA1), and undo the positive pure coefficient factor $\delta^2$.
All factors and radicals remain separated between the two fields.
Complex states split into real and imaginary evaluations.

This closes the formerly missing infinitesimal-$r$, unbounded-$e$
slice, and incorporates the previous unit-$e$ theorem. Together with
the inherited unit-$r$ results, it covers all positive native $r$
when $\eta$ is a unit and $e$ is bounded below. It does not yet
claim the infinitesimal-$e$ slice.


\section{\texorpdfstring{The outer and unit-state join for infinitesimal $e$}{The outer and unit-state join for infinitesimal e}}
\label{component:p05a86be0046c}

\label{p05a86be0046c:sec:small-e-outer}

\begin{lemma}
Use the unchanged physical native C1 source with $c=2$,
$\lambda=n=0$, positive original openings, and nonnegative original
event coefficients. Suppose $k,\zeta,e=k\sqrt{\eta/\zeta}$ are positive
infinitesimals, $\eta$ is a positive bounded unit, and
$r=1-a\eta>0$. Then every independent projective test and state is
tensor SOS except possibly the chart
\[
 u,v,w>0\text{ infinitesimal},\qquad Z=1,
 \qquad w=(W+uY)/(1+u^2).
 \tag{SEO1}
\]
There is no restriction on the positive relative scales of $a,r,J$
or on $k^2/r$. This lemma supplies a join, not a certificate for the
remaining chart (SEO1).
\end{lemma}

\paragraph{The small weak-event scale.}
Put $\kappa=k^2$, $K=1+\kappa$, $H=r+\eta$,
$V_c=\kappa\eta\zeta$, $S=1-\kappa\eta+tV_c$,
$R=1+\kappa r$, $T=K-\kappa\zeta t$,
$M=R-t$, and $J=T-a$. Let $P_4,P_d$ be the short native
polynomials printed in the centered-A outer lemma. All of
$K,H,S,P_4,P_d$ are positive bounded units. Set
\[
 L=KHS,\quad N=KH(1+V_c)-J\eta P_d,\quad
 d_0=4\eta P_d,\quad A_0=2\eta^2P_4,\quad D_0=2N.
\]
After exact cancellation of the native $r$ factors, the source is
\[
 d=\frac{ad_0}{L},\qquad A=\frac{aA_0}{rL},\qquad
 D=\frac{aD_0}{JL},\qquad
 \frac{m}{ae^2}=
 \frac{2(1-t\zeta)(1+V_c)}{MS}.                 \tag{SEO2}
\]
Here $\bar N=1$. Native $\beta\ge0$ gives the exact bound
\[
 \frac{t}{e^2}\le\frac{2c_b}{\eta c_d}
 =\frac{K[1+(1-t)\zeta]H}{\eta P_d}.
\]
Thus $t$ is infinitesimal, $M\sim1$, and $m/(ae^2)\sim2$.
In particular
\[
 \frac{k^2}{e^2}=\frac\zeta\eta
 \quad\hbox{is infinitesimal}.                       \tag{SEO3}
\]
The ratios $d/a,A/a,D/a$ have positive real lower bounds, while
$d/a$ is a bounded unit. These statements allow arbitrary positive
opening rates and literal $\beta=0$.

\paragraph{Choose the critical parameter without a small denominator.}
There are two source cases. If $a$ is a bounded unit, keep $A,D,d$
and replace $c=2$ by the critical number
\[
 c_* =\frac{d(d+2A)(d+2D)}{2(4AD-d^2)}.
 \tag{SEO4}
\]
This is positive and bounded, and $c_*-2=O(k^2)$.
Here is an exact short verification of the potentially delicate bound.
Define
\[
 Q=4A_0D_0-rJd_0^2,
\]
so $\bar Q=16\bar\eta^2\bar a\bar H>0$. Then
\[
 c_*-2=
 \frac{ad_0(d_0r+2A_0)(d_0J+2D_0)-4QL}{2QL}.
 \tag{SEO5}
\]
After $r=1-a\eta$ is inserted, its polynomial numerator vanishes
identically at $\kappa=0$, with $\zeta$ still free. It is therefore
$\kappa$ times a polynomial in bounded native variables; its
denominator is a positive unit. No inverse of $r$ or $J$ occurs.

If $a$ is infinitesimal, then $r,J\sim1$. Divide the whole parent
by $a$ and replace the original $A$ by
\[
 A_* =\frac{d^2(4+d+2D)}{2(8D-2dD-d^2)}.       \tag{SEO6}
\]
The denominator is positive. In this case
$A_*/a,D/a,d/a$ are positive bounded units, $c/a=2/a$ is
unbounded, and $(A_*-A)/a=O(k^2)$. To verify this without a
rate comparison, put
\[
 B_0=8D_0L-a(2d_0D_0+d_0^2J).
\]
Its residue is $16\bar H>0$, and
\[
 \frac{A_*-A}{a}=
 \frac{r d_0^2\{4JL+a(d_0J+2D_0)\}-2A_0B_0}
 {2rLB_0}.                                           \tag{SEO7}
\]
Again the numerator vanishes identically at $\kappa=0$, so it has
an exact polynomial factor $\kappa$, with a unit denominator.

Use $\sigma=1$ in the first case and $\sigma=a$ in the second,
and henceforth normalize all parent coefficients by $\sigma$.
Let $\epsilon$ be the positive part of the critical increase in the
chosen normalized coefficient. From (SEO2)--(SEO7),
\[
 \epsilon\text{ is infinitesimal},\qquad
 \frac{\epsilon}{m/\sigma}=O(\zeta)\text{ is infinitesimal}.
 \tag{SEO8}
\]
If the chosen critical value is smaller than the actual coefficient,
retain the positive excess packet instead. In either source case
the chosen critical parent obeys
\[
 2c(4AD-d^2)=d(d+2A)(d+2D),                  \tag{SEO9}
\]
with the letters in (SEO9) now denoting the normalized critical
coefficients. In the first case $c,d$ are units and $A,D$ are
bounded below by positive units; in the second $A,D,d$ are units
and $c$ may be unbounded.

\paragraph{A critical-wall frame permitting the large coefficients.}
Write $D_u=1+u^2$, $D_v=1+v^2$ and
\[
 B=X+uZ,\quad U=uW,\quad V=vZ,\quad J_t=U-Y,
 \quad C=4D+2d,
\]
and define the pure-coefficient numbers
\[
 y=\frac d{2A},\quad
 b=\frac{4AD-d^2}{AC},\quad s=1-b,\quad q=y+b.
\]
The critical identity gives the useful exact relations
\[
 b=\frac C{4c+C},\quad Cs=4cb,\quad
 cb=\frac{d(2A+d)}{4A}.                         \tag{SEO10}
\]
In both source cases, $y,b,s,q$ are bounded and $b,s>0$.
The entire signed core plus the three critical packets is exactly
\begin{align*}
 F_*={}&4cD_vJ_t^2+4AD_uD_vY^2
       +C(J_t+uvB-D_uV)^2+2dD_vB^2\\
       &+2dD_u\{v^2Y^2-2YV-V^2\}.              \tag{SEO11}
\end{align*}

In the first source case use
\[
 H_A=\sqrt A(Y-yV),\qquad
 H_D=\sqrt C(J_t+uvB-(b+u^2)V).
 \tag{SEO12}
\]
The inverse is $Y=yV+H_A/\sqrt A$ and
$J_t=(b+u^2)V-uvB+H_D/\sqrt C$.
All coefficients in (SEO11) after this substitution are bounded.
Indeed the large $A,C$ occur only in the normalized squares and in
the bounded combinations $Ay^2$, $\sqrt A\,y$, $Cs^2$,
$\sqrt C\,s$; use $Cs=4cb$ and bounded $c$ for the last two.
At $u=v=0$, the hard block on $(B,H_A,H_D)$ is
\[
 \operatorname{diag}(2d,4,1+4c/C)\succ0.
\]
On its exact zero section the scalar coefficient of $V^2$ is
\[
 g u^2+\gamma v^2+4cu^4+g(y+2)u^2v^2+4cu^4v^2,
 \quad g=\frac{d(2A+d)}A,\quad
 \gamma=\frac{Dd(2A+d)^2}{A^2(2D+d)}.             \tag{SEO13}
\]
Both $g,\gamma$ are positive bounded units. Every hard-$V$ entry
has angular degree at least two, with bounded coefficients; this is
an exact full-parent polynomial statement. Allocate it on $uV,vV$,
leaving a factor $u$ or $v$ in its coefficient. The resulting exact
five-root Gram on $(B,H_A,H_D,uV,vV)$ therefore has strict residue.

In the second source case set $L_c=4c+C$ and instead use
\[
 H_A=\sqrt A(Y-yV),\qquad
 H_c=\sqrt{L_c}(J_t-bV).                         \tag{SEO14}
\]
Its coefficients are bounded because $A,D,d$ are units and
$c/L_c$, $cb/\sqrt{L_c}$ and $cb^2$ are bounded by (SEO10).
Its hard block is $\operatorname{diag}(2d,4,1)$.
The exact zero-section scalar is
\[
 g u^2+\gamma v^2+C u^4+gy\,u^2v^2.
 \tag{SEO15}
\]
The same complete hard-$V$ allocation gives a strict five-root
Gram on $(B,H_A,H_c,uV,vV)$. Thus this is a full critical-wall
certificate in both source cases, even with the stated unbounded
coefficients. Every radical and denominator used in the roots is
pure coefficient; no mixed inverse is used.

\paragraph{All unit-$W$ states at the small-angle corner.}
Normalize $W=1$, with $Z,Y,X$ bounded, and take $u,v$ infinitesimal.
Retain the complete original $\alpha,\beta$ packets. Put
\[
 R_m=v-X+uvY-uZ=v(1+uY)-B,\qquad
 V=\frac{Z}{1+uY}(R_m+B).                        \tag{SEO16}
\]
The displayed denominator is a positive pure-testing unit.
In the first source case the only negative correction to the
critical parent is $-4\epsilon D_vJ_t^2$; (SEO12) and (SEO16)
express $\sqrt\epsilon J_t$ as a bounded combination of the five
critical roots and $\sqrt{m/\sigma}\,R_m$, with every coefficient
having zero residue. In the second case the only negative
correction is $-4\epsilon D_uD_vY^2$; use
$Y=yV+H_A/\sqrt A$ and (SEO16) for the identical conclusion.
The only possibly delicate coefficient is a bounded testing
multiple of $\sqrt{\epsilon/(m/\sigma)}$, and it is infinitesimal
by (SEO8). The extra original $m$ square contributes the strict
diagonal $4$. Hence the complete perturbed Gram remains strict.
This includes unit or infinitesimal $Z$, literal $Z=0$, and both
signs of the bounded testing state. It uses the entire original
weak $m$ event, without a false positive-unit lower bound on $m$.

\paragraph{Finite outer angles without the weak event.}
The same critical parent also gives the two finite outer charts.
In (SEO11) regard $B,Y,J_t,V$ as independent roots. Exactly,
\[
 F_*(0,v)=F_*(0,0)
       +v^2\{2dB^2+(4A+2d)Y^2+4cJ_t^2\}.       \tag{SEO17}
\]
The zero-angle form has three positive hard pivots and null ray
$(B,Y,J_t,V)=(0,y,b,1)V$. On this null ray the additional
coefficient is $\gamma v^2>0$. The normalized hard frames
(SEO12) or (SEO14) keep every coefficient bounded. Thus when
$u$ is infinitesimal and $v$ a bounded unit the full four-root
residue is strict. The infinitesimal critical-parameter deficit
is absorbed directly on these four roots; retain the entire
$m,\alpha,\beta$ packets.

For bounded-unit $u$ and infinitesimal $v$, set $V'=D_uV$.
The other exact identity is
\[
 F_*(u,0)=F_*(0,0;V')
        +u^2\{4AY^2+2d(V')^2/D_u\}.             \tag{SEO18}
\]
Use the same normalized critical roots with $V'$ in place of $V$.
In the first source case include $uvB$ in the $H_D$ root so that
the unbounded $D$ event remains isolated for nonzero small $v$.
The extra null mass in (SEO18) is positive in residue because
$d$ and $u$ are units. The bounded full four-root Gram is strict,
and the deficit is again a zero-residue correction. These two
identities include every state, not merely the normal restriction.

For completeness the same argument covers both finite angles of
nonzero residue, without a separate physical-source admission.
The exact determinant of (SEO11) on $(B,Y,J_t,V)$ is
\[
 64cdD_u^2D_v^2
 \{A(2D+d)u^2+D(2A+d)v^2\}.                    \tag{SEO20}
\]
The determinant in the normalized frame (SEO12) is this number
divided by $AC$, namely
\[
 32cdD_u^2D_v^2
 \left\{u^2+\frac{D(2A+d)}{A(2D+d)}v^2\right\}.
\]
It is a positive unit whenever at least one finite angle has
nonzero residue. The hard restriction $V=0$ is explicitly
\[
 \left(4D_uD_v+\frac{2dD_uv^2}{A}\right)H_A^2
 +H_D^2+4cD_v(H_D/\sqrt C-uvB)^2+2dD_vB^2,
\]
which has a strict bounded residue. Hence its full residue Gram
is positive definite by (SEO20). In frame (SEO14), the determinant
is (SEO20) divided by $A(4c+C)$. Here $A,D,d$ are units and
$c/(4c+C)\sim1/4$, so it too is a positive unit. Its hard
restriction is
\[
 \left(4D_uD_v+\frac{2dD_uv^2}{A}\right)H_A^2
 +\frac{4cD_v}{4c+C}H_c^2
 +C(H_c/\sqrt{4c+C}+uvB)^2+2dD_vB^2,
\]
again strict. All entries of both full transformed Grams are
bounded as already proved. The infinitesimal deficit is a
zero-residue correction on these four roots. This settles every
finite angle pair outside the double-infinitesimal corner.

\paragraph{Projective infinity charts.}
For these charts normalize by $a$, clip the positive coefficients
$c/a,A/a,D/a$ at suitable positive real lower bounds, and retain
their positive original excess squares. Retain the whole original
$m,\alpha,\beta$ packets. The normalized $d/a$ is a unit.
All the infinity-infinity and infinity-zero four-root identities
in (BE2) and the subsequent outer routing apply without $m$.
For $u=\infty$ in residue and $v$ a bounded unit, the signed-core
residue itself retains
$2d\{(W+vX)^2+(vY-Z)^2\}$; together with $cW^2,AY^2$ it is
strict, so no $m(vY-Z)^2$ lower bound is needed.
For $v=\infty$ and $u$ a bounded unit the old signed-core and
$c,D$ root argument is unchanged.
At the remaining $(u,1/v)=(0,0)$ corner, replace the obsolete
state-$W$ mass argument by the exact root ideal
$(X,Y,U=uW,Z)^2$ after retaining $m$. In $q^2F(u,1/q)$ its
complete selected residue is
\[
 2dX^2+4c(U-Y)^2+(4A+2d)Y^2+4DZ^2\succ0.
 \tag{SEO19}
\]
Every remaining occurrence of $W$ carries $u$, so this is a full
bounded Gram including all states. Literal infinity is included
by these homogeneous identities. Literal zero axes have their
unchanged original-event SOS. Finite tests with both noncommuting
residues were included directly in (SEO20).

\paragraph{The residual state chart.}
Only $u,v$ infinitesimal remain. Normalize the state by a maximal
coordinate. On a negative pure-testing-core branch, $X,Y$ have
zero residue. The unit-$W$ chart was just proved for every bounded
$Z$. Otherwise $Z$ is a unit, so normalize $Z=1$ and $W$ is
infinitesimal. Then $w=(W+uY)/D_u$ is infinitesimal. The inherited
exact testing Schur identity and sign gauges give (SEO1).
Complex states split into their real and imaginary evaluations.
The exact checker accompanying this lemma verifies its source
cancellations, critical roots, full hard-normal orders, both finite
outer identities, and the repaired projective-corner identity.


\section{\texorpdfstring{Small $e$: a bounded full-parent contact certificate for every $r$}{Small e: a bounded full-parent contact certificate for every r}}
\label{component:pa7e744614f49}


\begin{proposition}[Local part of the small-$e$ theorem]
Use the unchanged physical native C1 source, with $k,\zeta,e$ positive
infinitesimals and $\eta$ a positive bounded unit. Every native positive
$r=1-a\eta$, and every positive opening rate, is allowed.
At every independent test with $u,v,w$ infinitesimal, $Z=1$, and $w>0$,
the entire original parent is tensor SOS. Literal $\beta=0$ is included.
\end{proposition}
This proposition is the contact part only. The independent outer/state
join is supplied separately; no assertion here identifies a testing rate
with a source rate.

\paragraph{The source entry forced by the original sign.}
Keep $K,Q,E,\ell,\delta$ from (CA3), with $E=b^2$ and
$\ell=Qr=K-aE$. Here $E$ is a positive unit, $0<\ell<K$, and no lower
bound on $a$ or $\ell$ is assumed. Put
\[
 \varepsilon=1/\delta,\qquad j=\delta k,\qquad t=\varepsilon^2\tau.
\]
Both $\varepsilon,j$ are infinitesimal: $e=\varepsilon K/Q$ and
$\zeta=j^2QE/K$. The original condition $\beta\ge0$ gives
\[
 \frac{t}{e^2}\le\frac{2c_b}{\eta c_d},
\]
whose right side is bounded after cancelling the common native $r$
factor. Thus $\tau$ is positive bounded. The exact chart is
\[
 k=\varepsilon j,\quad K=1+\varepsilon^2j^2,\quad
 Q=K^2-\ell\varepsilon^2j^2,\quad a=(K-\ell)/E,
\]
\begin{align*}
 \eta&=KE/Q,\quad r=\ell/Q,\quad H=r+\eta,\quad
 e^2=\varepsilon^2K^2/Q^2,\\
 \zeta&=j^2QE/K,\quad R=K^2/Q,\quad T=K-k^2t\zeta. \tag{SE1}
\end{align*}
All denominators displayed here are positive units. The cancelled
coquotients $\widetilde c_i=c_i/\ell$ are the same short exact formulas
as in (CA3); the cancellation introduces no inverse of $\ell$.
Their residues are $2(E+\ell)$ for $i=2,b,3,5,\alpha$, and $4$ for
$i=d$; $c_4$ has residue $2$. These formulas hold uniformly as
$\ell$ tends to zero or to one.

\paragraph{One bounded centered remainder, including both normals.}
Set
\[
 A_* =\frac{ArH}{1+rH},\qquad
 A-A_* =\frac{A}{1+rH}>0.                              \tag{SE2}
\]
Retain the entire square $4(A-A_*)D_uD_v(Y-r uW)^2$.
It and its full first derivative vanish at exact native contact.
Divide the remaining parent by the positive pure coefficient $a$;
every factor $a$ cancels before coefficients are used. Write
\[
 \begin{array}{ll}
 b_2=8\widetilde c_2e^2/(R\widetilde c_3),&
 b_b=8\widetilde c_be^2/(R\widetilde c_3),\\
 b_3=8/T,&b_5=8\widetilde c_5/(T\widetilde c_3),\\
 L=Ar/a=2\eta^2c_4/(Q\widetilde c_3),&
 d_0=d/a=2\eta\widetilde c_d/\widetilde c_3,\\
 A_0=A_*/a=LH/(1+rH),&
 \alpha_0=\alpha/a=2\varepsilon^4j^6Q^2\eta^2
       \widetilde c_\alpha/(K^2\widetilde c_3).
 \end{array}                                                     \tag{SE3}
\]
Every quantity in (SE3) is bounded with unit denominators, and
$b_2,b_b$ contain $\varepsilon^2$.
Use the original normals $N,M_B$, put $p=t/R$, $h=a/T$, and write
\begin{align*}
 q_A&=vD_uw,\qquad q_B=v(D_u-uN-puD_uv w),\\
 W&=w+u(M_B+hq_B)/D_u,\qquad Y=uw-(M_B+hq_B)/D_u.
\end{align*}
The entire bounded remainder is
\begin{align*}
G_s={}&b_2\{N^2-2q_AN+(1-p)q_A^2\}
 +b_bD_v\{N^2+2q_AN+pq_A^2\}\\
 &+b_3D_v\{M_B^2+2q_BM_B+hq_B^2\}
 +b_5\{M_B^2-2q_BM_B+(1-h)q_B^2\}\\
 &+4A_0D_uD_v(Y-r uW)^2
 +4LD_uD_v(2uWY-r u^2W^2)\\
 &+4\alpha_0D_u(vY-1)^2
 +2d_0D_u\{(vY-1)^2-D_v\}
 +\nu D_vN^2,\qquad \nu=d_0(R-t)/R.                 \tag{SE4}
\end{align*}
This is (SD4), after the indicated exact centered subtraction and
normalization, except that a positive part of the original $\Pi_2$ is
kept inside the remainder. Indeed the original coefficient of
$D_vN^2$ in $\Pi_2/a$ is
\[
 \frac{4\beta+2d}{a}\frac{R-t}{R}\ \ge\ 2\nu.
\]
Thus subtracting only $\Pi_2/a-\nu D_vN^2$ remains legitimate, even
when $\beta=0$. The other three original packets are retained unchanged.
All full first normals and every hard quadratic term are present in
(SE4); there are no terms of higher normal degree.

\paragraph{A conductor that never divides by $\varepsilon$.}
Define the separated monomials and differences
\[
 A=\varepsilon wv,\quad B=v,\quad C=Euw,\quad
 D_0=\varepsilon^2j^3E,\quad k=\varepsilon j,
\]
\[
 x=B-C,\quad y=A-kC,\quad z=uB-kC,\quad h_0=kC-D_0.
\]
The two standard radical pairings are obtained by dividing the original
four squared circuit monomials by $\delta$: they are
$A,uB,kC,D_0$, and their product is $(k^2BC)^2$.
Their strict circuit is
$T_4=y^2+z^2+h_0^2-2yz+2q_1^2$.
The exact contact is $(u,w,v)=(k,j,\varepsilon j^2E)$.
Besides the three basic conductor identities,
\[
 C(u-k)=z-ux,\quad
 \varepsilon C(w-j)=y-\varepsilon wx,\quad
 C(v-\varepsilon j^2E)=E(wz+jy-kwx),                \tag{SE5}
\]
use the additional, division-free identity
\[
 C v(w-j)=Ew(w-j)z+Ejw y-Ekw^2x.                    \tag{SE6}
\]
For a monic contact division
$S=(u-k)T_u+(w-j)T_w+(v-\varepsilon j^2E)T_v$,
whenever $T_w=\varepsilon U+vV$, the bounded conductor is
\[
 \begin{aligned}
 CS={}&x[-uT_u-\varepsilon wU-Ekw^2V-EkwT_v]\\
 &+y[U+EjwV+EjT_v]
 +z[T_u+Ew(w-j)V+EwT_v].
 \end{aligned}                                      \tag{SE7}
\]
No division by $\varepsilon,j,k,r,a,t$, or a contact difference occurs.

Here is the full rational membership argument required for applying
(SE7), not only a finite truncation. At zero normals let
$V_0=1-puv w$. Then
\[
 q_A=D_u wv,\quad q_B=D_uvV_0,\quad
 uW=C/E+u^2hvV_0,\quad Y=C/E-hvV_0.
\]
The scalar of (SE4) therefore has an exact raw allocation
\[
 f_AA^2+f_{BB}B^2+f_{BC}BC+f_{CC}C^2+\rho D_0^2.     \tag{SE8}
\]
Expand only the nonconstant part of $(vY-1)^2$ into $B,C$; its constant
has $\alpha_0/D_0^2=2\widetilde c_\alpha/\widetilde c_3$.
The coefficients in (SE8) have unit denominators independent of $w$.
Every dependence on $w$ in them comes through $V_0$, and hence has a
factor $\varepsilon^2$; modulo $\varepsilon$ they are independent of
$w$ and even in $u,v$. Both complete first normals are in $(A,B,C)$
with bounded coefficients by the same exact formulas.

Clear the finitely many unit denominators, chosen even in $u,v$ and
independent of $w$, and use the unchanged scalar square-lift algorithm.
For its first scalar conductor, the only $w$-dependent coefficient
$T_w$ is divisible by $\varepsilon$ (in fact by $\varepsilon^2$).
At $\varepsilon=u=0$, its residual polynomial $S$ is even in $v$,
and contact makes its constant zero; thus $S\in(v^2)$ there.
More explicitly, modulo $\varepsilon$ the first divided differences
satisfy $T_u\in(u)$, $T_w=0$, and $T_v\in(v)$ after the preceding
monic substitutions. It follows
that after this first conductor, the sole possible $w$ dependence
modulo $\varepsilon,u$ in the four second-stage remainders is of the
form $w v$ times a polynomial in $v$. Hence their $T_w$ all belong to
$(\varepsilon,v)$, exactly the hypothesis of (SE7). The complete
normal raw coefficients are independent of $w$ modulo $\varepsilon$,
so their one conductor has $T_w\in(\varepsilon)$. The retained contact
squares and the original native balances supply the exact contact
value and first-jet identities throughout. This proves the full
scalar-square and complete first-normal lifts over the bounded ring.

\paragraph{Finite residues of that full lift.}
Give $(\varepsilon,j,u,w,v,N,M_B)$ weights $(1,1,2,1,3,5,3)$;
$E,\ell,\tau$ have weight zero. Every positive-weight variable is
infinitesimal, while all weight-zero coefficients are bounded.
The contact substitutions preserve the grading; $x$ has weight three
and the other circuit roots weight five. Both (SE5) and (SE6) preserve
weights, so after the scalar collection through weight ten and full
first normals through weight five, every lifted remainder coefficient
has positive weight. Restore the unit denominators afterward.

Set $H_0=E+\ell$ and $V_*=1+\ell H_0$. Direct exact multiplication of
(SE1)--(SE4), followed by (SE7), gives scalar residue
\[
       \frac{8(E+2)}{H_0V_*}x^2+8T_4,                         \tag{SE9}
\]
and complete first-normal residues
\[
       L_N=0,\qquad L_M=\frac{16E(H_0+1)}{H_0V_*}\,x.
\]
The hard block has residue
\[
 \operatorname{diag}\left(\frac{4E}{H_0},
           \frac{8(E^2+2E\ell+2\ell^2+2)}{V_*}\right).       \tag{SE10}
\]
These statements include coefficient corrections of positive weight;
they are not asserted as literal equalities of the full rational scalar.
The exact checker reconstructs (SE4), verifies the full finite lifts,
and compares their residues with (SE9)--(SE10).
The only hard-soft coupled determinant is
\[
 \det G_{M,x}
 =\frac{64(E^2+2E\ell+4E+4\ell)}{H_0^2V_*}>0.              \tag{SE11}
\]
Together with $4E/H_0>0$ and the strict two-pairing circuit this gives
the full bounded strict eight-root Gram. All original positive packets,
the centered square, and the positive pure-coefficient factor $a$ are
then restored. All inverses and radicals remain separated between the
two independent real closed fields.


\section{\texorpdfstring{Completion of the unit-$\eta$, infinitesimal-$e$ source domain}{Completion of the unit-eta, infinitesimal-e source domain}}
\label{component:p29aeda9b68a4}


\begin{theorem}[Small $e$, every positive native $r$, all independent tests]
Use the unchanged physical native C1 source with $c=2$, $\lambda=n=0$,
all original event coefficients nonnegative, and all original openings
$a,t,M,J$ positive. Suppose $k,\zeta$ are positive infinitesimals,
$\eta$ is a positive bounded unit, $e=k\sqrt{\eta/\zeta}$ is a positive
infinitesimal, and $r=1-a\eta>0$. Then the entire original parent is
tensor SOS at every independent projective test and every real or
complex state. This includes every relative rate of $k^2/r$, every
positive opening rate, and literal $\beta=0$.
\end{theorem}

\begin{proof}
The exact original $\beta$ sign supplies the bounded opening scale
$t/e^2$ in (SE1) and (SEO2), without a lower bound on $t$.
The outer/state lemma Section~\ref{component:p05a86be0046c} proves every
independent test outside the chart $Z=1$ with $u,v,w>0$ infinitesimal.
Its critical-source deficits have bounded exact coquotients after a
single factor $k^2$ is removed. The weak original $m$ event absorbs
those deficits because $k^2/e^2=\zeta/\eta$ is infinitesimal.
The two source cases $a$ unit and $a$ infinitesimal exhaust all positive
native $a$, and the homogeneous endpoint identities include every
projective corner and every state normalization.

On the remaining chart, the proposition in
Section~\ref{component:pa7e744614f49} gives the bounded full-parent
remainder (SE4). Its centered cap is uniform in $r$; the reserved
positive part of $\Pi_2$ keeps the full $N$ pivot positive even at
$\beta=0$. The division-free conductor (SE5)--(SE8) applies to the
entire rational scalar and both complete first normals. Its finite
residues (SE9)--(SE11) give the strict eight-root Gram. Restore all
retained original packets, the centered square, and the positive pure
coefficient normalization. All testing and source inverses and radicals
remain separated. The real and imaginary evaluations add to give the
complex-state assertion.
\end{proof}

Together with the bounded-positive-$\delta$ theorem and the inherited
unit-$r$ results, this completes every positive $e$ and every positive
native $r$ when $k,\zeta$ are infinitesimal and $\eta$ is a positive
bounded unit. It makes no assertion about infinitesimal or unbounded
$\eta$, or about another native C1 source branch.


\section{\texorpdfstring{Bounded unit native state scale and bounded unit $e$}{Bounded unit native state scale and bounded unit e}}
\label{component:p69995d829bb2}


\begin{theorem}
Use the unchanged physical native C1 source, $c=2$, $\lambda=n=0$,
with all original event signs and positive openings $a,t,M,J$.
Suppose $k,\zeta$ are positive infinitesimals, and
\[
 \eta,\quad r=1-a\eta,\quad e=k\sqrt{\eta/\zeta}
\]
are positive bounded units. Then the complete original parent is tensor
SOS at every independent projective test and every real or complex state.
The openings may be arbitrarily small, and $\beta=0$ is included.
\end{theorem}

\paragraph{The double-infinity corner.}
With $p=1/u$, $q=1/v$, the exact homogeneous parent is
\begin{align*}
F_{\infty\infty}={}&4m(pW-pqX+Y-qZ)^2
 +4c(1+q^2)(W-pY)^2+4A(1+p^2)(1+q^2)Y^2\\
&+(4D+2d)(qW+X-pqY-pZ)^2
 +(4\alpha+2d)(1+p^2)(Y-qZ)^2\\
&+(4\beta+2d)(1+q^2)(pX+Z)^2
 -2d(1+p^2)(1+q^2)Z^2.
\end{align*}
At $Z=1$, after retaining $\alpha,\beta$, its complete quadratic part is
\[
 4cW^2+2dX^2+4D(X-p)^2+4AY^2+(4m+2d)(Y-q)^2.
 \tag{BE1}
\]
In particular the $(Y,q)$ determinant is $4A(4m+2d)>0$.
This is an exact full-parent expansion.

An even shorter full-parent certificate retains every event except
$c,A$. The remaining polynomial belongs to $(W,X,Y,qZ)^2$, with bounded
coefficients at $p,q=0$. Its residue Gram is the quadratic form
\[
 4cW^2+2dX^2+4AY^2+2d(Y-qZ)^2.                 \tag{BE2}
\]
It is strictly positive when $c/d,A/d$ are bounded below by positive
units. All cross terms of the complete polynomial are allocated on
these four roots, so (BE2) handles arbitrary states and both literal
projective endpoints, without a normal restriction.

\paragraph{Uniform lower bounds and outer tests.}
Write $H=r+\eta$. The short coquotients have residues $2rH$ for
$c_2,c_b,c_3,c_5,c_\alpha$, $2$ for $c_4$, and $4r$ for $c_d$.
The source identities imply
\[
 \frac md=\frac{e^2c_2}{M\eta c_d},\quad
 \frac cd=\frac{c_3}{a\eta c_d},\quad
 \frac Ad=\frac{\eta c_4}{c_d},\quad
 \frac Dd=\frac{2c_5-J\eta c_d}{2J\eta c_d}.
 \tag{BE3}
\]
Every ratio is bounded below by a positive real number. For the last,
$J+a=T\sim1$ and
\[
 \overline{2c_5-J\eta c_d}=4\bar r(\bar H-\bar J\bar\eta)
 =4\bar r>0.
\]
No reciprocal of an opening is required to be bounded. Choose positive
real lower bounds for the four ratios, cap these four coefficients at
the corresponding multiples of $d$, and retain their nonnegative
excess event squares. All selected-parent Gram coefficients divided
by $d$ are then bounded.

The physical-source theorem covers every noncommuting residue test.
For $u$ infinitesimal and $v$ a bounded unit, the selected residue form is
\[
 4m(vW-X)^2+2d(1+v^2)(X^2+Y^2)
 +4(c+A)(1+v^2)Y^2+4D(Y+vZ)^2\succ0.
\]
For $u=\infty$ in residue and $v$ a bounded unit, retain the signed
core at that commuting residue and use the four positive squares
$cW^2$, $AY^2$, $m(vY-Z)^2$, $D(W+vX)^2$; their common kernel is zero.
For $v=\infty$ in residue and $u$ a bounded unit, use the core
$2d(1+u^2)(X^2+Y^2)$ and the $c(uW-Y)^2,D(uX-Z)^2$ packets.
The $(0,\infty)$ corner has the strictly positive state residue
\[
 4mW^2+2dX^2+\{4(c+A)+2d\}Y^2+4DZ^2.
\]
The $(\infty,\infty)$ corner is (BE2).
At $(\infty,0)$ retain all events except $c,A$ and use the exact
four-root allocation on
\[
 (pX+Z,\,W,\,Y,\,vX),\qquad p=1/u.
\]
Its residue is
\[
 2d(pX+Z)^2+2d(W+vX)^2+4cW^2+4AY^2\succ0.
\]
The full selected polynomial lies in the square of this displayed
root ideal; the pure-testing change $Z=(pX+Z)-pX$ makes this immediate.
The other coefficients have zero residue. These homogeneous
identities include literal infinities. Literal zero axes have a
nonnegative pure-testing core and are immediate original-event SOS.
Only bounded $u$ and infinitesimal $v$ remain.

\paragraph{State entry and the $W$ endpoint.}
Normalize the state by a maximal coordinate. Set
$D_u=1+u^2$ and $w=(W+uY)/D_u$. On the negative pure-core branch,
if $w$ is infinitesimal then $Z$ is a unit: otherwise the residue core
$2\{(X+uZ)^2+(uW-Y)^2\}$ forces every normalized state coordinate
to have zero residue. Normalize $Z=1$; the exact pure-testing Schur
identity in (SD5) gives $w>0$ after the standard angle gauges.

If $w$ is a unit and $u$ is a unit, a zero of that core residue has
$Y=uw\ne0$. The positive lower bound for $A/d$ gives a strict state
residue after clipping the positive events. Thus only $u,v$
infinitesimal and $W\sim w$ a unit need an additional state check.
When $Z$ is infinitesimal, retain all events except clipped $m,A$ and
use the roots $(X,Y,uW,vW)$. Since $Z/W$ is a pure-testing
infinitesimal, the complete selected polynomial has residue
\[
 2d\{X^2+(uW-Y)^2\}+4m(vW-X)^2+4AY^2\succ0.
 \tag{BE4}
\]
Every remaining coefficient contains $u$, $v$ or $Z/W$; there is no
constant or first-normal term outside this frame. This includes
literal $Z=0$. If $Z$ is a unit, normalize $Z=1$ and use the next
raw scalar frame with $w$ a unit.

\paragraph{The exact remainder and its complete normals.}
Use exactly (SD3)--(SD8), their four positive packets, and the positive
clearing factor $\mathcal C$. Put
\[
 \delta=\Delta R/K,\quad E=b^2=\eta K/R,\quad
 A_0=wv,\ B_0=\delta v,\ C_0=\delta Euw,
 \ D_0=\delta^2k^3E,\quad x=B_0-C_0.
\]
Here $\delta,E$ are positive bounded units, and $a,t$ are bounded.
The factored five-term scalar is in
$(A_0,B_0,C_0)^2+(D_0^2)$, and the two complete first normals are
in $(A_0,B_0,C_0)$. These are the full formulas, without restricting
normals or replacing native coefficients by leading surrogates.

At $k=v=0$ and at the relevant small-$w$ or small-$u$ endpoint, put
$r_0=1-aE$ and $H_0=r_0+E$. Direct collection of the short formulas
gives the raw scalar residue
\[
 8r_0H_0D_u^2 A_0^2+
 8D_u\{x^2+u^2r_0H_0B_0^2\}+8r_0H_0D_uD_0^2.
 \tag{BE5}
\]
The hard block on $(N,\delta M_B)$ is
\[
 \operatorname{diag}(16r_0H_0,16r_0H_0+8E^2/D_u),
\]
and the complete first normals have residues $0,16Ex$.
The coupled $(\delta M_B,x)$ determinant is
$128D_ur_0H_0>0$.

If $u$ is a unit and $w$ infinitesimal, (BE5) is a strict Gram on
$(A_0,x,B_0,D_0,N,\delta M_B)$. If $u$ is infinitesimal and $w$ a
unit, use the exact relations
\[
 B_0=(\delta/w)A_0,\qquad C_0=(\delta/w)A_0-x.
\]
The factor $\delta/w$ is a product of a pure-coefficient unit and a
pure-testing unit; it is not a mixed inverse. The resulting residue
is the strict form
\[
 8r_0H_0A_0^2+8x^2+8r_0H_0D_0^2
 +16r_0H_0N^2+(16r_0H_0+8E^2)(\delta M_B)^2
 +16E(\delta M_B)x.                                  \tag{BE6}
\]
This supplies precisely the unit-$w$ chart for which the older
$\Delta$-small retained-current argument is not applicable.

For transparency, the only nontrivial raw coefficient collection in
(BE5) is
\[
 r_0H_0D_u+a^2E^2+r_0E(2a-1)=1+u^2r_0H_0,
 \qquad aE+r_0=1.
\]
The cross coefficient is $-16D_u(aE+r_0)=-16D_u$.
In the full normal formula $S_B$ has residue $-16E$ and $S_C$ has
residue $16E$. These prove the claimed full normal residues directly.

\paragraph{The sole contact chart.}
For $u,w,v$ infinitesimal and $Z=1$, use the unchanged monic
contact--conductor identities and the two-pairing circuit. Use weights
\[ (\delta,k,b,a,t,u,w,v):\quad(0,1,0,0,0,1,1,2). \]
The exact source collection is
\[
 (E_0)_{\le6}=8\rho T_4-8(\rho+E)k^2x^2,
 \quad (L_1)_{\le3}=0,\quad (L_2)_{\le3}=16Ex,
 \qquad \rho=r_0H_0-1.
\]
The source/contact substitutions preserve this grading. After the
displayed terms are removed, the scalar and normal lift coefficients
have positive weights, hence zero residues. The term $k^2x^2$ also
has a zero-residue coefficient. Thus the soft leading form is
$8x^2+8r_0H_0T_4$, with the same strict two-pairing circuit frame.
The hard block and its only first-normal coupling are those of (BE6),
and their determinant is $128r_0H_0>0$. Retain the old positive
$g_4^2$ square. The full eight-root remainder Gram is strict.
This invokes only the polynomial monic conductor already proved;
no opening or contact difference is inverted.

Finally restore the four complete positive packets, the retained
positive square, and all pure-coefficient and pure-testing factors.
The cases above exhaust both projective angle lines and every
projective state. Complex states split into their real and imaginary
evaluations. Neither this theorem nor the correction of (BE1) asserts
coverage when $r$ is infinitesimal, $\eta$ unbounded, or $e$
infinitesimal; those are separate source domains.


\section{\texorpdfstring{Outer tests and state entry when the native $r$ is infinitesimal}{Outer tests and state entry when the native r is infinitesimal}}
\label{component:p364e14238053}

\label{p364e14238053:sec:centered-a-outer}

\begin{lemma}
Use the unchanged physical native C1 source, normalized by $c=2$ and
$\lambda=n=0$, with all original event signs and positive openings
$a,t,M,J$. Suppose $k,\zeta,r=1-a\eta$ are positive infinitesimals and
$\eta,e=k\sqrt{\eta/\zeta}$ are positive bounded units.
Then the outer-test and state-entry arguments reduce the all-tests
problem to precisely
\[
 u,v,w>0\text{ infinitesimal},\qquad Z=1,
 \qquad w=(W+uY)/(1+u^2).
 \tag{CAO1}
\]
All other independent projective tests and states admit a full-parent
tensor SOS. No comparison of $k^2$ with $r$ is made, and arbitrarily
small positive openings and literal $\beta=0$ are included.
\end{lemma}

\paragraph{Source ratios with the common factor cancelled.}
Here $r$ always denotes the native source parameter, not the event root
$uW-Y$. Put
\[
 K=1+k^2,\quad H_c=r+\eta,\quad V_c=k^2\eta\zeta,\quad
 S_3=1-k^2\eta+tV_c,
\]
The two required short native polynomials are
\begin{align*}
P_4={}&1+k^2+\eta k^2(a-k^2-1)-\zeta k^2t\\
 &+V_c[-\zeta k^2t^2+2k^2t+a(1-t)+t],\\
P_d={}&(1+k^2)^2-a\eta k^4
 +\zeta k^2[(k^2+1)-(k^2+2)t]\\
 &+V_c[-\zeta k^2t^2+k^2a(t-1)+k^2t+k^2+1].
\end{align*}
They are positive bounded units with residue one. The exact source
identities, after cancelling the common factor $r$ in the coquotients,
are
\begin{align*}
 d&=\frac{4a\eta P_d}{KH_cS_3},&
 Ar&=\frac{2a\eta^2P_4}{KH_cS_3},\\
 \frac md&=\frac{e^2K(1-t\zeta)(1+V_c)H_c}{2M\eta P_d},&
 \frac cd&=\frac{KH_cS_3}{2a\eta P_d},\\
 \frac Ad&=\frac{\eta P_4}{2rP_d},&
 \frac Dd&=\frac{KH_c(1+V_c)-J\eta P_d}{2J\eta P_d}.
 \tag{CAO2}
\end{align*}
The positive openings obey
$M+t=1+k^2r$ and $J+a=1+k^2-k^2\zeta t$; hence $a,t,M,J$ are
bounded. Since $a\eta=1-r$, one has
\[
 \bar\eta\ge1,\qquad \bar d=4/\bar\eta,\qquad Ar\sim2,
 \qquad A\text{ positive unbounded}.
 \tag{CAO3}
\]
The numerator of $D/d$ in (CAO2) has residue
$\bar H_c-\bar J\bar\eta=1$. Therefore $m/d,c/d,D/d,A/d$ all have
positive real lower bounds. This conclusion uses no inverse bound on
an opening. Moreover,
\[
 \begin{cases}
 \displaystyle \bar D=\frac{2}{\bar\eta(\bar\eta-1)},
       &\bar\eta>1,\\[4pt]
 D\text{ positive unbounded},&\bar\eta=1.
 \end{cases}
 \tag{CAO4}
\]

\paragraph{The previously proved outer routing still applies.}
Clip $m,c,A,D$ at positive real multiples of $d$, retaining all their
nonnegative excess original event squares. The resulting coefficients
divided by $d$ are bounded. The outer-test argument in
Section~\ref{component:p69995d829bb2}, through (BE2)--(BE4), uses only
these lower bounds and the independent physical-source admission at
noncommuting residue tests. Its exact homogeneous identities and
four-root ideal memberships are unchanged. It therefore includes all
projective infinity charts, the $(0,\infty)$, $(\infty,0)$ and
$(\infty,\infty)$ corners, and the unit-$W$, infinitesimal-$Z$ state
chart. The original-event SOS on literal zero axes is unchanged as
well. This invokes none of the bounded-$r$ scalar or normal conclusions
(BE5)--(BE6). Only bounded $u$ and infinitesimal $v$ remain.

\paragraph{An exact centered root identity.}
Write $D_u=1+u^2$, $D_v=1+v^2$ and introduce
\[
 h=A^{-1/2},\quad L=Ar,\qquad
 H=\sqrt A(Y-ruW),\quad B=X+uZ,\quad U=uW,\quad V=vZ.
 \tag{CAO5}
\]
Thus $Y=rU+hH$, while $h,r$ are infinitesimal and $L\sim2$.
All coefficient radicals in (CAO5) lie in the coefficient field;
$B,U,V$ lie in the testing field. The root $H$ is their finite tensor
combination, with no mixed inverse or mixed radical.

For clarity, the complete original finite parent used here is
\begin{align*}
F={}&4m(vW+uvY-X-uZ)^2+8D_v(uW-Y)^2
       +4AD_uD_vY^2\\
 &+(4D+2d)(uW+uvX-Y-vZ)^2
   +(4\alpha+2d)D_u(vY-Z)^2\\
 &+(4\beta+2d)D_v(X+uZ)^2-2dD_uD_vZ^2.
\end{align*}

Choose $0<D_*\le D$ bounded, and let $m_*=0$ or choose a positive
real $m_*\le m$, as specified below. Retain the four nonnegative
original event contributions with weights $m-m_*,D-D_*,\alpha,\beta$.
For $C=4D_*+2d$, the exact remaining parent is
\begin{align*}
F_*={}&8D_v[(1-r)U-hH]^2
 +4D_uD_vH^2+4LrD_uD_vU^2+8LhD_uD_vUH\\
&+C[(1-r)U-hH+uvB-D_uV]^2+2dD_vB^2\\
&+2dD_u\{v^2(rU+hH)^2-2(rU+hH)V-V^2\}
 +4m_*R_m^2,                                      \tag{CAO6}\\
R_m={}&vW+uv(rU+hH)-B.
\end{align*}
In particular, the centered $A$ contribution is the exact identity
\[
 4AD_uD_vY^2
 =4D_uD_vH^2+4LrD_uD_vU^2+8LhD_uD_vUH.
\]
This retains its complete hard root. The last two coefficients have
zero residue. Formula (CAO6) follows directly from the original parent
and includes the complete signed core; it is not an eliminated scalar
or a restriction to state normals. Every displayed coefficient is
bounded. The term involving $Z$ in that core has already cancelled
to the displayed expression in $V=vZ$.

\paragraph{Every state when $u$ is a bounded unit.}
Set $m_*=0$, so the entire original $m$ event is retained. Formula
(CAO6) is an exact Gram allocation on $(H,B,U,V)$. At infinitesimal
$v$ its residue is
\[
 4D_uH^2+2dB^2+8U^2+C(U-D_uV)^2-2dD_uV^2.
 \tag{CAO7}
\]
Residues are understood in this and the following determinant
calculations. The $U,V$ determinant is
\[
 D_u\{8CD_u-2d(8+C)\}.
\]
If $\bar\eta>1$, take $D_*=D$, which is bounded by (CAO4).
Then $C=8/(\bar\eta-1)$ in residue and this determinant is
\[
 \frac{64u^2(1+u^2)}{\bar\eta-1}>0.
\]
If $\bar\eta=1$, choose a sufficiently large positive real $D_*$;
this is permitted because $D$ is unbounded. The determinant is then
\[
 32(1+u^2)(D_*u^2+2u^2-2)>0.
\]
The first diagonal is $8+C>0$, so the complete four-root residue
Gram is positive definite in both cases. The cap is a real constant
chosen from the nonzero testing residue of $u$; no division of tensor
elements is introduced. The bounded full Gram is therefore tensor
SOS, and restoring the retained events proves the claim for every
state in this chart.

\paragraph{The remaining unit-$w$ state chart.}
Suppose $u,v$ are infinitesimal and $w$ is a unit. Then $W$ is a unit
on the negative-core branch. When $Z$ is infinitesimal the inherited
(BE4) argument already applies. Otherwise normalize the state so that
$Z$ is a unit and put $q=W/Z$, a pure-testing unit. Choose a positive
real lower bound $m_*\le m$. Now
\[
 R_m=qV+uv(rU+hH)-B,
\]
so (CAO6) is again an exact bounded Gram on $(H,B,U,V)$, with residue
\[
 4H^2+2dB^2+4m_*(qV-B)^2
       +8U^2+C(U-V)^2-2dV^2.                    \tag{CAO8}
\]
The $B$ pivot is positive. Its Schur complement adds
\[
 \mu V^2,\qquad \mu=\frac{4dm_*q^2}{d+2m_*}>0
\]
to the $U,V$ block. This is just a residue determinant computation;
no mixed inverse is used to define a root or a full-parent coefficient.
If $\bar\eta>1$, take $D_*=D$ as before. The determinant becomes
\[
 \frac{8\bar\eta\mu}{\bar\eta-1}>0.
\]
If $\bar\eta=1$, choose a sufficiently large positive real cap
$D_*$ so that $\mu(C+8)-64>0$. The unboundedness of $D$ permits this
choice. Thus the complete residue Gram in (CAO8) is strict, including
all hard-soft terms, and restoring the retained events gives the
full-parent tensor SOS.

\paragraph{Completion of state entry.}
Normalize the state by a maximal coordinate and first remove the
nonnegative pure-testing-core branch, which has its direct original
SOS. On the negative-core branch with $v$ infinitesimal, the residue
core is
\[
 2\{(X+uZ)^2+(uW-Y)^2\}.
\]
Consequently these two roots have zero residue. If $w$ and $Z$ were
both infinitesimal, the exact inverse relations
\[
 W=w+\frac{u(uW-Y)}{D_u},\quad
 Y=uw-\frac{uW-Y}{D_u},\quad X=(X+uZ)-uZ
\]
would make every normalized state coordinate infinitesimal. Hence
small $w$ forces $Z$ to be a unit. The bounded-unit $u$ chart was
settled by (CAO7), and all unit-$w$ states by (CAO8) and (BE4).
Normalize the only remaining states by $Z=1$. The same exact
pure-testing Schur identity (SD5) and testing sign gauges give
$u,v,w>0$. This is exactly (CAO1). Literal axes and literal state
zeros have already been included. Complex states split into their
real and imaginary evaluations because the parent is real.



\section{\texorpdfstring{An infinitesimal native $r$ with unit $\eta$ and unit $e$}{An infinitesimal native r with unit eta and unit e}}
\label{component:p3763434080c2}


\begin{theorem}[The centered-$A$ unit-$e$ slice]
Use the physical native C1 source with its unchanged original events,
$c=2$, $\lambda=n=0$, all positive openings $a,t,M,J$, and all original
nonnegative signs, including literal $\beta=0$.
Suppose $k,\zeta,r=1-a\eta$ are positive infinitesimals, while
$\eta$ and $e=k\sqrt{\eta/\zeta}$ are positive bounded units.
Then the full parent is tensor SOS at every independent projective test
and every real or complex state. No bound on $k^2/r$ is imposed.
\end{theorem}

The theorem is restricted to unit $e$. Its proof uses $\delta^{-1}$
as a bounded pure coefficient number; it does not assert the cases of
infinitesimal or unbounded $e$.

\paragraph{Centering the sole large event.}
The exact native contact has $Y_c=r u_cW_c$. Consequently
\[
 L_A=Y-r uW
\]
vanishes at contact with no asymptotic replacement. The native formulas
give $a$ a positive unit, $A$ unbounded and $Ar$ a positive unit of
residue $2$. Put $A_*=1/a$, a positive bounded pure coefficient unit.
The complete packet
\[
             4(A-A_*)D_uD_v L_A^2                         \tag{CA1}
\]
is positive and has zero value and full first derivative at the exact
native contact. Retain it as a square.

It is essential to center this packet before taking residues.
The remaining $A$ channel is exactly
\[
 4A_*D_uD_vL_A^2+
 4(Ar)D_uD_v\{2uWY-r u^2W^2\}.                         \tag{CA2}
\]
Only $Ar$ and $r$ occur in its coefficients. In particular this is a
bounded expression for every positive rate of $r$ relative to $k$.

\paragraph{A source chart with no hidden $k^2/r$ quotient.}
Keep the exact contact variables $\delta,k,E=b^2$ from (SD6), and put
\[
 K=1+k^2,\quad \ell=K-aE=Qr>0,\quad
 Q=K^2-\ell k^2,\quad a=(K-\ell)/E.
\]
Here $\delta,E,a$ are positive bounded units and $k,\ell$ are positive
infinitesimals. The exact source substitution is
\[
 \eta=KE/Q,\quad r=\ell/Q,\quad
 \Delta=\delta Q/K,\quad \zeta=\delta^2k^2QE/K,
 \quad R=K^2/Q,\quad T=K-k^2t\zeta.                    \tag{CA3}
\]
All denominators in (CA3) are positive pure coefficient units.
For $i=2,b,3,5,\alpha,d$ let $\widetilde c_i=c_i/\ell$,
defined by cancelling the explicit $L_c=\ell/Q$ factor in the short
native formulas. This cancellation gives rational expressions with
unit denominators; it does not leave an inverse of $\ell$.
For example, writing $H_c=r+\eta$ and $V_c=k^2\eta\zeta$,
\[
 \widetilde c_3=\frac{2KH_cQ_c(1-k^2\eta+tV_c)}Q,
 \qquad \widetilde c_d=\frac{4Q_cP_d}Q.
\]
Their residues are $2E$ for $i=2,b,3,5,\alpha$, and $4$ for $i=d$;
$c_4$ has residue $2$.

Define the following bounded pure coefficient quantities:
\[
 \begin{array}{ll}
 B_2=\dfrac{8a\widetilde c_2}{R\widetilde c_3\Delta^2},&
 B_b=\dfrac{8a\widetilde c_b}{R\widetilde c_3\Delta^2},\\[4pt]
 B_3=8a/T,&B_5=\dfrac{8a\widetilde c_5}{T\widetilde c_3},\\[4pt]
 L=Ar=\dfrac{2a\eta^2c_4}{Q\widetilde c_3},&
 d=\dfrac{2a\eta\widetilde c_d}{\widetilde c_3},\\[4pt]
 \alpha=\dfrac{2a\Delta^2k^6\eta^2\widetilde c_\alpha}
                    {\widetilde c_3}.&
 \end{array}                                                     \tag{CA4}
\]
These formulas explicitly account for the cancellation needed below.

\paragraph{The exact full remainder, before any collection.}
Retain the four original (SD3) packets, divided by the positive source
factor $\Delta^2\Lambda$, and retain (CA1). Use the unchanged pure
testing coordinate $w=(W+uY)/D_u$ and the hard normals $N,M_B$.
For $Z=1$ put $p=t/R$, $h=a/T$ and
\[
 q_A=vD_uw,\quad q_B=v(D_u-uN-puD_uv w),
\]
\[
 W=w+u(M_B+hq_B)/D_u,\qquad
 Y=uw-(M_B+hq_B)/D_u.
\]
The complete remainder is exactly
\begin{align*}
\mathscr G={}&B_2\{N^2-2q_AN+(1-p)q_A^2\}
 +B_bD_v\{N^2+2q_AN+pq_A^2\}\\
&+B_3D_v\{M_B^2+2q_BM_B+hq_B^2\}
 +B_5\{M_B^2-2q_BM_B+(1-h)q_B^2\}\\
&+4A_*D_uD_v(Y-r uW)^2
 +4LD_uD_v(2uWY-r u^2W^2)\\
&+4\alpha D_u(vY-1)^2
 +2dD_u\{(vY-1)^2-D_v\}.                         \tag{CA5}
\end{align*}
There are no normal terms above degree two. Every coefficient of
(CA5) is bounded. The only testing denominator is a power of the
positive unit $D_u$; every source denominator in (CA3)--(CA5) is a
positive unit. The possibly unbounded opening ratios remain exclusively
in the four retained positive packets.

\paragraph{The exact low-weight collection.}
The outer/state lemma in Section~\ref{component:p364e14238053} reduces the
independent tests to $Z=1$ and $u,w,v$ infinitesimal, with $w>0$ on
the negative pure-core branch. Use the unchanged circuit monomials
\[
 A_0=wv,\quad B_0=\delta v,\quad C_0=\delta Euw,
 \quad D_0=\delta^2k^3E,
\]
\[
 x=B_0-C_0,\quad y=A_0-kC_0,\quad
 z=uB_0-kC_0,\quad h_0=kC_0-D_0.
\]
Let $q_1,q_2$ be the same two separated radical pairings as in the
contact--conductor proof, and let
\[
 T_4=y^2+z^2+h_0^2-2yz+2q_1^2.
\]
Give $(\delta,k,E,\ell,t,u,w,v)$ weights $(0,1,0,1,0,1,1,2)$.
Then $x$ has weight two and $y,z,h_0,q_1,q_2$ have weight three.
Denote the scalar of (CA5) by $S$, and its two complete first normals
by $L_N,L_M$. Direct multiplication of (CA3)--(CA5) through the
indicated weights gives
\[
 \begin{aligned}
 S_{\le6}&=\frac{4P}{E^4\delta^2}x^2
 -\frac{16(k+u)}{E^2\delta^2}xz
 +\frac8{E\delta^2}T_4,\\
 (L_N)_{\le3}&=0,\\
 (L_M)_{\le3}&=
 \frac{8\{E^2+2E-2\ell(E+1)\}}{E^2\delta}\,x,
 \end{aligned}                                                   \tag{CA6}
\]
where the short coefficient polynomial is
\[
 \begin{split}
 P={}&E^2(E+4)-E(E^2+6E+4)\ell+2(E+1)(E+2)\ell^2\\
 &+E^2(8-3E)k^2+4E^2ku+E^2(E+4)u^2.
 \end{split}
\]
The full hard-normal block, from the quadratic part of (CA5), has
residue
\[
       \operatorname{diag}\left(\frac{16}{E\delta^2},
                                \frac{4(E^2+4)}E\right).          \tag{CA7}
\]
Here and in the residue Grams, $E,\delta$ denote their positive
residues. Formula (CA6), rather than a large expansion, is the finite
coefficient identity used by the proof. Its 37 raw scalar monomials
collect to the displayed three terms. The exact checker reconstructs
the six cancelled native coquotients, the entire expression (CA5),
and both complete normals before making this collection.

\paragraph{Why all higher terms have harmless coefficients.}
The scalar of (CA5), after clearing the stated units, belongs to
$(A_0,B_0,C_0)^2+(D_0^2)$. This follows directly from $q_A=wvD_u$,
$q_B|_{N=M_B=0}\in(v)$, and $Y,W$ in (CA5): the centered scalar
$Y-r uW$ lies in $(v,uw)$; its square is in the same raw monomial
square ideal since $\delta,E$ are units. The sole constant state term
comes from $\alpha$, which has the factor $k^6$, and is a bounded
unit multiple of $D_0^2$. Both full first normals belong to
$(A_0,B_0,C_0)$ by the same displayed expression.

All retained packets vanish to first order at the exact contact
\[ (u,w,v,N,M_B)=(k,\delta k,\delta k^2E,0,0). \]
Thus the scalar belongs to the square of its contact ideal and the
first normals belong to that ideal. The unchanged monic
contact--conductor identities therefore apply, over the coefficient
ring in which only the displayed positive units have been inverted.
They never divide by $k$, $\ell$, $M$, $J$, $t$, or a contact difference;
the displayed $1/a$ is the inverse of a known positive unit.
The contact substitutions preserve the new weights. After (CA6) is
subtracted, scalar weights are at least seven and normal weights at
least four. Their lifted coefficients consequently have positive
weights, hence zero residues. Clearing the finitely many unit
denominators before this polynomial lift and restoring them afterward
changes no residue sign and introduces no mixed inverse.

\paragraph{Strictness, full normals, and restoration.}
The leading scalar is
\[
 \frac{4(E+4)}{E^2\delta^2}x^2+\frac8{E\delta^2}T_4.
\]
The standard two-pairing circuit supplies a strict Gram on
$(y,z,h_0,q_1,q_2)$. The only nonzero first-normal residue is
$8(E+2)x/(E\delta)$, the coefficient of the full normal $M_B$.
Accordingly the coupled $(M_B,x)$ Gram is
\[
 \begin{pmatrix}
 4(E^2+4)/E&4(E+2)/(E\delta)\\
 4(E+2)/(E\delta)&4(E+4)/(E^2\delta^2)
 \end{pmatrix},\qquad
 \det=\frac{256}{E^3\delta^2}>0.                      \tag{CA8}
\]
Together with the independent $N$ pivot in (CA7), this gives a
bounded strict eight-root Gram for the entire remainder (CA5).
Restore (CA1), all four original positive packets, and all separated
normalization factors. The outer/state lemma exhausts the remaining
tests; real and imaginary state evaluations give the complex case.

This proves precisely the stated unit-$e$, unit-$\eta$,
infinitesimal-$r$ slice. Combined with the bounded-unit-$r$ theorem,
it covers every positive $r$ when $\eta,e$ are bounded units and
$k,\zeta$ infinitesimal. It makes no assertion about unbounded
$\eta$ or either nonunit alternative for $e$.


\section{\texorpdfstring{A unit native state scale: the same circuit with a different mass}{A unit native state scale: the same circuit with a different mass}}
\label{component:p890eea5f538b}


\begin{theorem}[Unit $\eta$ and unbounded $e$]
Use the unchanged physical native C1 source, its original event signs,
and the positive shapes and gaps $a,t,M,J$. Suppose $k,\zeta$ are
positive infinitesimals, $\eta$ is a positive bounded unit,
$r=1-a\eta$ is a positive bounded unit, and $e=k\sqrt{\eta/\zeta}$
is unbounded. Then the full original parent is tensor SOS at every
independent projective test and every real or complex state.
No positive lower bound on $a,t,M,J$ is imposed.
\end{theorem}

Write $H=r+\eta$ and $\Delta=1/e$. The seven native coquotients
have residues $2rH$ for $c_2,c_b,c_3,c_5,c_\alpha$, $2$ for
$c_4$, and $4r$ for $c_d$. All are positive units. In particular the
four retained packets of the contact--conductor proof remain positive
and the remainder has bounded coefficients. We reuse their exact
definitions, not asymptotic substitutes for the original source.

\paragraph{Outer tests.}
The exact formulas give
\[
 \frac md=\frac{e^2c_2}{M\eta c_d},\qquad
 \frac\beta d=\frac{e^2c_b}{t\eta c_d}-\frac12.
\]
Both ratios are unbounded. Apply the exact full-parent Alice reflection
(AS3)--(AS4), which exchanges $m\leftrightarrow c$ and
$\beta\leftrightarrow D$ while leaving $v$ unchanged. The already
proved two-strong-current identity then covers every test with $1/v$
bounded, for every Alice angle and every state. Its homogeneous form
includes both literal Alice endpoints, so no inverse-angle point is lost.

For $v$ infinitesimal and Alice infinity, use the critical-current
identity with $g_c=cD/(c+D)$, $\epsilon=d/g_c$ and
$\gamma=Ag_c/d^2$. Its actual sufficient conditions on this chart
are bounded $\epsilon$ and positive-unit $\gamma$: in its exact
four-root identity the correction $L_0$ has factors $c_u$ or $v$,
where $c_u=(1+u^2)^{-1/2}$ is infinitesimal. Thus $\sqrt\epsilon$
need only be bounded. The residue Gram is
$\operatorname{diag}(1,2,2,4\bar\gamma)$.
Use the exact relation $J+a=T$ with $T-1$ infinitesimal, and the
cancelled original coefficient formula
\[
 \frac Da=\frac{2c_5/c_3-J\eta c_d/c_3}{J}.
\]
Its numerator has residue $2/\bar H$. The unchanged native formulas give
\[
 \frac d a\sim\frac{4\eta}{H},\quad
 \frac A a\sim\frac{2\eta^2}{rH},\quad
 \frac D a\sim\frac2{JH},\qquad
 \epsilon\sim\frac{2\eta(JH+a)}H,\quad
 \gamma\sim\frac{H}{4r(JH+a)}.
\]
The identity
$\overline{JH+a}=1+\bar\eta(1-\bar a)^2$ shows that the two required
conditions hold even when $a$ or $J$ is infinitesimal. The calculation
uses the cancelled harmonic coefficient $g_c/a$, so no lower bound
on an opening is inserted. Homogeneous angle normalization includes
literal Alice infinity. It remains to treat bounded $u$ and
infinitesimal $v$.

\paragraph{State entry.}
Normalize the state by a coordinate of maximal absolute value and put
$D_u=1+u^2$ and $w=(W+uY)/D_u$. If $w$ is a unit, the exact five-root
current identity (UW2)--(UW3) applies with $Q=W+uY=D_uw$ a unit.
Its first four diagonal residues are $8rH$ and its fifth is
$8D_u$; both negative corrections have factors $\Delta$ or
$\Delta^2\eta$. Thus this chart is strict.

If $w$ is infinitesimal, the nonnegative pure-testing core is a direct
exit. On the negative-core branch $Z$ is a unit. Indeed its residue
at $v=0$ is $2\{(X+uZ)^2+(uW-Y)^2\}$; if $Z,w$ and both currents
had zero residue, every normalized state coordinate would have zero
residue. Normalize $Z=1$. The exact core Schur identity gives $w>0$
after the usual positive-angle sign gauges. All transformations here
belong solely to the testing field.

\paragraph{The ordinary-$u$ chart needs no contact division.}
Use the exact source change
\[
 \delta=\Delta R/K,\quad b^2=\eta K/R,\quad K=1+k^2,
 \quad Q=K+ak^2b^2,\quad T_n=K^2-t\delta^2k^4Qb^2.
\]
All source denominators are positive units, $\delta,k$ are
infinitesimal and $b^2\sim\eta$. Let
\[
 A_0=wv,\quad B_0=\delta v,\quad C_0=\delta b^2uw,
 \quad D_0=\delta^2k^3b^2,\qquad x=B_0-C_0.
\]
The five-term exact scalar belongs to
$(A_0,B_0,C_0)^2+(D_0^2)$ and both complete first normals belong
to $(A_0,B_0,C_0)$. Collecting these displayed factors before taking
residues is legitimate over the bounded coefficient ring. At
$\delta=k=w=v=0$, keeping bounded $u$ free, the scalar Gram is
\[
 8rH D_u^2A_0^2+
 8D_u\{x^2+u^2rH B_0^2\}+8rH D_uD_0^2.
 \tag{UE1}
\]
The hard block on $(N,\delta M_B)$ is
$\operatorname{diag}(16rH,16rH+8\eta^2/D_u)$, and the complete
first normals have residues $0$ and $16\eta x$. The latter is the
coefficient of $\delta M_B$ in the full polynomial, so the symmetric
Gram off-diagonal entry between $\delta M_B$ and $x$ is $8\eta$.
The coupled $(\delta M_B,x)$ block has determinant
\[
 (16rH+8\eta^2/D_u)\,8D_u-(8\eta)^2=128D_urH>0.
 \tag{UE2}
\]
If $u$ is a unit, the remaining $B_0$ pivot in (UE1) is strict,
as are the $A_0,D_0,N$ pivots. Every omitted coefficient has zero
residue. This proves the chart by a bounded strict Gram on these
six roots, without any mixed inverse. Equation (UE1) follows directly
from $a\eta+r=1$: its $(B_0,C_0)$ part before collection is
$8D_u\{(1+u^2rH)B_0^2-2B_0C_0+C_0^2\}$.

\paragraph{The infinitesimal-$u$ chart uses the old conductor.}
Now $\delta,k,u,w,v$ are infinitesimal, while $b,a,t$ are bounded.
Use the unchanged roots
\[
 y=A_0-kC_0,\quad z=uB_0-kC_0,\quad h=kC_0-D_0,
 \quad q_1=a_1a_2-a_3a_4,
\]
and the second pairing $q_2=a_1a_3-a_2a_4$ of the four separated
radicals. Set
\[
 T_4=\sum_{i=1}^4a_i^4-4\prod_{i=1}^4a_i
     =y^2+z^2+h^2-2yz+2q_1^2,
 \qquad \rho=(1-ab^2)(1+(1-a)b^2)-1.
\]
Give $(\delta,k,b,a,t,u,w,v)$ weights $(1,1,0,0,0,1,2,3)$.
The root $x$ has weight four and the other soft roots have weight
five. Substitution at the exact contact preserves this grading.
Writing $E_0$ for the complete scalar correction relative to the
old leading polynomial, the five-term factorization and the two
complete normal formulas give
\[
 \begin{aligned}
 E_0&=8\rho T_4-8(\rho+b^2)k^2x^2+O_{\rm wt}(12),\\
 L_1&=O_{\rm wt}(6),\qquad
 L_2=16b^2x+O_{\rm wt}(6).
 \end{aligned}\tag{UE3}
\]
These identities require only multiplication of the short source
polynomials. In particular the residue correction is the single
factor $\rho$ multiplying the same circuit, rather than a new
polynomial positivity problem.

The unlocalized monic contact--conductor lift applies unchanged.
After the terms explicitly displayed in (UE3) are removed, scalar
coefficient weights are at least two and normal coefficient weights
are at least one. Therefore they all have zero residue. There is no
division by $k,b,\delta,a,t$ or a contact difference. The leading
strict soft Gram is now
\[
                  8x^2+8rH T_4.
\]
The two-pairing circuit identity supplies a strict Gram on
$(y,z,h,q_1,q_2)$, and $rH$ is a positive unit. The complete hard
block is $\operatorname{diag}(16rH,16rH+8\eta^2)$; its only
nonzero first-normal residue is $16\eta x$. Its coupling to the
$8x^2$ pivot has determinant $128rH>0$. Retain the old positive
$g_4^2$ square. The resulting eight-root remainder Gram is strict.

Finally restore all four positive packets, the positive square just
retained, and the positive pure-coefficient multiplier. The state and
angle cases above exhaust the independent tests. Complex states split
into their real and imaginary quadratic evaluations. The result
closes the entire stated source domain, rather than only an optimizing
or matching testing chart.



\section{\texorpdfstring{Ordinary actual Alice angle, with arbitrary state ratio}{Ordinary actual Alice angle, with arbitrary state ratio}}
\label{component:p9f00f949fc8c}

\label{p9f00f949fc8c:sec:c1-arbitrary-state-ratio}

The positive bounded-unit assumption on $s=g/f$ can be removed from
the arbitrary-opening C1 theorem. The new ingredient is a different
grading of the already proved polynomial contact--conductor identities.
No new ideal membership or division is required.

\begin{theorem}
Use the unchanged native C1 source, normalized by $c=2$ and
$\lambda=n=0$. Suppose $k$ is a positive bounded unit, $f,g>0$ are
infinitesimal, and
\[
 \begin{gathered}
 \eta=f^2,\quad\zeta=g^2,\quad K=1+k^2,\qquad
 R=K-ak^2\eta,\quad T=K-k^2t\zeta,\\
 a,t>0,\qquad M=R-t>0,\qquad J=T-a>0.
 \end{gathered}
 \tag{AS1}
\]
Then the full original parent is tensor SOS at every independent
projective test and every real or complex state. The ratio $s=g/f$
is arbitrary positive. No positive lower bound on $a,t,M,J$ and no
boundedness of original event coefficients is assumed.
\end{theorem}

The native source and every retained packet are exactly those of
(OA2), (AO3)--(AO5), with $s=g/f$ and $\Delta=s/k$. Since
$a<T<K$ and $t<R<K$, the shapes are bounded, while $R,T$ are units
of residue $K$. The six coquotients $c_2,c_b,c_3,c_5,c_4,c_\alpha$
have residue $2K$, and $c_d$ has residue $4K^2$.
All original events are positive. For the two subtracted expressions,
this follows from
\[
 \frac Da=\frac1J\left(\frac{2c_5}{c_3}
                   -J\eta\frac{c_d}{c_3}\right),\qquad
 \frac\beta a=\frac1{s^2t}
      \left(\frac{2k^2c_b}{c_3}-t\zeta\frac{c_d}{c_3}\right).
 \tag{AS2}
\]
The parentheses have positive unit residues $2$ and $2k^2$.

\paragraph{Exact reflection reduces to bounded $s$.}
The whole-parent reflection is
\[
 u'=1/u,\quad v'=v,\quad
 (W',X',Y',Z')=(vW-X,-W-vX,Z-vY,Y+vZ),
 \tag{AS3}
\]
with $(\lambda\ n)(m\ c)(A\ \alpha)(D\ \beta)$ on events and
\[
 u^2F(\mathrm{ref}(c);1/u,v;\psi')=(1+v^2)F(c;u,v;\psi).
 \tag{AS4}
\]
This is an identity of the full parent, including the signed core.
At the native contact set
$H_f=1+(1-a)\eta$, $H_g=1+(1-t)\zeta$, and
$V=1+k^2\eta\zeta$. Both $H_f,H_g$ are positive units. After
normalizing the reflected contact by its positive $Z'$ coordinate,
the exact reflected native parameters are
\[
\begin{gathered}
 k'=1/k,\qquad s'=\frac{H_g}{sH_f},\qquad
 \eta'=\frac{k^2\zeta H_f}{H_g},\qquad
 \zeta'=\frac{k^2\eta H_g}{H_f},\\
 a'=\frac{M}{k^2H_f},\quad t'=\frac{J}{k^2H_g},\qquad
 M'=\frac{aV}{k^2H_g},\quad J'=\frac{tV}{k^2H_f}.
\end{gathered}\tag{AS5}
\]
For example, the reflected state has
$r'=Q_c/H_g$, $\theta'=\eta J/H_f$, where
$Q_c=1-k^2\zeta(1-a\eta)$; the differences
$H_g-Q_c=\zeta M$ and $H_f-E=\eta J$ prove (AS5).
The reflected coefficient vector, after positive normalization to
$c'=2$, is the reflected native source: this follows from the same
six contact equations and their nonzero selected minor used for the
original native source. All displayed shapes and gaps are positive.
Thus unbounded $s$ becomes infinitesimal $s'$, with both
$\eta',\zeta'$ infinitesimal and $k'$ a unit. Unit $s$ is the
previous arbitrary-opening theorem. It remains to prove the case
$s$ infinitesimal.

\paragraph{Outer angles require only a positive current mass.}
Put $D_1=2c_5/c_3-J\eta c_d/c_3$, a unit of residue $2$.
For $g_c=cD/(c+D)$ the exact cancellation
\[
 \frac{g_c}{a}=\frac{2D_1}{2J+aD_1},\qquad
 2J+aD_1=2T+a(D_1-2)
 \tag{AS6}
\]
gives
\[
 d/c,\ d/D\text{ infinitesimal},\qquad
 \gamma=\frac{Ag_c}{d^2}\text{ a positive bounded unit},
 \qquad\bar\gamma=\frac1{4\bar K^3}.
 \tag{AS7}
\]
The two-strong-current identity therefore covers every test with
$1/v$ bounded, including Bob infinity, for every Alice angle and state.
For $v$ infinitesimal and Alice unbounded, use the exact critical-current
identity with residue diagonal
\[
 \operatorname{diag}(1,2,2,4\bar\gamma-\bar c_u^2),
 \qquad c_u=(1+u^2)^{-1/2}.
 \tag{AS8}
\]
Here $\bar c_u=0$ and $\bar\gamma>0$, so it is strict.
The former sufficient assumption $\bar\gamma\ge1/4$ is unnecessary
on this particular reciprocal-Alice chart. Its same exact identity
and homogeneous normalization also include literal Alice infinity.
Consequently only bounded $u$ and infinitesimal $v$ remain.

\paragraph{The complementary state chart.}
Normalize the testing state by a coordinate of maximal absolute value.
Put $D_u=1+u^2$ and $w=(W+uY)/D_u$, a pure testing coordinate.
If $w$ is a unit, the retained-current identity (UW2)--(UW3) applies
with $Q=W+uY=D_uw$ a unit. It requires only bounded $u,v,Z$, not
infinitesimal $u$ or unit $k$. The five roots are
\[
 r_2,\ B,\ \Delta r,\ \Delta r_5,\ \Delta\eta Y.
\]
Their positive diagonal coefficients are
$H_2,H_bD_v,H_3D_v,H_5,4RTc_4D_uD_v$, all units.
The only two negative corrections have factors
$\Delta Z/Q$ and $\Delta^2\eta Z^2/Q^2$, hence zero residue.
Restoring the retained packets proves this chart.

Otherwise $w$ is infinitesimal. Literal axes and a nonnegative pure
testing core are direct original-event SOS exits. On the negative-core
branch, sign gauges give $u,v>0$. The coordinate $Z$ must be a unit:
at $v$ of zero residue the
core has residue $2((X+uZ)^2+(uW-Y)^2)$. If $Z,w$ and these two
currents all had zero residue, the normalized state would have every
coordinate of zero residue. Thus a negative-core test with small $w$
has unit $Z$; normalize $Z=1$ by a pure testing factor. The
pure-testing Schur identity now gives $w>0$; before this normalization
its conclusion is $wZ>0$.

\paragraph{The same conductor, with $k$ of weight zero.}
In this last chart use exactly (SD3)--(SD8), retaining all four
positive packets. Set
\[
 \delta=\Delta R/K,\qquad b^2=\eta K/R,
 \qquad Q=K+ak^2b^2,\quad T_n=K^2-t\delta^2k^4Qb^2,
 \quad\mathcal C=K^6Q^4T_n.
\]
Here $\delta,b$ are positive infinitesimals, and $k,a,t,u$ are bounded.
The exact source substitution is unchanged:
\[
 \eta=Kb^2/Q,\quad\Delta=\delta Q/K,\quad
 \zeta=\delta^2k^2Qb^2/K,\quad R=K^2/Q,\quad T=T_n/K.
\]
Let $E_0,L_1,L_2$ be the scalar correction and two full first normals
in the contact--conductor proof. Its polynomial construction places
\[
 E_0\in(x,y,z,h,q_1)^2,\qquad L_1,L_2\in(x,y,z),
 \tag{AS9}
\]
where
\[
 \begin{gathered}
 A_0=wv,\ B_0=\delta v,\ C_0=\delta b^2uw,\ D_0=\delta^2k^3b^2,\\
 x=B_0-C_0,\quad y=A_0-kC_0,\quad z=uB_0-kC_0,\quad h=kC_0-D_0.
 \end{gathered}
\]
These are universal polynomial identities; their proof never assumed
that $k$ was small or inverted it.

Give $(\delta,k,b,a,t,u,w,v)$ weights $(1,0,1,0,0,0,1,3)$.
All four monomials $A_0,B_0,C_0,D_0$, and all the soft roots, have
weight $4$. The substitutions at contact
$(u,w,v)=(k,\delta k,\delta k^2b^2)$ preserve this grading.
The factored five-term scalar gives
\[
 \operatorname{wt}E_0\ge10,\qquad
 \operatorname{wt}L_1\ge8,\qquad
 \operatorname{wt}L_2\ge6.\tag{AS10}
\]
Here is a short way to read these bounds. One has
$Q-K=O(2)$, $T_n-K^2=O(4)$, $V_n-K^2=O(4)$,
$D_v-1=O(6)$, $Y_n=O(1)$ and $B_n+KT_n=O(4)$.
The six cleared coquotients equal $2K^3Q^3+O(2)$, while
$C_d=4K^4Q^3+O(2)$. In the five-term scalar formula, discard terms
of weight at least $10$; its four weight-eight squares and cross term
are precisely $8K^2D_u\mathcal C P$. For the normals use
\[
 K(C_b-C_2)=\delta^2k^2Qb^2C_3,
 \qquad Q(C_3-C_5)=-k^2Kb^2C_b,
\]
and $S_B,S_C=O(2)$ in their displayed factorizations. This gives
(AS10) without an expanded source calculation.

Monic contact division preserves these weights. The conductor lifts
therefore have coefficient weights at least $2,4,2$, respectively.
Every such polynomial coefficient contains a positive power of
$\delta,b,w$, or $v$, and has zero residue. There are no negative
exponents, including on $a,t$; the only source denominators used before
this polynomial construction are the displayed positive bounded units.

The leading scalar remains exactly
\[
 P=D_u(A_0^2+B_0^2)+K(C_0^2+D_0^2)-2K^2B_0C_0
   =x^2+g_4^2+\sum_{i=1}^4a_i^4-4\prod_{i=1}^4a_i.
\]
Here $g_4=uwv-\delta^2k^4b^2$ and
\[
 \begin{gathered}
 a_1=\sqrt{wv},\quad a_2=\sqrt{\delta uv},\quad
 a_3=b\sqrt{\delta kuw},\quad a_4=\delta k^{3/2}b,
 \\
 q_1=a_1a_2-a_3a_4,\qquad q_2=a_1a_3-a_2a_4.
 \end{gathered}
\]
Their radicals separate between coefficient and testing fields;
$y=a_1^2-a_3^2$, $z=a_2^2-a_3^2$, $h=a_3^2-a_4^2$.
The universal two-pairing circuit identity gives a strict constant Gram
on $(x,y,z,h,q_1,q_2)$. Retain $g_4^2$ as a positive square; its
coefficients need not have zero residue when $k,u$ are units.
The complete hard block on $(N,\delta M_B)$ has positive residue
$16\bar K^{14}I_2$ after multiplication by $\mathcal C$.
Indeed $\Delta/\delta=Q/K\sim1$, $H_i\sim8K^2$, and the remaining
hard-block terms have factors $\eta$ or $v$. The full first normals
are exactly those in (AS9), not only their restrictions at contact.
Thus (AS9)--(AS10) give a bounded strict eight-root remainder Gram.
Restore $g_4^2$ and the four positive packets, then undo the positive
pure source multiplier $\mathcal C\Delta^2\Lambda$.

\paragraph{Exhaustion and precise advance.}
The source reflection, unit-$s$ theorem, two outer-angle identities,
unit-$w$ current identity, and small-$w$ conductor cover every source
and every independent projective test under (AS1). Pure testing
normalizations and pure source multipliers are restored separately;
no mixed inverse or radical occurs. Complex states split into real
and imaginary parts.

Together with the complete infinitesimal-$k$ theorem and the exact
orientation $0<k\le1$, this proves the all-tests assertion whenever
the actual native entry has both $\eta$ and $\zeta$ infinitesimal,
for every positive $k$ in that orientation and every state ratio.
The remaining native complement is therefore: $k$ infinitesimal with
$\zeta$ infinitesimal and $\eta$ a unit or unbounded; $k$
infinitesimal with $\zeta$ non-infinitesimal; and $k$ a unit with
at least one of $\eta,\zeta$ non-infinitesimal. Each is still
intersected with physicality and the complements of all other inherited
admissions and symmetry images. No whole-global completion is asserted.


\section{\texorpdfstring{Ordinary-angle C1 with arbitrary positive shape openings}{Ordinary-angle C1 with arbitrary positive shape openings}}
\label{component:p64b054e815d0}

\label{p64b054e815d0:sec:c1-arbitrary-openings}

The unit assumptions on $a,t,M,J$ in
Theorem~\ref{p18a57c9c271d:thm:c1-ordinary-angle} can all be removed. The large
positive packets must be retained before taking a bounded Gram residue.

\begin{theorem}[Ordinary angles and arbitrary positive openings]
Use the unchanged native C1 source (OA2), with $c=2$ and
$\lambda=n=0$. Assume only that $k,s=g/f$ are positive bounded units,
$\eta=f^2>0$ is infinitesimal, and
\[
\begin{gathered}
 \zeta=g^2=s^2\eta,\quad K=1+k^2,\quad
 R=K-ak^2\eta,\quad T=K-k^2s^2\eta t,\\
 a,t>0,\quad M=R-t>0,\quad J=T-a>0.
\end{gathered}\tag{AO1}
\]
Then the full original parent is tensor SOS at every independent
projective test and every real or complex state. No lower bound on
$a,t,M,J$, and no boundedness assumption on the original event
coefficients, is required.
\end{theorem}

Since $a<T<K$ and $t<R<K$, both shapes are bounded, and $R,T$ have
residue $K$. In particular
\[
       p=t/R\in(0,1),\qquad h=a/T\in(0,1)              \tag{AO2}
\]
are bounded at every opening rate. The seven native coquotients have
the same residues as in (OA3). All original events are positive:
$A/a$ and $\alpha/a$ are $\eta^2$ times positive units, while
$d/a$ is $\eta$ times a positive unit; positivity of the remaining
events also follows from the explicit bounds below.

\paragraph{Retain the four large packets exactly.}
Put $\Delta=s/k$, $\Lambda=RTc_3/(2a)$, and
\[
 H_2=4Tc_2,\quad H_b=4Tc_b,\quad H_3=4Rc_3,\quad H_5=4Rc_5.
\]
All four have residue $8K^2$. For a bounded finite testing angle define
\[
 D_u=1+u^2,\quad D_v=1+v^2,\quad
 x=X+uZ,\quad y=uW-Y,\quad
 q_A=v(W+uY),\quad q_B=v(Z-uX),
\]
\[
                    N=x-pq_A,\qquad L=y-hq_B.
\]
The four nonnegative packets are
\[
 \Pi_1=H_2\frac tM N^2,\quad
 \Pi_2=H_b\frac MtD_vN^2,\quad
 \Pi_3=\Delta^2H_3\frac JaD_vL^2,\quad
 \Pi_4=\Delta^2H_5\frac aJL^2.                         \tag{AO3}
\]
Their possibly unbounded pure source coefficients stay in their
separate square roots. They never enter the bounded remainder Gram.
The exact whole-parent identity is
\begin{align*}
 G:=\Delta^2\Lambda F-\sum_{i=1}^4\Pi_i
 ={}&H_2\{N^2-2q_AN+(1-p)q_A^2\}
     +H_bD_v\{N^2+2q_AN+pq_A^2\}\\
 &+\Delta^2H_3D_v\{L^2+2q_BL+hq_B^2\}
     +\Delta^2H_5\{L^2-2q_BL+(1-h)q_B^2\}
     +\mathcal W,                                      \tag{AO4}\\
 \mathcal W={}&4RT\Delta^2\eta^2c_4D_uD_vY^2
   +4RT\Delta^4k^6\eta^2c_\alpha D_u(vY-Z)^2\\
 &+2RT\Delta^2\eta c_dD_u\{(vY-Z)^2-D_vZ^2\}.
\end{align*}
This is the original four-packet identity (SD3)--(SD4), valid without
unit shapes. For example, the first subtraction uses only
$R=M+t$ and $x=N+(t/R)q_A$; its result is
$H_2\{N^2-2q_AN+(M/R)q_A^2\}$. The other three follow in the
same way from $R=M+t$ and $T=J+a$.

\paragraph{A bounded coordinate change with no mixed inverse.}
Let $w=(W+uY)/D_u$, which lies entirely in the testing field. Then
\[
\begin{gathered}
 q_A=D_uv w,\quad
 X=N+pD_uv w-uZ,\quad
 q_B=v(D_uZ-uN-puD_uv w),\\
 W=w+\frac{u(L+hq_B)}{D_u},\qquad
 Y=uw-\frac{L+hq_B}{D_u}.                             
\end{gathered}
\tag{AO5}
\]
These are exact inverse formulas. Their only testing denominator is
the positive pure testing unit $D_u$; their source coefficients use
only $R^{-1},T^{-1}$. In particular no determinant involving $p,h,u,v$
is inverted. Set $\widehat G=D_uG$. It is a polynomial in
$N,L,u,w,v,Z$ with bounded coefficients, since this multiplication
clears the denominators in (AO5).

Suppose $v$ is infinitesimal and $u$ is a finite positive unit.
The exact polynomial $\widehat G$ has a bounded Gram on
\[
                    N,\ L,\ vw,\ vZ,\ \eta w,\ \eta Z.
                                                               \tag{AO6}
\]
After expanding only the squares in (AO4), every term contains a
product of two listed roots. The coefficient errors carry $v$, $\eta$,
or a native coquotient difference of zero residue. To display the
Gram residue, use the purely formal substitutions
$N=\eta\xi$, $L=\eta\chi$, $v=\eta\nu$ and read the coefficient of
$\eta^2$. It is
\[
 \mathcal G=16K^2D_u\xi^2+16\Delta^2K^2D_u\chi^2
                         +8\Delta^2K^3D_u^2S,        \tag{AO7}
\]
where
\[
 S=u^2w^2+s^2k^4Z^2+
   \frac{D_u}{K}\nu^2\left(\frac{k^2}{s^2}w^2+Z^2\right)
   -2Ku\nu wZ.                                       \tag{AO8}
\]
There is no shape parameter in this residue other than the fixed
units $k,s$. In particular no assertion that $v/\eta$ is bounded
has been made. Formula (AO7) follows directly from
$H_2,H_b,H_3,H_5\sim8K^2$, $c_4,c_\alpha\sim2K$, $c_d\sim4K^2$:
the $p$ and $h$ terms in the two scalar pairs sum to one, while the
two linear normal coefficients have zero residue.

\paragraph{Finite unit Alice angles: the same scalar, fewer hard roots.}
The capacity identity (OA10), with $W$ replaced by $w$, proves a
strict six-root Gram unless $\bar u=\bar k$. At the matched angle,
\[
 S=(\nu Z-kw)^2+
       \left(\frac{k}{s}\nu w-sk^2Z\right)^2.
\]
Define
\[
 r_1=vZ-\eta kw,\qquad r_2=(k/s)vw-\eta sk^2Z.
\]
The exact relation
\[
 wr_1-(s/k)Zr_2=\eta k(s^2Z^2-w^2)                    \tag{AO9}
\]
closes every state with $\overline{s^2Z^2-w^2}\ne0$ by adding its
identically zero squared relation to the Gram on $(N,L,r_1,r_2,\eta)$.
The multiplier in this relation is never inverted.

Normalize the testing state. If $N$ or $L$ has nonzero residue,
$\widehat G$ has positive residue and is already certified. Otherwise
the inverse formulas force $w$ or $Z$ to be a unit. When (AO9) has
zero multiplier residue, $Z$ is a unit; normalize $Z=1$.
The negative-core branch has $w>0$ by the pure testing core Schur
calculation printed in the ordinary-angle proof. Thus only
$\bar w=\bar s$ remains. The nonnegative-core branch and literal
axes are direct original-event SOS exits.

\paragraph{The exact contact with arbitrary shape rates.}
The native contact (OA4), under (AO5), is
\[
 (u_c,w_c,v_c,N_c,L_c)
       =\left(k,\frac{sR}{K},ks\eta,0,0\right).        \tag{AO10}
\]
All four packets (AO3) and all their first derivatives vanish there.
The original native contact identities consequently give exact value
and first-derivative zero for $\widehat G$ as well. This remains true
when a packet coefficient is unbounded: each coefficient is an element
of the source field multiplying a square that vanishes exactly.

Use the five roots
\[
       N,\ L,\ v-v_c,\ \eta(w-w_c),\ \eta(u-k).       \tag{AO11}
\]
The complete polynomial $\widehat G$ belongs to
$(N,L,v)^2+\eta(N,L,v)+(\eta^2)$, by (AO4)--(AO5).
Translation by $v_c=\eta ks$ preserves this inclusion. Its constant
normal coefficient is $\eta^2$ times a bounded polynomial in the
two soft deviations; its linear normal coefficients are $\eta$ times
such polynomials. Exact contact and stationarity make their constant
and the required first coefficients vanish. Finite monic Taylor
division therefore gives an exact bounded Gram on (AO11), just as in
(OA14), with no division by any testing/source comparison.
Higher Taylor terms have zero residue after this allocation.

The residue Gram is now particularly simple. Its first two blocks are
$16K^3$ and $16\Delta^2K^3$. Its remaining block is
$8\Delta^2K^5$ times the half-Hessian of $S$ at
$(\nu,w,u)=(ks,s,k)$ and $Z=1$. The latter is exactly (OA15), with
positive leading principal minors
\[
                         K,\qquad4k^4,\qquad12s^2k^6/K.
\]
Thus the entire five-root Gram is strict independently of all four
shape-opening rates.

\paragraph{The outer charts without bounded event weights.}
These charts admit a simpler uniform retained-packet comparison.
The exact source formulas give
\begin{align*}
 m/a&\ge\frac{2k^2c_2}{s^2Rc_3},&
 c/a&\ge\frac2T,\\
 D/a&\ge\frac{2c_5}{Tc_3}-\frac{\eta c_d}{c_3},&
 \beta/a&\ge\frac{2k^2c_b}{s^2Rc_3}-\frac{\eta c_d}{c_3}.
                                                               \tag{AO12}
\end{align*}
All four right sides are positive bounded units. Choose a positive
real $\mu$ smaller than their residues. In $F/a$, retain the four
nonnegative event packets with weights
$m/a-\mu,c/a-\mu,D/a-\mu,\beta/a-\mu$. The remaining complete
parent has the four corresponding weights all equal to $\mu$,
and the unchanged weights
\[
 d/a=\eta\,d_1,\quad A/a=\eta^2A_2,\quad
 \alpha/a=\eta^2\alpha_2,
\]
where $d_1,A_2,\alpha_2$ are positive units. This is an exact
decomposition at the present test, not a claim that the capped source
is globally physical.

If the homogeneous Bob sine has nonzero residue, the four event
amplitudes have trivial common state kernel, as proved in the
ordinary-angle theorem. Their bounded positive Gram absorbs the
infinitesimal signed core. If Bob is small and $u$ is infinitesimal,
use the six roots
$(x,y,vW,vZ,\eta uW,\eta Z)$ of the same theorem. The four hard
weights $\mu$ give a strict four-root block and the two weak diagonal
coefficients are positive units. The signed soft cross term has an
additional factor $u$, so its Gram coefficient has zero residue.
At the reciprocal Alice endpoint put $q=1/u$ and use
\[
                x'=qX+Z,\ y'=W-qY,\ vY,\ vX,
                       \eta Y,\ \eta qX.
\]
The hard residue is
$4\mu\{(vY-x')^2+x'^2+y'^2+(y'+vX)^2\}$;
the weak diagonals have positive unit coefficients. Assign the signed
soft cross term as a multiple of $q(vX)(\eta Y)$, which has zero
residue. These are the exact finite and reciprocal monomial
allocations already proved in the ordinary-angle theorem; they use
no shape inverse. Literal infinity follows in the homogeneous chart.

These outer charts and the finite-unit Alice analysis exhaust all
independent tests. Restore the retained packets, then the positive
pure factors $D_u$, $\Delta^2\Lambda$, or $a$ used in the respective
charts. Every radical and inverse is confined to one field. Complex
states split into their real and imaginary parts.



\section{\texorpdfstring{An ordinary-angle C1 family: all independent tests}{An ordinary-angle C1 family: all independent tests}}
\label{component:p18a57c9c271d}

\label{p18a57c9c271d:sec:c1-ordinary-angle}

This section closes an actual source family outside the small-$k$
uniform theorem. All inverses below are either real constants or
bounded units in one of the two independent real closed fields.
In particular no quotient of a testing infinitesimal by a source
infinitesimal is used.

\begin{theorem}[Regular ordinary-angle C1 source]
\label{p18a57c9c271d:thm:c1-ordinary-angle}
Use the seven native C1 coquotients and original coefficients of
the source-definition appendix, with $c=2$ and $\lambda=n=0$.
Suppose
\[
 k,s,a,t,M,J\text{ are positive bounded units},\qquad
 \eta>0\text{ is infinitesimal},\qquad
 f^2=\eta,\quad g^2=s^2\eta,
\]
\[
 K=1+k^2,\quad R=K-ak^2\eta,\quad T=K-k^2s^2\eta t,
 \quad M=R-t>0,\quad J=T-a>0.                 \tag{OA1}
\]
Then the full original parent is tensor SOS at every independent
projective test, for every real or complex state.
\end{theorem}

Here $s=g/f$ is the actual contact state ratio. Its being a unit is
part of the hypothesis. The source normalization and all six original
retained event coefficients are preserved. In particular this theorem
includes $k=s=1$, $a=t=1/2$, and arbitrary positive infinitesimal $\eta$.

\paragraph{Exact source and its leading scales.}
Put $e=k/s$ and $\Delta=s/k$ in the native formulas. To fix completely
which source is meant, they are
\[
\begin{gathered}
 d=\frac{2a\eta c_d}{c_3},\quad
 m=\frac{2ak^2c_2}{s^2Mc_3},\quad
 A=\frac{2a\eta^2c_4}{c_3},\quad
 \alpha=\frac{2as^2k^4\eta^2c_\alpha}{c_3},\\
 D+\frac d2=\frac{2ac_5}{Jc_3},\qquad
 \beta+\frac d2=\frac{2ak^2c_b}{s^2tc_3},\qquad c=2.
\end{gathered}
\tag{OA2}
\]
The letter $D$ here denotes the fifth event, whereas $D_u=1+u^2$.
Setting $\eta=\zeta=0$ in the short native polynomials gives
$c_2=c_b=c_3=c_5=c_4=c_\alpha=2K$ and $c_d=4K^2$.
Consequently the following positive bounded source expressions have
the same residues as the indicated original coefficients:
\[
\begin{gathered}
 m_0=\frac{2ak^2}{s^2(K-t)},\quad
 c_0=2,\quad D_0=\frac{2a}{K-a},\quad
 \beta_0=\frac{2ak^2}{s^2t},\\
 d/\eta\sim4aK,\qquad A/\eta^2\sim2a,\qquad
 \alpha/\eta^2\sim2as^2k^4.                         
\end{gathered}
\tag{OA3}
\]
The symbol $\sim$ here means equality of residues, not equality of
the source. Each displayed inverse is a pure source unit inverse.
All original events are positive under (OA1), as follows from their
positive leading coefficients.

The exact contact is
\[
 u_c=k,\quad v_c=ks\eta,\quad
 (W_c,X_c,Y_c,Z_c)=(s,-k(1-ts^2\eta),ks(1-a\eta),1).
                                                        \tag{OA4}
\]
Its value and all first derivatives vanish. This assertion uses the
native identities (SD9)--(SD11), which are polynomial identities valid
for arbitrary $k$, not an asymptotic extension of a small-$k$ assertion.
The accompanying small checker repeats all six native balances.
For actual canonical C1 contacts, the source transport and nonzero
selected contact minor in (CE1)--(CE2) identify (OA2) with the original
source; no new source-entry assumption is implicit here.

\paragraph{A four-square reduction.}
Write
\[
 x=X+uZ,\quad y=uW-Y,\quad D_u=1+u^2,\quad D_v=1+v^2.
\]
The exact complete parent, including every signed term, is
\begin{align*}
 F={}&4m\{v(D_uW-uy)-x\}^2+4cD_vy^2
       +4AD_uD_v(uW-y)^2\\
 &+(4D+2d)\{y+v(ux-D_uZ)\}^2
       +4\alpha D_u\{v(uW-y)-Z\}^2\\
 &+(4\beta+2d)D_vx^2
       +2dD_u\bigl(\{v(uW-y)-Z\}^2-D_vZ^2\bigr).
                                                        \tag{OA5}
\end{align*}
The last line includes the essential $2dD_u(vY-Z)^2$ term.
No event is capped or discarded.

For a moment use formal variables $\xi,\chi,\nu$ and substitute
$x=\eta\xi$, $y=\eta\chi$, $v=\eta\nu$ into (OA5).
The coefficient of $\eta^2$ is
\begin{align*}
 \mathcal F={}&4m_0(\nu D_uW-\xi)^2+4\beta_0\xi^2
       +4c_0\chi^2+4D_0(\chi-\nu D_uZ)^2\\
 &+8aD_uu^2W^2+8as^2k^4D_uZ^2
       -16aKD_uu\nu WZ.
                                                        \tag{OA6}
\end{align*}
This is only a polynomial way to read a bounded Gram residue; it does
not substitute or compare any actual testing and coefficient rates.
Completing the first four squares gives
\[
 \mathcal F=4(m_0+\beta_0)
  \left(\xi-\frac tK\nu D_uW\right)^2
 +4(c_0+D_0)\left(\chi-\frac aK\nu D_uZ\right)^2
 +8aD_uS,                                               \tag{OA7}
\]
where
\[
 S=u^2W^2+s^2k^4Z^2+
   \frac{D_u}{K}\nu^2\left(\frac{k^2}{s^2}W^2+Z^2\right)
   -2Ku\nu WZ.                                         \tag{OA8}
\]

For bounded $u,v$ with $v$ infinitesimal, (OA5) is an exact bounded
quadratic Gram on
\[
                 x,\ y,\ vW,\ vZ,\ \eta W,\ \eta Z.    \tag{OA9}
\]
To see the assertion directly, expand only the displayed squares in
(OA5). Each monomial contains a product of two listed roots; all
remaining factors are bounded. Terms omitted in (OA6), after this
allocation, carry $v$, $\eta$, or a source coefficient difference of
zero residue. Hence its Gram residue is exactly (OA6)--(OA8), with
$\eta\nu$ read as $v$. This observation also applies to the modified
endpoint root lists below.

\paragraph{Angles away from the actual contact.}
Sign gauges allow $u,v\geq0$. Suppose first that $u$ has positive
finite residue. In the four roots
$\nu W,\nu Z,W,Z$, split the cross term in (OA8) between the pairs
$(\nu W,Z)$ and $(\nu Z,W)$, using their identical products.
The two positive cross-term capacities are
\[
 \sqrt{D_u/K}\,k^3,\qquad \sqrt{D_u/K}\,u.
\]
Their sum is strictly larger than $Ku$ unless $\bar u=\bar k$, since
the exact factorization is
\[
 (1+u^2)(u+k^3)^2-K^3u^2
  =(u-k)^2(k^4+2k^3u+2ku+u^2).                         \tag{OA10}
\]
Thus, when $\bar u\ne\bar k$, choose a real splitting constant
strictly within the two residue capacities. The resulting two
two-by-two blocks, together with the two completed hard squares,
give a strict residue Gram on (OA9).

If $u$ is infinitesimal, replace $\eta W$ in (OA9) by $\eta uW$.
Assign the entire cross term to $(vW,\eta Z)$; its coefficient has
zero residue. The four soft diagonal coefficients are positive units.
The same monomial allocation in (OA5) proves an exact strict Gram.

The reciprocal Alice endpoint requires no source reflection. Put
$h=1/u$, multiply the homogeneous parent by the pure testing square
$h^2$, and use $x'=hX+Z$, $y'=W-hY$. Its six roots are
\[
                 x',\ y',\ vY,\ vX,\ \eta Y,\ \eta hX.
                                                        \tag{OA11}
\]
When $h$ is infinitesimal, the four strong residue terms are
$4m_0(vY-x')^2+4\beta_0x'^2+4c_0y'^2+4D_0(y'+vX)^2$.
The two weak diagonal reserves are $8a(\eta Y)^2$ and
$8as^2k^4(\eta hX)^2$. The signed cross term is assigned as a
multiple of $h(vX)(\eta Y)$, so has zero residue; all other corrections
also have zero residue. This is a strict six-root Gram, including
the literal endpoint $h=0$ by the homogeneous formula.

\paragraph{Matched angle but unmatched state.}
It remains to consider $\bar u=\bar k$. At $u=k$, (OA8) factors as
\[
 S=(\nu Z-kW)^2+
       \left(\frac{k}{s}\nu W-sk^2Z\right)^2.          \tag{OA12}
\]
Use the exact roots
\[
 r_1=vZ-\eta kW,\qquad
 r_2=\frac{k}{s}vW-\eta sk^2Z,
\]
together with the two completed hard roots
$h_1=x-(t/K)D_uvW$, $h_2=y-(a/K)D_uvZ$, and $\eta$.
All six roots in (OA9) are bounded linear combinations of these five.
The exact relation
\[
 Wr_1-\frac{s}{k}Zr_2
   =\eta k(s^2Z^2-W^2)=p\eta                         \tag{OA13}
\]
shows that, if $\bar p\ne0$, adding the identically zero square
$(p\eta-Wr_1+(s/k)Zr_2)^2$ makes the five-root Gram strict.
This operation never inverts $p$.

After pure testing state normalization, any state with positive
parent residue already has a one-root certificate. Otherwise the
four strong events force $\bar x=\bar y=0$, so $W$ or $Z$ is a unit.
If $\bar p=0$, then $Z$ is a unit and we may put $Z=1$.
The remaining choices are $\bar W=\pm\bar s$. The negative choice
is a direct nonnegative-core exit. Indeed set
$w=(W+uY)/D_u=W-uy/D_u$. The pure testing core Schur calculation has hard block
\[
 \begin{pmatrix}1+v^2+u^2v^2&uv\\uv&(D_u+v^2)/D_u\end{pmatrix}\succ0
\]
and positive Schur denominator
$\mathcal D=u^2v^4+u^2v^2+u^2+v^4+2v^2+1$.
Its remaining scalar has positive factor $uvD_uD_v/\mathcal D$
and is proportional to
\[
 D_uuv w^2-2D_uD_vw+uv.
\]
For $u,v>0$ it is nonnegative if $w\leq0$. Hence a negative-core
test must have $w>0$; with $\bar y=0$ it cannot have
$\bar W=-\bar s$. A nonnegative core is itself a pure testing
square after evaluation, and each original event is already a square.
Literal $u=0$ or $v=0$ is included in this direct exit.

\paragraph{The contact: five roots and a finite Taylor argument.}
We have reduced to $Z=1$ and
\[
 u-k,\ W-s,\ v,\ x,\ y\text{ infinitesimal}.
\]
Set
\[
 x_c=\eta kts^2,\quad y_c=\eta ksa,\quad v_c=\eta ks,
\quad U=u-k,\quad B=W-s,
\]
\[
 N_1=x-x_c,\quad N_2=y-y_c,\quad N_3=v-v_c.
\]
The exact roots for this chart are
\[
                 N_1,\ N_2,\ N_3,\ \eta U,\ \eta B.   \tag{OA14}
\]
Here is the full allocation argument, including the normal terms.
As a polynomial in $x,y,v$, (OA5) belongs to
\[
              (x,y,v)^2+\eta(x,y,v)+(\eta^2)
\]
over the bounded polynomial coefficient ring in $u,W$ and the source
unit inverses already listed. Since $x_c,y_c,v_c$ are $\eta$ times
bounded source elements, the shifted polynomial belongs to
$(N_1,N_2,N_3)^2+\eta(N_1,N_2,N_3)+(\eta^2)$.
Its constant normal coefficient is $\eta^2$ times a polynomial in
$U,B$ whose constant and linear terms vanish by exact contact and
angle stationarity. Its linear normal coefficients are $\eta$ times
polynomials in $U,B$ whose constant terms vanish by exact state and
Bob stationarity. These conclusions can be read in the coefficient
field, where $\eta\ne0$; the resulting quotients are bounded by the
displayed factorization, not by cancellation of a mixed element.
Monic Taylor division by $U,B$ therefore allocates these terms to
products of the roots (OA14). Terms of normal degree at least two
are allocated to their first two normal factors. This is an exact
finite Gram construction with bounded entries.

All terms beyond weighted degree two carry one of the infinitesimal
deviations $N_i,U,B$, or a coefficient difference of zero residue.
The Gram residue is consequently the weighted Hessian of (OA6) at
\[
 (u,W,\nu,\xi,\chi)=(k,s,ks,kts^2,ksa),\qquad Z=1.
\]
Its $\xi,\chi$ block is positive definite by (OA7).
The Schur complement is $8aK$ times the half-Hessian of $S$ in
$(\nu,W,u)$ at this point. That half-Hessian is
\[
 \begin{pmatrix}
 K&k(k^2-1)&s(k^2-1)\\
 k(k^2-1)&k^2K&sk(k^4+1)/K\\
 s(k^2-1)&sk(k^4+1)/K&s^2K
 \end{pmatrix},                                      \tag{OA15}
\]
with leading principal minors
\[
                    K,\qquad4k^4,\qquad12s^2k^6/K.
                                                        \tag{OA16}
\]
These are positive units. The five-root Gram is strict, which proves
the contact chart without assuming that $v/\eta$, $x/\eta$, or
$y/\eta$ is bounded.

\paragraph{Exhaustion and restoration.}
For completeness, if the homogeneous Bob sine $v_1$ has nonzero
residue, the four strong packet amplitudes have trivial common state
kernel at every projective Alice angle. Indeed $x=y=0$ and the other
two amplitudes then give
$v_1(u_0W+u_1Y)=v_1(u_1X-u_0Z)=0$; since
$u_0^2+u_1^2$ is a testing unit, all state coordinates vanish.
Their four-root Gram absorbs the infinitesimal signed coefficient;
all other event squares may be retained. Thus only the Bob-small
chart used above remains. The finite and reciprocal Alice charts,
the two state branches, and the exact contact chart exhaust it.
Every strict bounded residue Gram is tensor SOS by the established
strict-Gram lemma. All projective multipliers are pure testing squares
and are restored by separate field square roots. Complex states split
into their real and imaginary parts.

\paragraph{Exact advance and remaining scope.}
The theorem closes the regular ordinary-$k$, comparable-$f,g$ source
slice (OA1), including the previous explicit $k=1$ entry obstruction.
It supplies a general mechanism: two completed strong squares, a
scalar capacity factorization locating the unique angular degeneracy,
and finite weighted contact division with a strict five-root Hessian.
It makes no claim for a nonunit $s$, a nonunit shape or gap among
$a,t,M,J$, non-small $\eta$ or $\zeta$, or the still-unclosed
small-$k$ unit/unbounded-$\eta$ source rows.


\section{\texorpdfstring{The C1 source coordinates and the $k^2\eta$ bound}{The C1 source coordinates and the k2eta bound}}
\label{component:pa133fae35e52}


Throughout, $d$ is the original positive signed coefficient. The event
coefficient at the fourth cycle position is denoted $A_{\rm ev}$, and
$D_{\rm ev}=e_5$. These must not be confused with actual state coordinates.
The source has the canonical omissions $\lambda=n=0$ and is normalized
by $c=e_3=2$. We use the canonical actual-contact sign entry, including
its admitted single-zero $D_{\rm ev}$ or $\beta$ boundaries. Explicitly,
the retained $m,c,A_{\rm ev},\alpha$ are positive, and $D_{\rm ev},\beta$
are positive or exactly one of them is zero. The noncoercive further-zero
faces are already separate exits. The sign and nonsingularity input is the native-entry theorem EN in Chapter~\ref{ch:entry}, proved for every finite noncommuting physical contact in this exact coefficient domain.

\paragraph{Exact transport from the actual contact.}
Write its positive-angle contact as
\[
 u_c=k,\quad v_c=\ell,\qquad
 \psi_c=(S,k(\theta-1),kSr,1),\quad
 S>0,\quad 0<r,\theta<1.
\]
The exact C1 reflection lets us orient $0<k\le1$. Define source-only
parameters
\[
 \eta=\frac{\ell}{kS},\quad \zeta=\frac{S\ell}{k},\quad
 f=\sqrt\eta,\quad g=\sqrt\zeta,\quad
 e=\frac{k}{S}=\frac{kf}{g},\quad
 a=\frac{1-r}{\eta},\quad t=\frac\theta\zeta.
                                                               \tag{CE1}
\]
Thus the classifier $\mathfrak T=k/(S\ell)$ is exactly $\zeta^{-1}$,
and its old state-ratio variable $B$ is exactly $\eta^{-1/2}$.
Put $\kappa=k^2$, $K=1+\kappa$. The original amplitude signs imply
\[
 R=1+\kappa r,\quad T=1+\kappa(1-\theta),\quad
 M=R-t>0,\quad J=T-a>0,
\]
\[
 E=1-\kappa\eta(1-\theta)>0,\qquad
 Q_c=1-\kappa\zeta r>0.                              \tag{CE2}
\]
Indeed the positive second amplitude is $k\zeta M$, the negative
fifth amplitude is $-\ell J$, the state gap is $SE$, and the remaining
state gap is $Q_c$. In particular $0<a,t<2$, without any shape-unit
hypothesis.

The seven native coquotients used by the uniform theorem are valid
under (CE1), not just asymptotically. Their constructed contact is
\[
 (W,X,Y,Z)=(\Delta k,-k(1-t\zeta),\Delta k^2(1-a\eta),1),
 \qquad \Delta=e^{-1}=S/k,
\]
which is exactly the original contact. The native contact identities
prove all six kernel and angular equations. To make the nonsingularity
premise explicit, put $Y_0=Sr$, $h_0=S+k\ell(\theta-1)$,
$R_2=\ell(S+k^2Y_0)-k\theta$, and
$R_5=k(S-Y_0)+k^2\ell(\theta-1)-\ell$. The inherited selected
six-column contact system, on $(d,m,A_{\rm ev},D_{\rm ev},\alpha,\beta)$
with $c=2$, has determinant
\[
 -8192\theta k^2\ell(1+k^2)^2(1+\ell^2)^2\mathcal D,
\]
\[
 \mathcal D=Y_0(1+k^2)(kY_0+\ell)(1-k\ell Y_0)h_0R_2(-R_5)>0.
\]
Each factor is strictly positive by the stated source sign theorem.
Thus the normalized coefficients equal the unique original source
coefficients in that inherited system. This
establishes full original-coefficient transport, including the exact
zero openings; it does not use an archived completion label.

\paragraph{A bounded polynomial consequence of the original fifth event.}
Assume now
\[
                 k\text{ infinitesimal},\qquad
                 \zeta=\mathfrak T^{-1}\text{ infinitesimal}.     \tag{CE3}
\]
Set $q=\kappa\eta$. Since $t<2$ and $E=1-q(1-t\zeta)>0$, $q$ is
bounded and $\bar q\le1$. The original fifth-event formula is
\[
 D_{\rm ev}=\frac{a}{Jc_3}(2c_5-J\eta c_d).
\]
Using $L_c=r$, $H_c=r+\eta$ and $V_c=q\zeta$, its sign is precisely
\[
 \mathcal B_5:=K(1+q\zeta)(r+\eta)-J\eta P_d\ge0.       \tag{CE4}
\]
Although $\eta$ need not be bounded, this whole expression has the
following bounded polynomial form. Define
\begin{align*}
 H={}&2+\kappa r+\zeta[K-(\kappa+2)t]
     +q\zeta(-\zeta t^2+t+1)
     +\kappa\zeta(1-r)(t-1),\\
 P_d={}&1+q\zeta+\kappa H.
\end{align*}
Then exact cancellation using $a\eta=1-r$ and $q=\kappa\eta$ gives
\[
 \mathcal B_5=Kr(1+q\zeta)+(1-r)P_d
       +q\{\zeta t(1+q\zeta)-(K-\kappa\zeta t)H\}.     \tag{CE5}
\]
Every variable on the right is bounded. Hence (CE4) gives
\[
       1-2\bar q\ge0,\qquad 0\le\bar q\le\tfrac12.     \tag{CE6}
\]
This is obtained before dividing by $1-q$ or by a coefficient opening.
In particular $E$ is a positive unit with residue at least $1/2$.

The fourth coquotient polynomial likewise has the bounded expression
\[
\begin{split}
 P_4={}&K(1-q)+\kappa(1-r)-\kappa\zeta t\\
 &+q\zeta(-\kappa\zeta t^2+2\kappa t+t)
   +\kappa\zeta(1-r)(1-t).
\end{split}                                                       \tag{CE7}
\]
Thus $\bar P_d=1$ and $\bar P_4=1-\bar q\ge1/2$.
These exact low-degree identities are the physicality input needed
for the entry, rather than an estimate on a truncated coefficient.



\section{\texorpdfstring{The unbounded-$e$ C1 source away from the zero-$W$ state}{The unbounded-e C1 source away from the zero-W state}}
\label{component:p7ef919c24082}

\label{p7ef919c24082:sec:c1-unbounded-e-w-unit}

The complementary state chart has a short retained-current proof.  It
does not require a contact division or a comparison of coefficient and
testing rates.

\begin{theorem}[Unbounded-$e$, $W$ a unit]
Use the unchanged canonical uniform C1 source of
Section~\ref{p05180564ed05:sec:productrepair}, with $k,f,g>0$ infinitesimal,
$e=kf/g$ unbounded, and the original positive bounded shapes and gaps
$a,t,M=R-t,J=T-a$.  Retain all original nonnegative event coefficients.
Then the entire parent is tensor-SOS for every independent test with
$u,v,X,Y$ infinitesimal, $Z$ bounded, and $W$ a unit.
\end{theorem}
\begin{proof}
Normalize $W=1$ using the testing field and put $\Delta=1/e$.
Evaluate the seven short source polynomials with their native
$K=k^2$, $\eta=f^2$, $\zeta=g^2$ and write $c_b=c_\beta$.
Thus
\[
 R=1+k^2(1-af^2),\qquad T=1+k^2-k^2g^2t,
\]
and their residues are $2,2,2,2,2,2,4$ in the original order,
irrespective of the size of $e$.  These are the same polynomial source
identities as before, now used on the unbounded-$e$ domain.
Write $D_u=1+u^2$, $D_v=1+v^2$ and
\[
 Q=W+uY,\quad B=X+uZ,\quad r=uW-Y,\quad
 r_2=vQ-B,\quad r_5=r-v(Z-uX).
\]
Thus $Q$ is a pure-testing unit.  All appearances of $Q^{-1}$ below
are inverses in the testing field, with bounded residue one.
Let $H_2=4Tc_2$, $H_b=4Tc_b$, $H_3=4Rc_3$, $H_5=4Rc_5$.
The exact original source, with $g=\Delta kf$, is
\[
\begin{gathered}
c=2,\quad d=\frac{2af^2c_d}{c_3},\quad
m=\frac{2ae^2c_2}{Mc_3},\quad
A=\frac{2af^4c_4}{c_3},\quad
\alpha=\frac{2a\Delta^2k^6f^4c_\alpha}{c_3},\\
D+\frac d2=\frac{2ac_5}{Jc_3},\qquad
\beta+\frac d2=\frac{2ae^2c_b}{tc_3},\qquad
\Lambda=\frac{RTc_3}{2a}.
\end{gathered}                                                    \tag{UW1}
\]
The letter $D$ in (UW1) is the original coefficient $e_5$; it is
distinct from the displayed testing factors $D_u,D_v$.
Directly combining the three positive signed-core squares with their
matching events gives the full-parent identity
\[
\begin{split}
\Delta^2\Lambda F={}&
 H_2\frac RM r_2^2+H_b\frac RtD_vB^2
 +\Delta^2H_3\frac TaD_vr^2
 +\Delta^2H_5\frac TJr_5^2\\
&+4RT\Delta^2f^4c_4D_uD_vY^2
 +4RT\Delta^4k^6f^4c_\alpha D_u(vY-Z)^2\\
&+2RT\Delta^2f^2c_dD_u
       (v^2Y^2-2vYZ-v^2Z^2).
\end{split}                                                       \tag{UW2}
\]
Use $R/M=1+t/M$, $R/t=1+M/t$, $T/a=1+J/a$ and
$T/J=1+a/J$.  Retain the four positive excess packets, the complete
alpha packet, and the positive $v^2Y^2$ part of the last line.
The remaining expression has the five-root frame
\[
 z=(r_2,B,\Delta r,\Delta r_5,\Delta f^2Y).
\]
Its diagonal part is
\[
 H_2z_1^2+H_bD_vz_2^2+H_3D_vz_3^2+H_5z_4^2
       +4RTc_4D_uD_vz_5^2.
\]
The two complete negative terms are, exactly,
\[
 -4RTc_dD_u\frac{\Delta Z}{Q}\,z_5(z_1+z_2)
 -2RTc_dD_u\frac{\Delta^2f^2Z^2}{Q^2}(z_1+z_2)^2,       \tag{UW3}
\]
because $vQ=z_1+z_2$.  Every displayed Gram coefficient is bounded.
The correction coefficients in (UW3) have zero residue, while the
diagonal residue is $8I_5$.  The bounded strict-residue Gram lemma
therefore applies.  Restoring the six retained nonnegative packets and
undoing the positive pure-coefficient factor $\Delta^2\Lambda$ proves
the assertion.  No shape or gap is inverted in the bounded Gram;
the possibly unbounded gap ratios occur only in retained squares.
\end{proof}

\begin{corollary}[Precise remaining unbounded-$e$ test domain]
For this source, the only independent testing domain not covered by
the preceding theorem and the existing uniform outer-angle identities
is
\[
 Z=1,\qquad u,v,W,X,Y\text{ infinitesimal}.
\]
\end{corollary}
\begin{proof}
The two strength ratios $d/c,d/D$ and the original harmonic mass
formula (UA1) depend only on $k,f,g$ being infinitesimal.  In particular
the mass has residue $1/4$ for unbounded $e$ as well.  Hence the two
outer-angle lemmas apply without change, including homogeneous
infinite-angle endpoints.  The literal axes give a nonnegative
pure-testing core.  Normalize a nonzero state by a coordinate of
maximal absolute value.  If $X$ or $Y$ is a unit, that core has positive
residue $2(X^2+Y^2)$.  Otherwise $X,Y$ are infinitesimal; the theorem
covers $W$ a unit, and its complement has $Z$ a unit and $W$
infinitesimal.  Pure-testing $Z$ normalization gives the stated domain.
\end{proof}


\section{\texorpdfstring{Two outer-current identities and their source hypotheses}{Two outer-current identities and their source hypotheses}}
\label{component:p6cad7544d1b8}

\label{p6cad7544d1b8:sec:c1-uniform-all-tests}

The two state charts can be joined without importing an all-tests atlas
label. We first record two short outer-angle identities.

\begin{lemma}[Two strong currents]
Suppose the original event coefficients are nonnegative, $d>0$, and
$d/c,d/D$ are infinitesimal, where $c=e_3$, $D=e_5$.
Then every testing chart with $1/v$ bounded is tensor-SOS, for every
other angle and every state.
\end{lemma}
\begin{proof}
Put $g=cD/(c+D)$, $h=D/(c+D)$, $R=uW-Y$,
$C=v(uX-Z)$, $B=X+uZ$, and $N=R+hC$.
Let $D_u=1+u^2$, $D_v=1+v^2$ and set
\[
 K_0=2\sqrt g\,C,\qquad
 B_0=\sqrt{2d}\sqrt{D_v/D_u}\,B,\quad
 c_u=D_u^{-1/2},\quad t_u=uD_u^{-1/2},\quad
 \eta_v=\sqrt{D_v}/v.
\]
The complete selected parent is exactly
\[
\begin{split}
4cD_vR^2+4D(R+C)^2+d\mathcal H
={}&4(c+D)N^2+4cv^2R^2+2d(R+C)^2\\
 &+2dD_u(vY-Z)^2+(K_0,B_0)Q(K_0,B_0)^T,
\end{split}
\]
where
\[
 Q=\begin{pmatrix}
 1-\dfrac{d}{2g}c_u^2\eta_v^2&
 \sqrt{\dfrac{d}{2g}}t_u\eta_v\\
 \sqrt{\dfrac{d}{2g}}t_u\eta_v&1
 \end{pmatrix}.
\]
To check the identity, use $Z=(uB-C/v)/D_u$ and complete the two current
squares. Since $d/g=d/c+d/D$ is infinitesimal and $c_u,t_u,\eta_v$
are bounded, $\bar Q=I_2$. Restore every unselected positive event.
All roots separate between fields. Clearing the original pure-testing
homogeneous angle factors gives the same identity at literal infinity:
$\eta_v^2=1+(1/v)^2$, and $c_u,t_u$ are the normalized homogeneous
Alice coordinates. Thus no infinite-angle endpoint is excluded.
\end{proof}

\begin{lemma}[Critical current mass outside the small Alice chart]
Under the preceding strength assumptions, suppose
\[
 \gamma=\frac{A g}{d^2}\quad\hbox{is bounded},\qquad \bar\gamma\geq\tfrac14,
\]
where $A=e_4$ is the original event coefficient. Then every independent
test with $v$ infinitesimal and $1/u$ bounded is tensor-SOS, for every
state, including the homogeneous Alice-infinite endpoint.
\end{lemma}
\begin{proof}
Keep $R,C,B,N,g,h,c_u,t_u$ from the preceding proof, and put
$\epsilon=d/g$, $q=\sqrt\epsilon$ and
\[
 K_1=2\sqrt g\,C+\sqrt{d\epsilon}\,Y,\quad
 B_1=\sqrt d(B-uvY),\quad V_0=\sqrt d\,vY,
 \quad Y_0=\sqrt{d\epsilon}\sqrt{D_u}\,Y.
\]
Define
\[
 L_0=\frac{qc_u}{2}(K_1-c_uY_0)-t_uvB_1+c_uvV_0.
\]
Direct multiplication gives $L_0=\sqrt{dD_u}\,v(vY-Z)$ and the exact
full selected-parent identity
\[
\begin{split}
&4cD_vR^2+4D(R+C)^2+4AD_uD_vY^2+d\mathcal H\\
&\qquad=4(c+D)N^2+4cv^2R^2+2d(R+C)^2+\mathcal G,\\
\mathcal G={}&K_1^2+2D_v(B_1^2+V_0^2)
 +(4\gamma D_v-c_u^2)Y_0^2\\
&+2qv(K_1-c_uY_0)V_0-2L_0^2.
\end{split}
\]
The four-root Gram of $\mathcal G$ has residue
$\operatorname{diag}(1,2,2,4\bar\gamma-\bar c_u^2)$.
Its last entry is at least $1-\bar c_u^2=\bar t_u^2>0$ because $1/u$
is bounded. Every correction is polynomial in bounded
$q,v,c_u,t_u,\gamma$ and has zero residue. The strict-residue lemma and
restoration of the other events prove the claim. The displayed bounded
homogeneous factors also cover literal Alice infinity.
\end{proof}

\paragraph{The exact uniform-source premises.}
Use the original C1 source with $k,f,g$ positive infinitesimals,
$e=kf/g>0$ arbitrary, positive bounded shapes $a,t$, and
$R=1+k^2(1-af^2)$, $T=1+k^2-k^2g^2t$, $M=R-t>0$, $J=T-a>0$.
The native coquotients are those in the source dictionary above.

First derive the two outer-angle premises from the original coefficients.
Write $B_c=2c_5-Jf^2c_d$. The exact source gives
\[
 D=\frac{aB_c}{Jc_3},\qquad
 \frac dD=\frac{2Jf^2c_d}{B_c},\qquad
 \frac dc=\frac{af^2c_d}{c_3}.
\]
Since $\bar B_c=4$, both ratios are infinitesimal, independently of
any loss of residue of $a$ or $J$. Also, for the original harmonic
mass $g_B=cD/(c+D)$,
\[
 \frac{A g_B}{d^2}
 =\frac{c_4c_3B_c}{c_d^2(2Jc_3+aB_c)}.              \tag{UA1}
\]
The denominator is a unit: its residue is
$16\cdot4(\bar J+\bar a)=64$, since $J+a=T$ has residue one.
The numerator has residue $16$. Thus the original mass in (UA1) has
residue exactly $1/4$. This argument uses neither a strong first pair
nor beta positivity beyond the original source hypotheses.

The two lemmas therefore cover every testing pair outside
the double-infinitesimal angle chart. The complete local and state
join is proved in the uniform-source certificate UJ, and the all-$e$
extension in UE. No historical local certificate is used here.


\section{\texorpdfstring{Endpoint hermite}{Endpoint hermite}}
\label{component:p935c9d6318ea}
\paragraph{Deriving the scalar and first-normal inputs by monic division.}
Here is a finite construction of every coefficient in the preceding table. It uses the seven factored packets, differentiation, and division by monic polynomials. Thus neither a stored expansion nor a rational-function ideal calculation is an input.

Let $\mathcal R=\Q[a,t,b,e,k]$ and use an indeterminate $U$ for $u^2$ in this paragraph only. Put $K=k^2$. Substitute $(X,Y,v)=(X_g,Y_g,v_g)$ in the peeled polynomial $E$. Directly from the seven packets, the resulting expression has the form
\[
\frac{E_g}{e^2b^4}=e^2\mathcal A(U)+ekZ\mathcal B(U)+k^2Z^2\mathcal C(U),
\qquad \mathcal A,\mathcal B,\mathcal C\in\mathcal R[U].
\]
The displayed division means coefficientwise polynomial extraction: the constant, linear and quadratic coefficients in $Z$ have the factors $e^2$, $ek$ and $k^2$, respectively. No evaluated tensor is divided. Substituting the graph into the factored packets gives the five short contact and axis identities
\begin{gather*}
(\mathcal A+\mathcal B+\mathcal C)(K)=0,\qquad
(\mathcal B+2\mathcal C)(K)=0,\qquad
(\mathcal A+\mathcal B+\mathcal C)'(K)=0,\\
\mathcal A(0)=\mathcal B(0)=0.
\end{gather*}
Consequently the following monic divisions have zero remainder:
\begin{align*}
\mathcal N&=\frac{\mathcal A+\mathcal B+\mathcal C}{(U-K)^2},&
\mathcal D&=\frac{\mathcal B+2\mathcal C}{U-K},\\
\mathcal R_0&=\frac{\mathcal C+K\mathcal D+K^2\mathcal N}{U},&
\mathcal J&=\frac{-\mathcal D-2K\mathcal N}{U}.
\end{align*}
For the last two divisions, evaluate the numerators at $U=0$ and use the first two division identities together with $\mathcal A(0)=\mathcal B(0)=0$. This is a polynomial-ring argument; the symbols $K$ and $U$ are not inverted after evaluation. Set
\[
A_{00}=\mathcal R_0(u^2),\qquad A_{01}=u\mathcal J(u^2),
\qquad A_{11}=\mathcal N(u^2).
\]
Comparing the three coefficients of $1,Z,Z^2$ proves exactly
\[
\frac{E_g}{e^2b^4}
=A_{00}\,[u(e-kZ)]^2+A_{01}\,u(e-kZ)(eu^2-k^3Z)
 +A_{11}(eu^2-k^3Z)^2.
\]
Multiplying by $e^2b^4$ gives the scalar identity used above.

For each $\xi\in\{X,Y,v\}$, let $R_\xi$ be the graph restriction of $\partial_\xi E$, divided coefficientwise by $e^2b^2,eb^2,e^2b^2$, respectively. The first two are linear in $Z$. The third has one quadratic coefficient, supplied solely by the signed packet:
\[
[Z^2]R_v=-4R_cT_0c_d\,e k b^2u(1+u^2).
\]
Define
\[
C_X=C_Y=0,\qquad C_v=4R_cT_0c_d\,eb^2(1+u^2).
\]
Then $\widetilde R_\xi=R_\xi-C_\xi Z\,u(e-kZ)$ is linear in $Z$ and has the polynomial parity form
\[
\widetilde R_\xi=eu\,L_{\xi,0}(U)+kuZ\,L_{\xi,1}(U),
\qquad L_{\xi,i}\in\mathcal R[U].
\]
In particular, subtracting the quadratic term also leaves its corrected linear coefficient divisible by $ku$; this follows upon collecting the signed packet with the two current packets. The contact identity is
$\big(L_{\xi,0}+L_{\xi,1}\big)(K)=0$. Form the monic quotient
\[
J_\xi=\frac{L_{\xi,0}+L_{\xi,1}}{U-K}.
\]
The promised full first-normal coefficients are now defined by the short formulas
\[
P_\xi=(L_{\xi,0}-UJ_\xi)(u^2)+C_\xi Z,
\qquad Q_\xi=uJ_\xi(u^2).
\]
Indeed direct multiplication gives
$R_\xi=P_\xi u(e-kZ)+Q_\xi(eu^2-k^3Z)$.
This proves the full first-normal identities, including the quadratic-in-$Z$ term; no first-normal term has been omitted.

For completeness, the weighted splits have an equally explicit finite rule. In a polynomial expanded in $(e,k,b,u)$, place a monomial with $k$-exponent at least $r$ in the first bucket and divide it by $k^r$; place every remaining monomial in the second bucket and divide it by $e^s b^s$. Use $(r,s)=(4,2)$ for $A_{00}$ and $(r,s)=(2,1)$ for
\[
A_{01},\quad P_X,\quad P_Y,\quad P_v-16kP_0.
\]
Substitution of the seven short packets into the monic formulas shows that every monomial left for the second bucket has the required $e^s b^s$ factor. Equivalently, the only polynomial remainders to check are
\[
A_{00}\equiv0\pmod{(k^4,e^2b^2)},\qquad
A_{01},P_X,P_Y,P_v-16kP_0\equiv0\pmod{(k^2,eb)}.
\]
These are ordinary finite polynomial truncations: discard the monomials divisible by the displayed generators. The two bucket quotients have respectively residues $(-8,-8)$ for $A_{00}$ and $(0,0)$ for each of the other four entries; $A_{11}$ has residue $8$, and all $Q_\xi$ have residue zero. To obtain these numbers, after the divisions set $e=k=u=Z=0$, keeping $a,t,b$ free. In particular the new term $C_vZ$ belongs to $eb\,\mathcal R[u,Z]$ and its second-bucket quotient contains the infinitesimal testing factor $Z$. No assumption $\bar b=0$ is used.

This is a reproducible coefficient construction: the source consists of seven short packets, the scalar construction of four monic quotients, and each first-normal construction of one monic quotient. The finite checks accompany these formulas as a convenience; the definitions and the multiplication identities, rather than expanded coefficient arrays, are the proof objects.


\section{\texorpdfstring{Endpoint source}{Endpoint source}}
\label{component:p7d4e5164e84a}
\paragraph{Seven short source polynomials.}
The following definitions make the endpoint input reproducible without a stored expansion. Put
\[
K=k^2,\quad \eta=e^2b^2,\quad\zeta=k^2b^2,\quad V_c=K\eta\zeta.
\]
\[
L_c=1-a\eta,\qquad H_c=1+(1-a)\eta,\qquad Q_c=1-K\zeta+aV_c.
\]
Define
\begin{align*}
P_4={}&1+K+\eta K(a-K-1)-\zeta Kt
 +V_c[-\zeta Kt^2+2Kt+a(1-t)+t],\\
P_\alpha={}&1+K-a\eta K+\zeta[(1+K)-(2+K)t]
 +V_c[-\zeta t^2+a(t-1)+t],\\
P_d={}&(1+K)^2-a\eta K^2+\zeta K[(K+1)-(K+2)t]\\
&+V_c[-\zeta Kt^2+Ka(t-1)+Kt+K+1].
\end{align*}
The seven coquotients are
\begin{align*}
c_2&=2(1+K)(1-t\zeta)L_c(1+V_c)H_cQ_c,\\
c_\beta&=2(1+K)L_c[1+(1-t)\zeta]H_cQ_c,\\
c_3&=2(1+K)L_cH_cQ_c(1-K\eta+tV_c),\\
c_5&=2(1+K)L_cH_cQ_c(1+V_c),\\
c_4&=2Q_cP_4,\qquad c_\alpha=2L_cH_cP_\alpha,
\qquad c_d=4L_cQ_cP_d.
\end{align*}
Their residues are respectively $2,2,2,2,2,2,4$, so all are positive coefficient units on this domain. With $M=R_c-t$ and $J=T_0-a$, the unchanged source is
\begin{gather*}
\lambda=n=0,\quad c=2,\qquad
d=\frac{2a\eta c_d}{c_3},\quad m=\frac{2ae^2c_2}{Mc_3},\quad
A=\frac{2a\eta^2c_4}{c_3},\quad
\alpha=\frac{2aK^2\eta\zeta c_\alpha}{c_3},\\
D+\frac d2=\frac{2ac_5}{Jc_3},\qquad
\beta+\frac d2=\frac{2ae^2c_\beta}{tc_3},\qquad
\Lambda=\frac{R_cT_0c_3}{2a}>0.
\end{gather*}
The theorem retains the hypothesis $D,\beta\ge0$. Write
\[
r_2=v-X+uvY-uZ,\quad r_\beta=X+uZ,\quad r_3=u-Y,
\quad r_5=u+uvX-Y-vZ,\quad r_\alpha=vY-Z,
\]
and put $q_0=v_g$, $z_0=\eta u$. Then $G$ is exactly the sum of
\begin{align*}
G_2&=4T_0e^2c_2[(r_2+tq_0)^2-tR_cq_0^2],\\
G_\beta&=4T_0e^2c_\beta(1+v^2)[(r_\beta+Mq_0)^2-R_cMq_0^2],\\
G_3&=4R_cc_3(1+v^2)[(r_3+Jz_0)^2-T_0Jz_0^2],\\
G_5&=4R_cc_5[(r_5-az_0)^2-aT_0z_0^2],\\
G_4&=4R_cT_0\eta^2c_4(1+u^2)(1+v^2)Y^2,\\
G_\alpha&=4R_cT_0K^2\eta\zeta c_\alpha(1+u^2)r_\alpha^2,\\
G_d&=2R_cT_0\eta c_d(1+u^2)[r_\alpha^2-(1+v^2)Z^2].
\end{align*}
These seven factored expressions agree separately with the defining source. Their proof requires multiplication of short polynomials only.


\section{\texorpdfstring{A product identity repairs the endpoint at every coefficient rate}{A product identity repairs the endpoint at every coefficient rate}}
\label{component:p05180564ed05}
\label{p05180564ed05:sec:productrepair}
The new boundary mechanism is algebraic. A troublesome root is retained, and a second relation controls its product with an additional root. This supplies two positive pivots without comparing an infinitesimal in one field with an infinitesimal in the other.

\begin{lemma}[Product-relation completion]\label{p05180564ed05:lem:productrepair}
Let $z$ be a root vector over $T$, and $h,r\in T$. Suppose
\[
F=(z,h,r)^TQ(z,h,r),\qquad Q\in\Sym(B),\qquad
\bar Q=\operatorname{diag}(Q_0,-a,0),
\]
where $Q_0$ is a real positive-definite matrix and $a\in\R_{\ge0}$. Suppose $p\in B$ has nonzero residue, $A$ is a row with entries in $\mathfrak n$, and
\[
ph=Az+r,\qquad phr=(z,h,r)^TE(z,h,r),\quad E\in\Sym(\mathfrak n).
\]
Then $F\in\Sigma T^2$.
\end{lemma}
\begin{proof}
For a sufficiently large positive real $M$, add the exact zero
\[
M\big[p^2h^2+r^2-(Az)^2-2phr\big]=0.
\]
The residue Gram becomes $\operatorname{diag}(Q_0,M\bar p^2-a,M)\succ0$. The subtracted terms have zero residue by hypothesis. Lemma~\ref{p25af2c58da7e:lem:gram} applies. In particular $p$ is never inverted.
\end{proof}

\subsubsection*{Exact contact scope and full-polynomial input}
This application uses the canonical C1 contact model, with $\lambda=n=0$. Its source parameters are positive $e,k,b,a,t$ satisfying
\[
e,k\in\mathfrak m_c,\quad b\in B_c,\qquad
R_c=1+k^2(1-ae^2b^2),\quad T_0=1+k^2-tk^4b^2,
\]
\[
0<a<T_0,\qquad0<t<R_c.
\]
The signed weight satisfies $d>0$, and all eight event coefficients are nonnegative original physical C1 contact coefficients. This is a specified contact chart, not every coefficient vector with the same omissions. The independent testing chart is $W=1$ with $u,v,X,Y,Z$ infinitesimal. In this section there is no restriction on $k^3/e$ or $k/e$.

The complete seven-packet contact polynomial is denoted $G$. Its short factored definition and unchanged original coefficient vector are printed below. The accompanying checkers compare the entire polynomials, including the scalar and normal terms; no global coverage claim is inferred from those checks.

Define
\begin{gather*}
v_g=ekb^2u,\quad X_g=-uZ+tv_g,\quad Y_g=(1-ae^2b^2)u,\\
N_1=X-X_g,\quad N_2=Y_g-Y,\quad \delta v=v-v_g,\quad D_t=Z-uX,\\
U_0=1+u^2,\quad c_0=1-tk^2b^2,\quad P_0=T_0+u^2=U_0+k^2c_0,\\
S=u+uvX-Y-vZ+(T_0-a)e^2b^2u,\\
s_0=eb^2u(e-kZ),\qquad q=eb^2(eu^2-k^3Z).
\end{gather*}
\par\noindent\textit{Source inclusion. }The original input is reproduced once as Component p7d4e5164e84a (Section~\ref{component:p7d4e5164e84a}).\par
Direct substitution gives the exact nonlinear relation
\[
L=P_0s_0-S+N_2-D_t\delta v+uv_gN_1-uc_0q=0. \tag{E1}
\]
Peel the two complete current squares:
\[
E=G-H_3(1+v^2)N_2^2-H_5S^2,\qquad
H_3=4R_cc_3,\quad H_5=4R_cc_5,\quad \bar H_3=\bar H_5=8.
\]
Each peeled packet divided by $e^2$ is polynomial in bounded variables. The full graph scalar has the identity
\[
E_g=A_{00}s_0^2+A_{01}s_0q+A_{11}q^2,
\]
and the entire first normals, with $g_0=s_0/(eb^2)$ and $g_2=q/(eb^2)$ denoting their displayed polynomial expressions, satisfy
\[
E_X=e^2b^2(P_Xg_0+Q_Xg_2),\quad
E_Y=eb^2(P_Yg_0+Q_Yg_2),\quad
E_v=e^2b^2(P_vg_0+Q_vg_2).
\]
The exact polynomial splits and residues needed are
\begin{center}\small
\begin{tabular}{ccc}\toprule
Quantity&Polynomial split&Quotient residues\\\midrule
$A_{00}$&$k^4a_0+e^2b^2a_1$&$-8,-8$\\
$A_{01}$&$k^2c_1+ebc_2$&$0,0$\\
$P_X$&$k^2p_X+ebq_X$&$0,0$\\
$P_Y$&$k^2p_Y+ebq_Y$&$0,0$\\
$P_v-16kP_0$&$k^2p_v+ebq_v$&$0,0$\\\bottomrule
\end{tabular}
\end{center}
Also $\bar A_{11}=8$ and $\bar Q_X=\bar Q_Y=\bar Q_v=0$. These equalities remain valid with $b$ symbolic; boundedness, rather than infinitesimality, suffices for $b$.

\par\noindent\textit{Source inclusion. }The original input is reproduced once as Component p935c9d6318ea (Section~\ref{component:p935c9d6318ea}).\par

Exact finite Taylor expansion in $(N_1,-N_2,\delta v)$ accounts for every higher-normal term. Because $E/e^2$ is polynomial, every term of normal degree at least two allocates to $x=eN_1$, $y=e\delta v$, $N_2$ with bounded coefficients. Terms of degree at least three have an extra infinitesimal normal. Terms involving $N_2$ retain an $e$ factor. The only additional nonzero residue is
\[
B_0=\begin{pmatrix}16&-8\\-8&8\end{pmatrix}
\quad\text{on }(x,y).
\]
It is visible without expansion of the full Hessian: at $u=Y=Z=0$, the quadratic residue of $E/e^2$ is $8(X-v)^2+8X^2$, with $b$ still free.

\subsubsection*{The new identity and its nine-root proof}
If the signed testing core is nonnegative, its pure testing square root and the original event packets already give SOS. Otherwise the sign gauges $u,v\ge0$ preserve every packet, and the core Schur complement implies $u,v>0$ and $D_t>0$. To check the latter, after $Z=D_t+uX$ the Schur complement on $(W,D_t)$ is
\[
\frac{2(1+v^2)}{N_0}
\begin{pmatrix}
u^2v^2(1+u^2)&-uv(1+u^2)(1+v^2)\\
-uv(1+u^2)(1+v^2)&u^2v^2
\end{pmatrix},
\]
where $N_0=u^2v^4+u^2v^2+u^2+v^4+2v^2+1>0$. The eliminated block is positive definite. For $W=1,D_t\le0$, the remaining form is nonnegative, proving the claim.

Consequently $w=\sqrt{ek}\sqrt{D_t}\,\delta v$ is a legal separated root. Set
\begin{gather*}
h=k^2s_0-uq=e^2b^2u(k^2-u^2),\qquad z=ebs_0,\qquad r=k^2D_t\delta v,\\
A_0=k^2(S-N_2)-k^3b^2u^2x-uU_0q.
\end{gather*}
Multiplication of (E1) by $k^2$ gives $P_0h=A_0+r$. The additional exact identity is
\begin{align*}
P_0hr={}&ku^2(P_0-k^2)y(S-N_2)
-ku(P_0+c_0k^2u^2)yq\\
&+u^2(P_0-k^2)w^2-k^4b^2u^2T_0xy. \tag{E2}
\end{align*}
Every displayed coefficient has zero residue and is a polynomial in separately bounded factors. To derive (E2), substitute the two soft roots to obtain
\[
u(P_0-k^2)s_0-(1+c_0u^2)q
=ekb^2(k^2-u^2)(D_t+uN_1).
\]
Then use $D_t=U_0Z-uN_1-tu v_g$, $T_0=1+k^2c_0$, and (E1) to replace $S-N_2$ in (E2). The terms in $D_t(\delta v)^2$ cancel; the remaining terms give $P_0hr$. This is a whole-polynomial calculation.

Let $J_0$ denote the right-hand side of (E2). Thus
\[
Z_0=P_0^2h^2+r^2-A_0^2-2J_0=0.
\]
Add to the entire $G$ the exact zero expression
\[
-16kyL+16(ebL)^2+16Z_0.
\]
The first term cancels the exceptional $16kP_0$ first normal and supplies $16w^2$. Use $k^2s_0=h+uq$ in every other scalar and first-normal term. The scalar becomes
\[
a_0(h+uq)^2+a_1z^2+[c_1(h+uq)+c_2z]q+A_{11}q^2.
\]
Its only negative residue diagonals are $-8$ on $h,z$. The safe relation
\[
ebL=P_0z-ebS+ebN_2-bD_ty+ekb^3u^2x-ebu c_0q
\]
has only infinitesimal coefficients apart from $P_0$; hence $16(ebL)^2$ changes the $z$ diagonal to $8$. Every coefficient of $A_0$ is infinitesimal, as is every coefficient in (E2). Thus $16Z_0$ changes the $h$ diagonal to $8$ and supplies a new $r$ diagonal $16$, with no other residue change.

The full Gram on
\[
\boxed{(x,y,N_2,S,w,q,h,z,r)}
\]
therefore has residue
\[
\boxed{B_0\oplus\operatorname{diag}(8,8,16,8,8,8,16)\succ0.}
\]
Its first two LDL pivots are $16,4$; the other seven pivots are printed above. The strict-residue lemma proves SOS of the complete seven-packet polynomial for every coefficient ratio $k^3/e$.

\subsubsection*{Restoration and coverage}
The complete original-parent restoration is
\begin{align*}
\Lambda F={}&G
+4T_0e^2c_2\frac{t}{M}(r_2-Mq_0)^2
+4T_0e^2c_\beta\frac{M}{t}(1+v^2)(r_\beta-tq_0)^2\\
&+4R_cc_3\frac{J}{a}(1+v^2)(r_3-az_0)^2
+4R_cc_5\frac{a}{J}(r_5+Jz_0)^2.
\end{align*}
Every additional packet is a nonnegative coefficient factor times a pure-testing factor times a square. This identity is four instances of
\[
S^2+2(A+B)qS+(A+B)Aq^2+\frac BA S^2
=\frac{A+B}{A}(S+Aq)^2.
\]
Use respectively $(A,B,q)=(M,t,q_0),(t,M,q_0),(a,J,z_0),(J,a,-z_0)$ and $S=r_2-Mq_0,r_\beta-tq_0,r_3-az_0,r_5+Jz_0$. Comparing the seven resulting coefficients with the displayed source proves the identity. Since $\Lambda$ is a positive pure coefficient factor, the pure-factor restoration in \hyperref[component:p9b00698ff36a]{RFA13} (p.~\pageref{component:p9b00698ff36a}) recovers the original parent. Normalizing the original state to $W=1$ and undoing that normalization use only testing factors.

This removes the former $k^3/e$ restriction and includes positive bounded $b$. It does not prove entry of every physical C1 or C2 source into this domain. For unbounded $b$, the recovered inverse-$b$ theorem remains a different input with a ten-root, large polynomial certificate. The new short identity should not be presented as a replacement for that certificate or for the full source cover.


\section{\texorpdfstring{Uniform C1: the bounded-reciprocal repair and its scope}{Uniform C1: the bounded-reciprocal repair and its scope}}
\label{component:p41b39e2f94e1}


The preserved uniform-source theorem explicitly assumes $e$ infinitesimal.
The preserved contact--conductor theorem explicitly assumes $e$ unbounded.
This note proves the missing unit-$e$ case by extending the latter proof
to bounded $\Delta=e^{-1}$ and bounded testing $w$. The extension uses
the same exact complete parent, the same retained packets, and the same
unlocalized conductor. The new work is a weight estimate that remains
strict when both $\delta$ and $w$ have weight zero, followed by an
exhaustive testing-chart join. An old all-$e$ status assertion is not a
premise.

\begin{theorem}[Uniform C1, bounded reciprocal $e$]
\label{p41b39e2f94e1:thm:bounded-delta}
Let the coefficient and testing fields be independent real closed
extensions of $\mathbb R$. Use the original canonical uniform C1 source
defined below, with $k,f,g>0$ infinitesimal, $e=kf/g>0$, and
$\Delta=1/e$ bounded. The shapes $a,t>0$ are bounded and satisfy the
original strict gaps $M=R-t>0$, $J=T-a>0$. Retain every original
coefficient sign, including $D_{\rm e},\beta\geq0$, and every literal
coefficient zero. Then the complete parent is tensor-SOS at every
independent projective measurement pair and every real or complex
testing state.
\end{theorem}

Together with the preserved infinitesimal-$e$ theorem, this proves the
whole positive-$e$ uniform source. It does not prove an unlisted source
chart or close any remaining C14 or weak-C2 row. We use the already
admitted separated bounded strict-Gram lemma: a finite bounded separated
Gram with strictly positive real residue gives tensor-SOS. No inversion
of a mixed element and no comparison between an independent coefficient
and testing rate enters the proof.

\subsection*{Original source and complete parent}

Put $\kappa=k^2$, $\eta=f^2$, $\zeta=g^2$, $K=1+k^2$, and
\[
 V_c=\kappa\eta\zeta,\quad L_c=1-a\eta,\quad
 H_c=1+(1-a)\eta,\quad Q_c=1-\kappa\zeta+aV_c,
\]
\[
 R=1+k^2(1-a\eta),\qquad T=1+k^2-k^2\zeta t,
 \qquad M=R-t>0,\quad J=T-a>0.
\]
The native coquotients are defined by the following short polynomials:
\begin{align*}
 P_4={}&1+\kappa+\eta\kappa(a-\kappa-1)-\zeta\kappa t
 +V_c[-\zeta\kappa t^2+2\kappa t+a(1-t)+t],\\
 P_\alpha={}&1+\kappa-a\eta\kappa
 +\zeta[(1+\kappa)-(2+\kappa)t]
 +V_c[-\zeta t^2+a(t-1)+t],\\
 P_d={}&(1+\kappa)^2-a\eta\kappa^2
 +\zeta\kappa[(\kappa+1)-(\kappa+2)t]\\
 &+V_c[-\zeta\kappa t^2+\kappa a(t-1)+\kappa t+\kappa+1],\\
 c_2={}&2K(1-t\zeta)L_c(1+V_c)H_cQ_c,\\
 c_b={}&2KL_c[1+(1-t)\zeta]H_cQ_c,\\
 c_3={}&2KL_cH_cQ_c(1-\kappa\eta+tV_c),\\
 c_5={}&2KL_cH_cQ_c(1+V_c),\\
 c_4={}&2Q_cP_4,\qquad c_\alpha=2L_cH_cP_\alpha,
 \qquad c_d=4L_cQ_cP_d.
\end{align*}
Their residues are $2,2,2,2,2,2,4$. The original source is
\begin{gather}
 \lambda=n=0,\quad C=2,\quad
 d_{\rm B}=\frac{2a\eta c_d}{c_3},\quad
 m=\frac{2ae^2c_2}{Mc_3},\quad
 A=\frac{2a\eta^2c_4}{c_3},\quad
 \alpha=\frac{2a\Delta^2k^6\eta^2c_\alpha}{c_3},\notag\\
 D_{\rm e}+\frac{d_{\rm B}}2=\frac{2ac_5}{Jc_3},\qquad
 \beta+\frac{d_{\rm B}}2=\frac{2ae^2c_b}{tc_3},\qquad
 \Lambda=\frac{RTc_3}{2a}.                                      \tag{1}
\end{gather}
The original event $D_{\rm e}$ is distinct from the testing factor
$D=1+u^2$. For a testing state $(W,X,Y,Z)$ put
\[
 D=1+u^2,\quad D_v=1+v^2,\quad B=X+uZ,\quad Q_A=W+uY,
 \quad r=uW-Y,
\]
\[
 r_2=vQ_A-B,\qquad r_5=r-v(Z-uX).
\]
The original full parent on this source is
\begin{align*}
 F={}&4mr_2^2+8D_vr^2+4ADD_vY^2+(4D_{\rm e}+2d_{\rm B})r_5^2\\
 &+(4\alpha+2d_{\rm B})D(vY-Z)^2
 +(4\beta+2d_{\rm B})D_vB^2-2d_{\rm B}DD_vZ^2.
\end{align*}
With $H_2=4Tc_2$, $H_b=4Tc_b$, $H_3=4Rc_3$, $H_5=4Rc_5$,
equation (1) gives exactly
\begin{align}
 \Phi:=\Delta^2\Lambda F={}&H_2\frac RM r_2^2
 +H_b\frac RtD_vB^2+\Delta^2H_3\frac TaD_vr^2
 +\Delta^2H_5\frac TJr_5^2\notag\\
 &+4RT\Delta^2\eta^2c_4DD_vY^2
 +4RT\Delta^4k^6\eta^2c_\alpha D(vY-Z)^2\notag\\
 &+2RT\Delta^2\eta c_dD\{(vY-Z)^2-D_vZ^2\}.             \tag{2}
\end{align}

\subsection*{The \texorpdfstring{$Z$}{Z}-unit chart with bounded \texorpdfstring{$W$}{W}}

Assume $u,v,X,Y$ are infinitesimal, $Z$ is a unit, and the state has
bounded coordinates. Normalize $Z=1$ in the testing field. The new
argument allows arbitrary bounded $W$. Put
\[
 p=t/R,\quad h_s=a/T,\quad q_A=v(W+uY),\quad q_B=v(1-uX),
 \quad N=B-pq_A,\quad M_B=r-h_sq_B.
\]
Subtract only the four nonnegative packets
\[
 \Pi_1=H_2\frac tM N^2,\quad \Pi_2=H_b\frac MtD_vN^2,
 \quad \Pi_3=\Delta^2H_3\frac JaD_vM_B^2,
 \quad \Pi_4=\Delta^2H_5\frac aJM_B^2.
\]
No boundedness of their gap ratios is assumed. The exact remainder is
\begin{align}
 G:=\Phi-\sum_i\Pi_i={}&H_2\{N^2-2q_AN+(1-p)q_A^2\}
 +H_bD_v\{N^2+2q_AN+pq_A^2\}\notag\\
 &+\Delta^2H_3D_v\{M_B^2+2q_BM_B+h_sq_B^2\}\notag\\
 &+\Delta^2H_5\{M_B^2-2q_BM_B+(1-h_s)q_B^2\}
 +\mathcal W,                                                   \tag{3}
\end{align}
where $\mathcal W$ is the last three terms of (2), with $Z=1$.
Use the pure-testing coordinate $w=(W+uY)/D$. Its full inverse is
\begin{gather}
 X=N+pDvw-u,\qquad q_B=v(D-uN-puDvw),\notag\\
 W=w+\frac{u(M_B+h_sq_B)}D,\qquad
 Y=uw-\frac{M_B+h_sq_B}D.                               \tag{4}
\end{gather}
The only testing denominator is the positive unit $D$.

If the pure-testing signed core is nonnegative, its product with
$d_{\rm B}>0$ is already a separated square and all event packets are
nonnegative. Otherwise the sign symmetries allow $u,v>0$. Explicitly,
changing $u$ to $-u$ and $(Y,Z)$ to $(-Y,-Z)$, or changing $v$ to
$-v$ and $(X,Z)$ to $(-X,-Z)$, preserves every event square and the
core. The negative-core branch then forces $w>0$: completing the two
positive core squares leaves
\[
 \frac{uvDD_v}{\mathcal D}
 \{D u v\,w^2-2DD_vw+uv\},\qquad
 \mathcal D=u^2v^4+u^2v^2+u^2+v^4+2v^2+1>0,
\]
which is nonnegative for $w\leq0$. The eliminated matrix is
\[
 \begin{pmatrix}1+v^2+u^2v^2&uv\\uv&(D+v^2)/D\end{pmatrix}\succ0.
\]
Thus all radicals used below have nonnegative pure-field factors.
Literal angle axes already have a nonnegative core.

\subsubsection*{Polynomial source change and exact scalar}

Define pure-source quantities
\[
 \delta=\Delta R/K,\quad b^2=\eta K/R,\quad
 Q=K+ak^2b^2,\quad T_n=K^2-t\delta^2k^4Qb^2.
\]
Here $\delta$ is positive bounded, $k,b$ are infinitesimal, and
\begin{equation}
 \eta=Kb^2/Q,\quad \Delta=\delta Q/K,\quad
 \zeta=\delta^2k^2Qb^2/K,\quad R=K^2/Q,\quad T=T_n/K.       \tag{5}
\end{equation}
Every denominator in (5) is a positive source unit. After this
substitution put $C_i=K^2Q^3c_i$. These are polynomials: each native
monomial $c k^jf^{2r}g^{2s}a^\alpha t^\tau$ maps to
\[
 c\delta^{2s}k^{j+2s}b^{2(r+s)}a^\alpha t^\tau
 K^{2+r-s}Q^{3-r+s},\qquad -2\leq r-s\leq3.
\]
No negative exponent or shape inverse occurs. Put
\[
 V_n=K^2-tQuvw,\quad Y_n=KT_nuw-avV_n,\quad B_0=vY_n-KT_n,
\]
\[
 A_c=C_5T_n+aK(C_3D_v-C_5),\qquad \mathcal C=K^6Q^4T_n.
\]
Use different letters for the conductor monomials and the event weights:
\[
 A_m=wv,\quad B_m=\delta v,\quad C_m=\delta b^2uw,
 \quad D_0=\delta^2k^3b^2.
\]
Direct substitution of (4) into (3), retaining every term, gives
\begin{align}
 F_0:=(\mathcal CG)|_{N=M_B=0}={}&
 4KQT_n^2[C_2(K^2-tQ)+C_bD_vtQ]D^2 A_m^2\notag\\
 &+4Q^2A_cD^2V_n^2B_m^2\notag\\
 &+4K^3C_4DD_v(KT_nC_m-ab^2B_mV_n)^2\notag\\
 &+4KQ^2C_\alpha DB_0^2D_0^2\notag\\
 &+2K^2QC_dD\{b^2B_m^2(Y_n^2-K^2T_n^2+2aKT_nV_n)
                         -2K^2T_n^2B_mC_m\}.             \tag{6}
\end{align}
Define the leading polynomial and its complete error by
\[
 P=D(A_m^2+B_m^2)+K(C_m^2+D_0^2)-2K^2B_mC_m,
 \qquad E_0=F_0-8K^2D\mathcal CP.
\]
The full first coefficients on $(N,\delta M_B)$ are
\begin{align}
 S_B={}&-8aK^2b^4C_4D_vV_n
       +8\delta^2Q^2k^6b^4C_\alpha B_0+4KQb^2C_dB_0,\notag\\
 S_C={}&8K^3T_nb^2C_4D_v,\notag\\
 L_1={}&A_m[8K^3QT_n^2(C_bD_v-C_2)D]\notag\\
 &+B_m[-8\delta K^2Q^2uD\,vV_nA_c+\delta K^3auvS_B]
   +C_m[\delta K^3auvS_C],\notag\\
 L_2={}&B_m[8K^2Q^2T_nDV_n(C_3D_v-C_5)-K^2T_nS_B]
          -C_mK^2T_nS_C.                               \tag{7}
\end{align}
In particular $E_0\in(A_m,B_m,C_m)^2+(D_0^2)$ and
$L_1,L_2\in(A_m,B_m,C_m)$ without ideal division.

\subsubsection*{Exact contact and the unlocalized conductor}

Work in $\mathcal R=\Rc[u,w,v]$ and put
\[
 \begin{gathered}
 \mathfrak J=(u-k,w-\delta k,v-\delta k^2b^2),\\
 x=B_m-C_m,\quad y=A_m-kC_m,\quad z=uB_m-kC_m,
 \quad h=kC_m-D_0.
 \end{gathered}
\]
The original source contact is
\[
 u=k,\quad v=\Delta k^2\eta,\quad
 (W,X,Y,Z)=(\Delta k,-k(1-t\zeta),\Delta k^2L_c,1).
\]
It becomes $(u,w,v,N,M_B)=(k,\delta k,\delta k^2b^2,0,0)$.
For completeness, its exact contact and derivative identities follow
from the six native balances
\[
 c_b-c_2=\zeta c_3,\qquad c_3-c_5=-\kappa\eta c_b,
 \qquad c_2=(1-t\zeta)c_5,
\]
\begin{align*}
 0={}&2\kappa c_2+2(1+V_c)L_cc_4
         -2\kappa^2\zeta Q_cc_\alpha-Q_cc_d,\\
 0={}&\kappa L_cc_2+c_5+(1+V_c)L_c^2c_4+\kappa Q_c^2c_\alpha
         +\tfrac12(\kappa\zeta L_c^2-2L_c-\eta)c_d,\\
 0={}&\kappa Rc_2+t\kappa V_cc_b+aV_cc_3+Tc_5
       +KV_cL_c^2c_4-\kappa^2\zeta KQ_cL_cc_\alpha
       -\tfrac K2(Q_cL_c+\eta)c_d.
\end{align*}
They follow by multiplying the short $P_4,P_\alpha,P_d$ polynomials.
Substitution into the derivatives of (2) proves that its value and all
first derivatives vanish at the displayed contact. The four packets
also vanish there to second order. Consequently
\[
 E_0\in\mathfrak J^2,\qquad L_1,L_2\in\mathfrak J.       \tag{8}
\]
Here $P\in\mathfrak J^2$ follows also from its exact square identity
below. No asymptotic contact replacement is made.

Adjoin only the separated radicals
\[
 a_1=\sqrt{wv},\quad a_2=\sqrt\delta\sqrt{uv},\quad
 a_3=b\sqrt{\delta k}\sqrt{uw},\quad
 a_4=\delta k^{3/2}b,\qquad q=a_1a_2-a_3a_4.
\]
Only $q^2\in\mathcal R$ is needed in the polynomial lift. Direct
multiplication gives
\begin{align}
 C_m(u-k)&=z-ux,\notag\\
 C_m(w-\delta k)&=\delta y-wx,\notag\\
 C_m(v-\delta k^2b^2)&=b^2(wz+\delta ky-kwx),\notag\\
 r_*:=y+z-h-2kx
 &=(v-\delta k^2b^2)[(w-\delta k)+\delta(u-k)]
       -\delta kb^2(u-k)(w-\delta k),\notag\\
 kC_mr_*&=q^2-yz.                                      \tag{9}
\end{align}
For any polynomial $S$ successive monic divisions give
\[
 S=S_c+(u-k)T_u+(w-\delta k)T_w
              +(v-\delta k^2b^2)T_v.
\]
If $S_c=0$, (9) yields the explicit polynomial lift
\begin{equation}
 C_mS=x(-uT_u-wT_w-b^2kwT_v)
      +y(\delta T_w+\delta kb^2T_v)+z(T_u+b^2wT_v).      \tag{10}
\end{equation}
Thus
\begin{equation}
 (A_m,B_m,C_m)\cap\mathfrak J\subset(x,y,z),\qquad
 [(A_m,B_m,C_m)^2+(D_0^2)]\cap\mathfrak J^2
                      \subset(x,y,z,h,q)^2.            \tag{11}
\end{equation}
The restriction to the unmixed generator $D_0^2$ is substantive.
For clarity, here is the complete constructive argument. In the first
inclusion write $F=R_AA_m+R_BB_m+R_CC_m$ and substitute
$A_m=y+kC_m$, $B_m=x+C_m$. The remaining coefficient of $C_m$ is
$S=kR_A+R_B+R_C$. The contact equality implies $S_c=0$ by cancellation
in the polynomial domain $\Rc$, so apply (10).

For the second inclusion write
$F=\sum_{i\le j}f_{ij}m_im_j+\rho D_0^2$, with
$(m_1,m_2,m_3)=(A_m,B_m,C_m)$. Substitute the same formulas and
$D_0=kC_m-h$ to obtain
\[
 F=Q_2+C_m(xH_x+yH_y+zH_z+hH_h)+C_m^2S,
 \quad Q_2\in(x,y,h)^2,\quad H_h=-2k\rho,\quad H_z=0.
\]
The contact value implies $S_c=0$. Absorb $C_m^2S$ into the four
$H$ coefficients by (10). Since $F\in\mathfrak J^2$, its first
derivatives vanish at contact, while
\[
 \det\left(\frac{\partial(x,y,z)}{\partial(u,w,v)}\right)_c
        =-\delta^4k^4b^4\ne0.
\]
The relation $r_*\in\mathfrak J^2$ has gradient coefficients
$(-2k,1,1,-1)$. Put
\[
 \lambda=2k\rho_c,\qquad J_i=H_i-\lambda(-2k,1,1,-1)_i.
\]
The last $J_i$ has zero contact
value; the determinant and the first derivative equations imply the
other three do too. This uses the adjugate and cancellation only to
prove zero polynomial identities; no determinant inverse is a lifting
coefficient. Finally
\[
 F=Q_2+2\rho_c(q^2-yz)+\sum_{i=1}^4r_i(C_mJ_i),
 \qquad(r_1,r_2,r_3,r_4)=(x,y,z,h),
\]
and each $C_mJ_i$ is lifted by (10). This proves (11), including all
coefficient-zero strata, in the unlocalized ring.

\subsubsection*{The new positive-weight estimate}

Assign weights
\begin{equation}
 (\delta,k,b,a,t,u,w,v)=(0,2,6,0,0,2,0,16).             \tag{12}
\end{equation}
The zero-weight variables are bounded; every positive-weight variable
is infinitesimal in its own field. Positive weight of every polynomial
coefficient therefore makes it a bounded separated infinitesimal,
without any comparison of independent rates. A polynomial
$\Ow(s)$ has all monomials of weight at least $s$.

The short coquotients and (5) give
\begin{align*}
 C_i&=2K^6+\Ow(12)\quad(i=2,b,3,5,\alpha),&
 C_4&=2K^6+\Ow(16),& C_d&=4K^7+\Ow(12),\\
 Q&=K+\Ow(16),& T_n&=K^2+\Ow(20),& V_n&=K^2+\Ow(18),\\
 D_v&=1+\Ow(32),&Y_n&=\Ow(2),&B_0&=-K^3+\Ow(18),\\
 A_c&=2K^8+\Ow(12),&&\mathcal C=K^{12}+\Ow(16).
\end{align*}
The stronger $C_4$ estimate follows because every $\eta$-term in
$P_4$ also contains $\kappa$, and the clearing correction $Q-K$
has weight $16$. The two exact difference identities are
\[
 K(C_b-C_2)=\delta^2k^2Qb^2C_3,\qquad
 Q(C_3-C_5)=-k^2Kb^2C_b.
\]
Both differences therefore have weight at least $16$; since $K,Q$
have constant term one, this conclusion is polynomial and introduces
no inverse into a lifting coefficient.

The monomial weights of $(A_m,B_m,C_m,D_0)$ are $(16,16,14,18)$.
In the five lines of (6), keep the respective leading terms
\[
 8K^{14}D^2A_m^2,\quad 8K^{14}D^2B_m^2,\quad
 8K^{15}DC_m^2,\quad 8K^{15}DD_0^2,\quad
 -16K^{16}DB_mC_m.
\]
The first two errors have weight at least $44$. The third error has
weight at least $42$: its only possibly smaller cross term contains
$b^2B_mC_m$, of weight $42$. The fourth error has weight at least
$48$. In the fifth line the coefficient error on $B_mC_m$ has weight
at least $42$, and the separate $b^2B_m^2$ term has weight at least
$44$. Their sum is $8K^{14}DP+\Ow(42)$. Since $\wt(P)\geq28$
and $\mathcal C-K^{12}=\Ow(16)$, this proves
\[
 E_0=\Ow(42).
\]
Also $S_B,S_C=\Ow(12)$, and (7) gives
\begin{equation}
 L_1=\Ow(32),\qquad L_2=\Ow(26).                       \tag{13}
\end{equation}

It remains to check that the explicit conductor divisions preserve
strict positivity of coefficient weights. The contact substitutions
$u\mapsto k$, $w\mapsto\delta k$, $v\mapsto\delta k^2b^2$
never decrease weight; their weights change respectively as
$2\mapsto2$, $0\mapsto2$, $16\mapsto16$.
For $S=\Ow(s)$, its monic quotients obey
\[
 \wt(T_u)\geq s-2,\quad \wt(T_w)\geq s,\quad
 \wt(T_v)\geq s-16.
\]
The coefficients in (10) on $(x,y,z)$ consequently have weights
at least $(s-2,s-2,s-4)$.

For a first-normal polynomial of weight $p$, monomial division by
$(A_m,B_m,C_m)$ gives quotient weights at least
$(p-16,p-16,p-14)$. Hence $S=kR_A+R_B+R_C$ has weight at least
$p-16$, and the completed lift has coefficient weights at least
$p-20$. From (13) these bounds are $12$ for $L_1$ and $6$ for $L_2$.

For $E_0$ of weight at least $42$, its monomial quotients on
\[
 A_m^2,A_mB_m,A_mC_m,B_m^2,B_mC_m,C_m^2,D_0^2
\]
have weights at least $(10,10,12,10,12,14,6)$. Thus the initial
$Q_2$ coefficients have weight at least $6$, the initial
$(H_x,H_y,H_h)$ have bounds $(10,10,8)$, and $S$ has bound $10$.
The first use of (10) leaves bounds $(8,8,6,8)$ on
$(H_x,H_y,H_z,H_h)$. Since $\rho_c$ has weight at least $6$ and
$\lambda=2k\rho_c$ at least $8$, subtracting the gradient relation
leaves the same bounds on $(J_x,J_y,J_z,J_h)$. A second use of
(10) therefore leaves every scalar-lift coefficient with weight at
least $2$. The separate coefficient $2\rho_c$ of $q^2$ has bound
$6$. Every correction coefficient is thus strictly infinitesimal.

\subsubsection*{The full strict Gram and restoration}

The eleven original soft roots are $g_0=x$, $g_4=uwv-\delta^2k^4b^2$,
the six differences $a_i^2-a_j^2$, and the three differences
$a_ia_j-a_\ell a_m$ over the three pairings of four indices.
They satisfy the exact identity
\begin{equation}
 P=g_0^2+g_4^2+\frac13\sum_{i<j}(a_i^2-a_j^2)^2
       +\frac23\sum_{\{ij,\ell m\}}(a_ia_j-a_\ell a_m)^2. \tag{14}
\end{equation}
The conductor roots $x,y,z,h,q$ are respectively
$g_0,a_1^2-a_3^2,a_2^2-a_3^2,a_3^2-a_4^2,a_1a_2-a_3a_4$.
Thus (8)--(13) express the complete scalar correction and both full
first normals as infinitesimal perturbations on this same finite
root frame. The leading diagonal in (14), multiplied by
$8K^2D\mathcal C$, has strictly positive residue.

There are no state-normal terms above degree two. To display the
entire hard block before multiplication by $\mathcal C$, put
\[
 C_B=H_3D_vh_s+H_5(1-h_s),
\]
\[
 B_Y=4RT\Delta^2\eta^2c_4DD_v
      +4RT\Delta^4k^6\eta^2c_\alpha D\,v^2
      +2RT\Delta^2\eta c_dD\,v^2.
\]
The entries on $(N,\delta M_B)$ are
\begin{align*}
 A_{11}&=H_2+H_bD_v+\Delta^2C_Bu^2v^2
                         +B_Yh_s^2u^2v^2/D^2,\\
 A_{12}&=-\frac{\Delta^2}{\delta}(H_3D_v-H_5)uv
                         -\frac{B_Y}{\delta D^2}h_suv,\\
 A_{22}&=\frac{\Delta^2}{\delta^2}(H_3D_v+H_5)
                         +\frac{B_Y}{\delta^2D^2}.
\end{align*}
Every apparent negative power of $\delta$ cancels: $\Delta/\delta=Q/K$
is a source unit, and $B_Y/\delta^2$ is bounded infinitesimal even
when $\delta$ is a unit. Hence this exact hard block has residue
$16I_2$. Its denominators are only positive pure-source or
pure-testing units. Together with (14) and the positive-weight
corrections, the full thirteen-root Gram for $\mathcal CG$ is bounded
and has strictly positive residue. Apply the admitted separated
strict-Gram lemma, restore all four packets $\mathcal C\Pi_i$,
and undo the positive source multiplier $\mathcal C\Delta^2\Lambda$.
This proves the entire original parent in the $Z$-unit chart with
bounded $W$, retaining both strict gaps and every original event.

\subsection*{Every remaining testing chart}

\paragraph{Outer angle charts.}
The two preserved outer-angle identities require only
$d_{\rm B}/C,d_{\rm B}/D_{\rm e}$ infinitesimal and the harmonic
mass $\gamma=A[CD_{\rm e}/(C+D_{\rm e})]/d_{\rm B}^2$ bounded
with residue at least $1/4$. These hypotheses are independent of the
size of $e$. Indeed, put $B_c=2c_5-J\eta c_d$. The original source gives
\[
 D_{\rm e}=\frac{aB_c}{Jc_3},\quad
 \frac{d_{\rm B}}{D_{\rm e}}=\frac{2J\eta c_d}{B_c},\quad
 \frac{d_{\rm B}}C=\frac{a\eta c_d}{c_3},
\]
and
\[
 \gamma=\frac{c_4c_3B_c}{c_d^2(2Jc_3+aB_c)}.
\]
Here $\overline B_c=4$, $J+a=T$ has residue one, the denominator
has residue $64$, and the numerator has residue $16$. Thus both
strength ratios are infinitesimal and $\overline\gamma=1/4$.
The exact retained-current identities in
Section~\ref{component:p6cad7544d1b8}, Lemmas 1--2, cover respectively
$1/v$ bounded and $v$ infinitesimal with $1/u$ bounded, with all
homogeneous infinity endpoints. Consequently only $u,v$
infinitesimal remain. The literal axes have a nonnegative core.

\paragraph{State charts.}
Normalize a nonzero state by a coordinate of maximal absolute value.
All its coordinates are then bounded and at least one is a unit. If
$X$ or $Y$ is a unit, the signed core has positive residue
$2(X^2+Y^2)$ and gives a direct separated square; restore all event
squares. Otherwise $X,Y$ are infinitesimal. If $Z$ is a unit, the
preceding section applies, including bounded unit $W$.

The only remaining case has $Z$ infinitesimal and $W$ a unit.
Normalize $W=1$ and put $Q_A=1+uY$, a pure-testing unit. Use
\[
 z=(r_2,B,\Delta r,\Delta r_5,\Delta\eta Y).
\]
In (2), retain the four nonnegative excess packets obtained from
$R/M=1+t/M$, $R/t=1+M/t$, $T/a=1+J/a$, $T/J=1+a/J$,
the complete $\alpha$ packet, and the positive $v^2Y^2$ part of the
last line. The remaining diagonal is
\[
 H_2z_1^2+H_bD_vz_2^2+H_3D_vz_3^2+H_5z_4^2
                   +4RTc_4DD_vz_5^2,
\]
and its complete negative correction is exactly
\[
 -4RTc_dD\frac{\Delta Z}{Q_A}z_5(z_1+z_2)
 -2RTc_dD\frac{\Delta^2\eta Z^2}{Q_A^2}(z_1+z_2)^2.
\]
This is the preserved $W$-unit identity, now on the larger domain:
$\Delta$ is bounded and $Z$ is infinitesimal, so both corrections
are infinitesimal. The diagonal residue is $8I_5$. The separated
strict-Gram lemma and restoration of all six packets prove this
last chart. There is no need for $\Delta$ itself to be infinitesimal.

The zero state is immediate. A complex state splits into its real and
imaginary quadratic evaluations because the parent matrix is real
symmetric. Every normalization and projective clearing factor is
pure testing. This completes the proof of Theorem~\ref{p41b39e2f94e1:thm:bounded-delta}.



\section{\texorpdfstring{The uniform C1 source for every positive \(e\)}{The uniform C1 source for every positive e}}
\label{component:pf5d2f33f694c}

\label{pf5d2f33f694c:sec:uniform-c1-all-e-rederived}

Use the original canonical C1 source and the exact native transport
EN13, with every original coefficient sign and strict gap retained.
Suppose
\[
 k,f,g>0\text{ are infinitesimal},\qquad e=kf/g>0 .
 \tag{UE1}
\]
There is no additional restriction on the size of \(e\), either
source ratio \(g/k,g/f\), or either separately labelled testing field.
Then the entire original parent is tensor SOS at every independent
projective testing pair and every real or complex testing state.

Here are the actual proof inputs and the complete join.  EN13 gives
\(0<a,t<2\), \(r=1-af^2\sim1\), and the original strict
\(M=R-t>0,J=T-a>0\).  Split the positive pure-source \(e\) into
\[
 e\text{ infinitesimal}\quad\hbox{or}\quad \Delta=1/e\text{ bounded}.
 \tag{UE2}
\]
In the first case UJ1--UJ5 is the fresh full-parent theorem.  It
includes both source ratios \(g/k\) bounded and unbounded, both
state charts, both complete state normals and all angular/projective
joins; the actual bounded-ratio product-repair, fresh EZ and fresh
IW proof bodies are its inputs.

In the second case use the actual retained
Section~\ref{component:p41b39e2f94e1}, Theorem \emph{Uniform C1, bounded
reciprocal \(e\)}.  Its hypothesis is exactly UE1 with
\(\Delta\) bounded, and its proof is complete in that body.
For clarity, the scope extension used there is substantive.
After its pure-source substitution
\[
 \delta=\Delta R/(1+k^2),\quad b^2=f^2(1+k^2)/R,
\]
both \(\delta\) and the pure-testing rotated state coordinate \(w\)
are allowed to be bounded units.  In its polynomial ring the weights
\[
 (\delta,k,b,a,t,u,w,v)=(0,2,6,0,0,2,0,16)
 \tag{UE3}
\]
make the full scalar error and both complete first normals have weights
at least \(42,32,26\).  The displayed unlocalized monic conductor,
restricted to \((A_m,B_m,C_m)^2+(D_0^2)\), gives a scalar Gram
whose every correction coefficient has weight at least \(2\);
the first-normal coefficient bounds are \(12,6\).
The exact hard block on its two normalized state normals has
residue \(16I_2\), including every quadratic term.  Its explicit
eleven-soft-root identity supplies a strict positive diagonal.
Thus the entire parent, with all four retained packets restored,
has a bounded strict-residue Gram.  This is a finite polynomial
division proof; neither \(\delta^{-1}\) nor a mixed element is a
lifting coefficient.

The same body proves the complementary \(W=1,Z\) infinitesimal
chart with five exact current roots.  Its two full correction
coefficients contain respectively \(\Delta Z\) and
\(\Delta^2 f^2Z^2\), so vanish in residue even when \(\Delta\)
is a unit.  The original source signs and both strict gap openings
remain intact.

For the outer angle joins, no conclusion from the old infinitesimal-\(e\)
assertion is imported.  Use the exact whole-parent identities UJ3--UJ4.
Their source hypotheses hold on all of UE1.  Indeed the seven native
coquotients have residues \(2,2,2,2,2,2,4\), regardless of \(e\).
Writing \(B_c=2c_5-Jf^2c_d\sim4\), the original weights after
normalizing \(d=1\) satisfy
\[
 C^{-1},D^{-1}\text{ infinitesimal},\qquad
 A\frac{CD}{C+D}
 =\frac{c_4c_3B_c}{c_d^2(aB_c+2Jc_3)}\sim\frac14.
 \tag{UE4}
\]
Here \(a+J=T\sim1\), so the denominator has residue \(64\)
without an individual lower bound on either opening.
Those identities therefore cover every testing chart outside
\(u,v\) infinitesimal, including literal projective infinities.
Literal axes, the zero state, signs and complex states are covered
by their exact full-parent normalization identities.

The two alternatives UE2 exhaust every positive source \(e\).
All scalings are within one labelled field, and every subtracted
packet is restored.  This proves UE1 at all tests, with literal
\(\beta=0\) and every allowed original coefficient boundary.
The two real closed extensions are arbitrarily chosen labelled
factors; no algebraic-disjointness assumption is imposed.



\section{\texorpdfstring{Both uniform C1 source charts on the entire bounded \(Z=1\) state chart}{Both uniform C1 source charts on the entire bounded Z=1 state chart}}
\label{component:pa4f6d1e57f34}
% Fresh whole-parent rederivation, not an unavailable historical witness.

\label{pa4f6d1e57f34:sec:uniform-c1-z-new}

\paragraph{Domain and conclusion.}
Use the original physical uniform C1 source with $d=1$,
$\lambda=n=0$, $k,f,g,e=kf/g$ positive infinitesimals, and
all original nonnegative event weights. Retain
\[
 r=1-af^2,\quad R=1+k^2r,\quad
 T=1+k^2-tk^2g^2,\quad M=R-t>0,\quad J=T-a>0,
 \qquad a,t>0.                                                   \tag{EZ1}
\]
In particular $a,t$ are bounded, $r,R,T$ are positive units,
and $M,J$ may approach zero at arbitrary rates.
For every independent testing $u,v,X,Y$ infinitesimal and
$Z=1$, $W$ arbitrary bounded of either sign, the entire original
parent is a separated tensor SOS. This includes both bounded
and unbounded $b=g/k$, and every permitted $\beta=0$ face.
It does not claim the complementary $W=1,Z$-infinitesimal chart.

\subsubsection*{The exact source and a common set of bounded coefficients}

Use the seven short native coquotients
$c_2,c_\beta,c_3,c_5,c_4,c_\alpha,c_d$ in the retained
Section~\ref{component:p7d4e5164e84a}, now with $\eta=f^2,\zeta=g^2$.
Their residues are $2,2,2,2,2,2,4$.
In the present $d=1$ normalization the exact weights are
\[
\begin{gathered}
 C=\frac{c_3}{af^2c_d},\quad A=\frac{f^2c_4}{c_d},\quad
 m=\frac{k^2c_2}{g^2Mc_d},\quad
 \alpha=\frac{k^4g^2c_\alpha}{c_d},\\
 D+\tfrac12=\frac{c_5}{Jf^2c_d},\qquad
 \beta+\tfrac12=\frac{k^2c_\beta}{g^2tc_d}.
\end{gathered}                                                   \tag{EZ2}
\]
Every use below is of these actual weights. In particular
$\beta\geq0$ is retained, not replaced by a formal source condition.

The following two pure-source coordinate choices are exhaustive:
\[
\begin{array}{c|ccccc}
 &s&\omega&P&G&\kappa=P/G\\ \hline
 b\text{ bounded}&e&b^2&k^2&k^2&1\\
 b\text{ unbounded}&f&1&k^2&g^2&(k/g)^2
\end{array}                                                     \tag{EZ3}
\]
In both rows $f^2=s^2\omega$, $g^2=\omega G$,
$s,\sqrt P,\sqrt G$ are infinitesimal, and $\omega,\kappa$
are positive bounded. In the second row $\kappa$ is infinitesimal.
Here $\beta\geq0$ gives exactly
$t\leq2\kappa c_\beta/c_d$, so $t=\kappa\tau$ with
$\tau>0$ bounded; consequently $M$ is a unit in this row.

Define bounded source coefficients
\[
 c_*=\frac{c_3}{Tc_d},\quad d_*=\frac{c_5}{Tc_d},\quad
 a_*=\frac{c_4}{c_d},\quad \gamma_*=\frac{c_\alpha}{c_d},
 \qquad B_*=\omega\beta t
              =\frac{\kappa c_\beta}{c_d}-\frac{\omega t}{2}.
                                                               \tag{EZ4}
\]
The first four have residue $1/2$; $B_*\geq0$ may vanish.
Use the further table
\[
\begin{array}{c|ccc}
 &\mu&\chi&\psi\\ \hline
 b\text{ bounded}&c_2/(Rc_d)&t&tM\\
 b\text{ unbounded}&\kappa c_2/(Mc_d)&0&0 .
\end{array}                                                     \tag{EZ5}
\]
All entries are bounded. In the first row $\mu$ has residue
$1/2$; in the second $\mu$ is infinitesimal.

\subsubsection*{Exact removal of the source poles}

The original quarter-parent is
\[
\begin{split}
 F={}&mr_2^2+CD_vr_3^2+AD_uD_vY^2
 +(D+\tfrac12)r_5^2+(\alpha+\tfrac12)D_u(vY-1)^2\\
 &+(\beta+\tfrac12)D_v(X+u)^2-\tfrac12D_uD_v,
\end{split}                                                     \tag{EZ6}
\]
where $r_2=vW-X+uvY-u$, $r_3=uW-Y$,
$r_5=uW+uvX-Y-v$, $D_u=1+u^2$, $D_v=1+v^2$.
Put
\[
 z=sW,\qquad q_0=f^2uW^2=\omega uz^2,\qquad
 z_0=f^2uW=s\omega uz.
\]
Retain the following complete nonnegative packets of the original
polynomial:
\[
 \frac{CJ}{T}D_v(r_3-az_0)^2,\qquad
 \frac{(D+1/2)a}{T}(r_5+Jz_0)^2,\qquad
 \beta D_v(X+u-tq_0)^2.                                        \tag{EZ7}
\]
For bounded $b$, also retain
\[
                       \frac{mt}{R}(r_2-Mq_0)^2.              \tag{EZ8}
\]
Each vanishes at the actual source contact. They are literal
positive squares, even when $\beta=0$ or a source opening is small.
The first two removals use $a+J=T$ and the exact identity
\[
 \frac aT[(L+Jq)^2-TJq^2]+\frac JT(L-aq)^2=L^2.
\]
The last removal uses
$r_2^2-(t/R)(r_2-Mq_0)^2
=(M/R)(r_2^2+2tq_0r_2-tMq_0^2)$.
There is no mixed denominator.

Use all three unrestricted normal coordinates
\[
 N_x=X+u-tq_0,\qquad
 N_y=(Y-ruW)/s,\qquad N_v=(v-f^2uW)/s.
                                                               \tag{EZ9}
\]
Only $N_x$ is necessarily bounded. The exact bounds needed later are
$N_x,sN_y,sN_v\in\mathrm I$; they follow from the stated testing
chart and bounded $W$. No boundedness of $N_y,N_v$ is assumed.
Define
\[
 X_*=-u+tq_0+N_x,\qquad
 y_*=ruz+s^2N_y=sY,\qquad
 v_*=\omega uz+N_v=v/s,
\]
\[
 L_2=zv_*+uv_*y_*-X_*-u,\quad
 L_3=T\omega uz-N_y,\quad
 L_5=-N_y+v_*(uX_*-1).
\]
After multiplying the original quarter-parent by the positive
pure-source $\omega$ and removing precisely the scaled packets
EZ7--EZ8, the \emph{entire} remaining polynomial is
\[
\begin{split}
\mathcal F={}&
 \mu(L_2^2+2\chi q_0L_2-\psi q_0^2)
 +c_*D_v(L_3^2-TJ\omega^2u^2z^2)\\
 &+\omega^2a_*D_uD_vy_*^2
 +d_*(L_5^2-aT\omega^2u^2z^2)\\
 &+(\omega^2P^2G\gamma_*+\omega/2)
                         D_u(v_*y_*-1)^2
 +\frac{\omega}{2}D_v(X_*+u)^2\\
 &+B_*D_v\{2q_0(X_*+u)-tq_0^2\}
                         -\frac{\omega}{2}D_uD_v,
 \qquad D_v=1+s^2v_*^2 .
\end{split}                                                     \tag{EZ10}
\]
This is a short full-parent definition of every scalar and normal
coefficient used below. In particular no normal was set to zero
in obtaining it. All its coefficient arrays are bounded through
$a,J,M,t$; the unbounded original weights occur only in the
retained complete squares.

\subsubsection*{The entire scalar is a cubic circuit with finite corrections}

Set $U=u^2,V=z^2$. The polynomial obtained from EZ10 at
$N_x=N_y=N_v=0$ is exactly
\[
 \mathcal F_0=\omega^2\Phi(U,V),\qquad
 \Phi=\gamma_*P^2G(1+U)+UV\,H(U,V),                            \tag{EZ11}
\]
where $H$ has bounded finite polynomial coefficient arrays.
The apparent divisions by $\omega^2$ and $UV$ in this definition
are coefficientwise polynomial extractions: substitution in EZ10
shows those factors before any evaluation. No testing coordinate
is inverted.

At the actual source contact $u^2=P,z^2=G$ all three normal
coordinates and all retained packets vanish. The original contact
value and stationary equations give
\[
 \Phi(P,G)=\Phi_U(P,G)=\Phi_V(P,G)=0.                            \tag{EZ12}
\]
The exact native coquotients in EZ2 independently verify these
identities. Substituting EZ11, and cancelling nonzero pure-source
factors in the source field, gives
\[
 H(P,G)=-\gamma_*P(1+P),\quad
 H_U(P,G)=\gamma_*,\quad
 H_V(P,G)=\gamma_*\kappa(1+P).
\]
Write $\Delta=U-P$, $\Gamma=V-G$.
Finite monic division, first in $U-P$ and then in $V-G$,
therefore gives the complete scalar identity
\[
\begin{split}
 \Phi={}&\gamma_*\{\mathcal M+P\kappa U\Gamma^2\}\\
 &+UV\{\Delta^2H_{20}+\Delta\Gamma H_{11}
                                      +\Gamma^2H_{02}\},\\
 \mathcal M={}&U^2V+\kappa UV^2+P^2G-3PUV .
\end{split}                                                     \tag{EZ13}
\]
All quotients are finite bounded coefficient polynomials; no
inverse of $P,G$ is introduced by their monic construction.
The extra term $\gamma_*P\kappa U\Gamma^2$ is retained as a
nonnegative square.

In the inverse-$b$ row one further coefficient factor is needed:
\[
                         H_{02}=\kappa\widetilde H_{02},
 \qquad\widetilde H_{02}\text{ has bounded arrays}.              \tag{EZ14}
\]
Here is a finite source-division certificate, rather than an
estimate involving independent testing rates. Write
$P=\kappa G$, $t=\kappa\tau$ before substituting the native
coquotients into EZ10. Every denominator is among the pure-source
units $M,T,c_d$ and their factors. At formal $\kappa=0$,
\[
 P=t=\mu=B_*=0,\quad
 c_*=d_*=\frac{1+(1-a)s^2}{2},\quad
 a_*=\frac1{2(1-as^2)} .
\]
The formula gives $H(0,V)=0$ identically in $V$.
After clearing unit denominators, ordinary polynomial division
by $\kappa$ therefore has zero remainder. Apply the two monic
$V-G$ divisions to this quotient. This proves EZ14 with bounded
arrays. In the bounded-$b$ row $\kappa=1$, so it needs no such
factor.

\subsubsection*{Six explicit strict roots and their conductors}

The following identities apply to both rows of EZ3. Put
\[
 \mathsf h=G^{1/6},\quad\mathsf c=\kappa^{1/6},\quad
 \mathsf r=\sqrt[3]u,\quad\mathsf s=\sqrt[3]z,
\]
\[
 (L_1,L_2,L_3)=
 (\mathsf r^4\mathsf s^2,\
   \mathsf c^2\mathsf r^2\mathsf s^4,\
   \mathsf h^6\mathsf c^4),\qquad
 (f_1,f_2,f_3)=
 (\mathsf r^2\mathsf s,\
   \mathsf c\mathsf r\mathsf s^2,\
   \mathsf h^3\mathsf c^2).
\]
Define $R_{2i-1}=f_i(L_1-L_2)$ and
$R_{2i}=f_i(L_1-L_3)$ for $i=1,2,3$.
Then, exactly,
\[
 \mathcal M=\sum_{i=1}^3
       (R_{2i-1}^2-R_{2i-1}R_{2i}+R_{2i}^2).                  \tag{EZ15}
\]
Each pair has the fixed strictly positive Gram
$\left(\begin{smallmatrix}1&-1/2\\-1/2&1\end{smallmatrix}\right)$.
Real cube roots of signed testing quantities are permitted:
$z=sW$ has the separated cube root
$s^{1/3}\sqrt[3]W$. The $L_i=f_i^2$ identity is valid for
either sign of $u,z$. No mixed radical occurs.

For explicit finite allocation define
\[
 J_1=\mathsf r^3R_1+\mathsf c^2\mathsf r\mathsf s^2R_2,\qquad
 J_2=-\mathsf h^6R_1+\mathsf s^6R_2
                              -\mathsf h^3\mathsf s^3R_5.
\]
Direct multiplication gives
\[
\begin{aligned}
 uz\Delta&=J_1,&u^2z\Gamma&=J_2,\\
 \kappa uz\Gamma
   &=\mathsf c^2\mathsf r\mathsf s^2(R_2-R_1)
                                  -\mathsf c^3\mathsf s^3R_3,\\
 UV\Delta^2&=J_1^2,\\
 UV\Delta\Gamma
   &=\mathsf s^6R_2^2-\mathsf h^6R_1^2
             -\mathsf h^3\mathsf s^3R_1(R_5+R_6),\\
 \kappa UV\Gamma^2
   &=\mathsf s^6(2R_1^2-3R_1R_2+R_2^2+R_3^2)
                      +\mathsf h^3\mathsf s^3R_1(R_5-R_6).
\end{aligned}                                                   \tag{EZ16}
\]
All coefficients on the right are infinitesimal: $\mathsf r,
\mathsf s,\mathsf h$ are infinitesimal and $\mathsf c$ is
bounded. It is harmless that $\mathsf c=1$ in the first row.
Equations EZ13--EZ16 therefore put the entire scalar, apart
from its retained positive square, on the six rows $\omega R_i$
with a bounded Gram whose real residue is $\tfrac12$ times the
strict matrix in EZ15.

\subsubsection*{Every complete first normal and all higher terms}

Define $L_x,L_y,L_v$ as half the three derivatives of the
whole polynomial EZ10 at zero normals. Expanding its short
factored form gives exactly
\[
 L_x=\omega uz^2T_x(U,V),\qquad
 L_y=\omega uzT_y(U,V),\qquad
 L_v=\omega uzT_v(U,V),                                       \tag{EZ17}
\]
where each $T_i$ is a polynomial with bounded coefficient arrays.
This supplies a finite construction of every coefficient:
differentiate EZ10, set the three normals to zero, and extract
the displayed monomial. No derivative is dropped.
The exact actual-source contact equations give $T_i(P,G)=0$.
Consequently the monic divisions are
\[
 T_i=(U-P)A_i(U,V)+(V-G)B_i(V).                               \tag{EZ18}
\]
For inverse $b$, each $B_i=\kappa\widetilde B_i$ with
bounded arrays. Indeed the same source substitution used in
EZ14 gives $T_i(0,V)=0$ identically at $\kappa=0$.
Divide the cleared source polynomial by $\kappa$ and then by
the monic $V-G$. This is an exact coefficient construction;
it uses neither $1/u$, $1/z$, nor an evaluated mixed quotient.
For bounded $b$, $\kappa=1$ again suffices.

Combining EZ16--EZ18 proves that both complete state normals
and the complete testing-angle normal have infinitesimal
coefficients on the strict six-row list $\omega R_i$.
The additional factor $z$ in $L_x$ is bounded and infinitesimal.
This argument uses the full first-normal polynomial, not just
its value or linear approximation at contact.

The full normal expansion of EZ10 has a bounded quadratic Gram
whose real residue, in order $(N_x,N_y,N_v)$, is
\[
 \begin{pmatrix}
  \bar\mu+\bar\omega/2&0&0\\
 0&1&1/2\\
 0&1/2&1/2
 \end{pmatrix}\succ0.                                         \tag{EZ19}
\]
For bounded $b$, $\bar\mu=1/2$; for inverse $b$,
$\omega=1$ and $\bar\mu=0$. Both give a strictly positive
first pivot. The other block has pivots $1,1/4$.

Here is also a finite justification for all higher terms, since
the scaled normals $N_y,N_v$ may be unbounded. Expanding EZ10
in the three normals gives degree at most four. A monomial
of normal degree at least three containing $N_x$ has total
$N_y,N_v$ degree at most two. Assign those factors, and any
needed extra $N_x$ factor, to two Gram rows; the remaining
factor contains $N_x\in\mathrm I$. For a higher monomial
without $N_x$, each of the $N_y,N_v$ exponents is at most two,
and its coefficient has a factor $s^2$. If its degree is
three, the unassigned coefficient is a bounded multiple of
$s(sN_y)$ or $s(sN_v)$. If its degree is four, it is a
bounded multiple of $(sN_y)(sN_v)$ (or of one of their
squares). All are infinitesimal by EZ9. These support and
divisibility statements follow directly from the factored
quadratics in EZ10 and are checked on the entire polynomial.
Thus no unbounded scaled normal is hidden in a Gram coefficient.

The scalar, all first normals, and the complete hard block now
give a finite bounded Gram with strictly positive real residue.
The separated strict-Gram lemma supplies tensor SOS.
Restore the whole retained packets EZ7--EZ8 and the positive
square in EZ13, and divide by the pure-source positive $\omega$.
This proves the stated original-parent theorem on both source
rows, without a comparison of source and testing rates.

\paragraph{State and boundary scope.}
The signs of $W$ and of the testing angles were not restricted
in EZ10--EZ19. The polynomial and cube-root identities include
literal zero testing coordinates. If a bounded normalized
$Z=1$ state has $X$ or $Y$ a unit, the original pure-testing
core has positive residue at the double-small angle corner;
its positive source multiple is already a separated square,
and all original event packets are nonnegative squares.
Thus this proof supplies the entire $Z$-unit state chart.
Complex states split into their real and imaginary quadratic
evaluations. The two strong-current outer testing identities
remain separate retained inputs to a future uniform all-tests
join, as does the complementary $W$-unit chart.



\section{\texorpdfstring{The inverse-$b$ uniform C1 source: a complete local conductor}{The inverse-b uniform C1 source: a complete local conductor}}
\label{component:pabbcd0905dcd}
% Fresh derivation, 2026-10-22. No missing historical IW witness is imported.

\label{pabbcd0905dcd:sec:c1-inverse-b-local-new}

\paragraph{Statement and certificate ring.}
Let $K_c,K_t$ be independent real closed extensions of $\mathbb R$,
with bounded rings $B_c,B_t$ and infinitesimal ideals
$\mathfrak m_c,\mathfrak m_t$. The certificate ring is
$\mathcal T=K_c\otimes_{\mathbb R}K_t$, with bounded coefficient ring
$\mathcal B=B_c\otimes_{\mathbb R}B_t$ and zero-residue ideal
$\mathfrak n=\mathfrak m_c\otimes B_t+B_c\otimes\mathfrak m_t$.
We use the finite strict-residue Gram lemma in this ring.
All radicals and inverses below are pure source, pure testing, or
products of separately specified such factors.

Use the actual uniform native C1 stationary source, with
$\lambda=n=0$, $d>0$, and the original nonnegative event coefficients.
Normalize $d=1$ by a positive pure-source scale. Suppose
\[
 k,f,g>0\text{ are infinitesimal},\qquad k/g\in\mathfrak m_c,
 \qquad r=1-af^2,\quad h=1-tg^2,
\]
\[
 R=1+k^2r,\quad T_0=1+k^2h,\quad
 0<a<T_0,\quad 0<t<R,\quad \beta\geq0.                 \tag{IW1}
\]
All original strict gaps remain imposed. Then the complete original
parent is tensor-SOS whenever the independent positive testing angles
$u,v$ are infinitesimal, for every real or complex testing state.
The proof below works in the rotated chart $w=1,z>0$ with $z$ arbitrary;
in particular it includes the previously missing original
$W$-unit, $Z$-infinitesimal chart. It uses no bound on $g/f$ and no
comparison between any testing and source infinitesimals.
The exact sign, zero-axis and state reductions are given at the end.

\subsubsection*{The actual C1 source and bounded source arrays}

For clarity, the six weights used here come from the following
unchanged native formulas. Put $x_n=k^2$, $\eta=f^2$, $\zeta=g^2$,
$y_n=x_n\eta\zeta$, $K=1+x_n$, $G=h+\zeta$, and
\[
 M=R-t,\quad J=T_0-a,\quad E_n=1-x_n\eta h,
 \quad Q_n=1-x_n\zeta r,\quad \theta=t\zeta.
\]
Define
\[
 P_d=1+2x_nh+rx_n^2G+x_n\zeta
       +\eta x_n(K\zeta+x_nh\theta),
\]
\[
 P_4=T_0E_n+x_n(1-r)G,\qquad
 P_\alpha=Rh(1+y_n)+\zeta M E_n.
\]
The actual relative coefficients, now with $d=1$, are
\[
\begin{array}{ll}
 m=\dfrac{Kx_n(r+\eta)h(1+y_n)}{2M\zeta P_d},&
 C=\dfrac{K(r+\eta)E_n}{2(1-r)P_d},\\[4pt]
 A=\dfrac{\eta P_4}{2rP_d},&
 D+\tfrac12=\dfrac{K(1+y_n)(r+\eta)}{2J\eta P_d},\\[4pt]
 \beta+\tfrac12=\dfrac{Kx_n(r+\eta)G}{2\theta P_d},&
 \alpha=\dfrac{x_n^2\zeta(r+\eta)P_\alpha}{2Q_nP_d}.
\end{array}                                                     \tag{IW2}
\]
No auxiliary C2 stationary source is used.
Since $a,t$ are bounded and $k,f,g$ are infinitesimal,
$r,h,P_d,R,T_0$ have residue one. The inequality
$\beta+\tfrac12\geq\tfrac12$ in (IW2) gives
$t=O(k^2/g^2)$. Thus $t$ and $p=t/R$ are infinitesimal, and
$M=R(1-p)$ is a positive unit. The two independent openings
$a$ and $J=T_0-a$ are not capped.

Introduce the exact compact coordinates
\[
 j=\frac{T_0}{g^2R},\quad B=\frac{k^2g^2R}{T_0},\quad
 T=\frac{f^2T_0}{R},\quad p=\frac tR,\quad H=\frac a{T_0},
 \qquad x=jB,\quad y=TB,\quad q_z=jT.                    \tag{IW3}
\]
Here $j$ is unbounded, $0<H<1$, and
\[
 T,x,B,p,\sigma:=1/j,\delta:=j^2B\in\mathfrak m_c,
 \qquad \delta\sim k^2/g^2,\qquad p=O(\delta).
                                                               \tag{IW4}
\]
The quantity $q_z$ can have any positive size. The exact inverse is
\[
\begin{gathered}
 R=\frac{1+x}{1+HTx},\quad T_0=\frac{1+x}{1+Bp},\quad
 r=\frac{1-HT}{1+HTx},\quad h=\frac{j-p}{j(1+Bp)},\\
 f^2=\frac{T(1+Bp)}{1+HTx},\quad
 g^2=\frac{1+HTx}{j(1+Bp)},\quad a=HT_0,\quad t=pR.
\end{gathered}                                                   \tag{IW5}
\]
In particular both strict gaps are retained exactly.

Here is a short rational dictionary sufficient to check every source
estimate. Write $D_+=D+\tfrac12$, $b_+=\beta+\tfrac12$,
$\mu=m(1-p)$, and set
\[
\begin{split}
 P={}&1+B(j+1-p)+BT(1+Hj)\\
 &+B^2T(Hj+jp-p^2+p)+B^2HT^2j(1+Bp),\\
 L={}&1+T(1-H+Bp),\qquad E=1+B(HT+p-1),\\
 F_\alpha={}&j+1-2p+BHTj(1-p)+BTp(j-p).
\end{split}
\]
Substitution of (IW5) in (IW2) gives exactly
\[
\begin{split}
 \mu&=\frac{Bj(j-p)(1+y)(1+Bp)L}{2(1+x)P},\\
 C&=\frac{(1+Bp)L[1+BTp-BTj(1-H)]}{2HT(1+x)P},\\
 A&=\frac{T(1+Bp)[1+BT(2Hj-Hp+H-j+p)+B^2H^2T^2j]}
              {2(1-HT)P},\\
 D_+&=\frac{(1+y)(1+Bp)(1+HTx)L}{2T(1-H)(1+x)P},\\
 b_+&=\frac{Bj(1+Bp)L(j+1-p+HTx)}{2p(1+x)P},\\
 \alpha&=\frac{B^2LF_\alpha}{2EP}.
\end{split}                                                     \tag{IW6}
\]
The denominators $P,E,1+x,1-HT$ are positive bounded units.
After $j=1/\sigma$, $B=\delta\sigma^2$, each of the following
normalized arrays is a rational function with polynomial numerator
and denominator in the bounded variables $\sigma,\delta,T,p,H$;
its denominator has a nonzero real constant at
$\sigma=\delta=T=p=0$, independently of $H$:
\[
 \frac{2\mu}{\delta},\quad\frac{2b_+p}{\delta},\quad
 2TCH,\quad2TD_+(1-H),\quad\frac{2A}{T},\quad
 \frac{2\alpha}{B^2j}\quad\hbox{all have residue }1.
                                                               \tag{IW7}
\]
This finite polynomial criterion proves boundedness and the stated
residues at every relative rate and at both $H$ openings. It also
shows $\beta p\leq b_+p=O(\delta)$, including $\beta=0$.
A useful exact cancellation is
\[
                  K_d:=CH-D_+(1-H)=-\frac{b_+p}{j}.    \tag{IW8}
\]

\subsubsection*{The entire parent and both state normals}

Put $D_u=1+u^2$, $D_v=1+v^2$. The quarter-parent is
\[
\begin{split}
 \mathcal P={}&m(vW-X+uvY-uZ)^2+CD_v(uW-Y)^2+AD_uD_vY^2\\
 &+D_+(uW+uvX-Y-vZ)^2
 +(\alpha+\tfrac12)D_u(vY-Z)^2\\
 &+b_+D_v(X+uZ)^2-\tfrac12D_uD_vZ^2.
\end{split}                                                     \tag{IW9}
\]
On the rotated chart
$w=(W+uY)/D_u=1$, $z=(Z-uX)/D_u>0$, use the two unrestricted normals
\[
 N=\frac{X+uZ-pD_uv}{D_u},\qquad
 M_0=\frac{uW-Y-HD_uvz}{D_u}.
\]
Their exact inverse is
\[
 X=N+pv-uz,\quad Z=z+uN+puv,\quad
 W=1+u(M_0+Hvz),\quad Y=u-M_0-Hvz.                     \tag{IW10}
\]
Only the pure-testing unit $D_u$ is inverted. Retain the positive
packet $mpD_u^2N^2$, and put $G=(\mathcal P-mpD_u^2N^2)/D_u$.
Direct substitution, before any state specialization, gives
\[
\begin{split}
 G={}&\mu D_u[N^2-2vN+(1-p)v^2]
      +CD_uD_v(M_0+Hvz)^2+A D_v(u-M_0-Hvz)^2\\
 &+D_+D_u(M_0-(1-H)vz)^2\\
 &+(\alpha+\tfrac12)
       [u(1-p)v-(1+Hv^2)z-uN-vM_0]^2\\
 &+b_+D_uD_v(N+pv)^2-\tfrac12D_v(z+uN+puv)^2.
\end{split}                                                     \tag{IW11}
\]
In particular $G$ is exactly quadratic in the two normals:
\[
 G=F+2L_NN+2L_MM_0+H_NN^2+H_MM_0^2
                         +2(\alpha+\tfrac12)uvNM_0,
                                                               \tag{IW12}
\]
where $F=G|_{N=M_0=0}$ and
\[
\begin{split}
 h_0&=\mu+b_+\geq\tfrac12,\\
 H_N&=h_0D_u+b_+v^2+\beta u^2v^2+\alpha u^2,\\
 H_M&=D_u(CD_v+D_+)+AD_v+(\alpha+\tfrac12)v^2.
\end{split}                                                     \tag{IW13}
\]
The potentially unbounded $\beta$ is fully retained in (IW11)--(IW13).
There is no omitted higher-normal polynomial.

\subsubsection*{An exact separated scalar circuit}

Define the rows
\[
 U=u/\sqrt x,\quad V=v/\sqrt y,\quad Z_*=z/\sqrt{q_z},
 \quad Q_*=U-VZ_*,\quad B_*=V-UZ_*,\quad R_*=Z_*-UV,
\]
\[
 \Phi=U^2+V^2+Z_*^2+U^2V^2Z_*^2-4UVZ_*.
\]
None of $U,V,Z_*$ is assumed bounded. On the positive chart,
\[
 \Phi=(U-V)^2+Z_*^2(1-UV)^2+2UV(1-Z_*)^2             \tag{IW14}
\]
is a legal separated SOS; for example
$\sqrt{UV}=(xy)^{-1/4}\sqrt{uv}$.

For an explicit finite definition of the actual masses, put
$K_0=CH^2+D_+(1-H)^2$ and
\[
\begin{array}{ll}
 M_a=Ax,&M_b=[\mu(1-p)+b_+p^2]y,\\
 M_c=\alpha q_z,&
 M_{bc}=[K_0+AH^2+2H\alpha+H-\tfrac12]yq_z,\\
 M_{b^2}=b_+p^2y^2,&
 M_{b^2c}=H^2(A+C+\alpha+\tfrac12)y^2q_z,\\
 M_{ab^2c}=CH^2xy^2q_z,&M_{ab^2}=\beta p^2xy^2.
\end{array}
\]
Set
\[
\begin{gathered}
 f_c=M_b,\quad h_c=M_{ab^2},\quad
 k_c=(M_a-M_b+M_{ab^2c}-M_{bc}+h_c)/2,\\
 r_c=M_c-M_b-k_c,\quad q_c=M_a-f_c,\quad
 g_c=M_{ab^2c},\quad d_c=M_{b^2c}-k_c,
 \quad e_c=b_+p^2y^2-g_c.
\end{gathered}
\]
The complete scalar from (IW11) satisfies the exact identity
\[
\begin{split}
 F={}&q_cQ_*^2+r_cR_*^2+d_cV^2Q_*^2+e_cV^2B_*^2\\
 &+f_c\Phi+g_cV^2\Phi+k_cZ_*^2(V^2-1)^2+h_cV^2R_*^2.
\end{split}                                                     \tag{IW15}
\]
The last term is the extra nonnegative square supplied by the actual
$\beta$ event. The checker verifies all scalar monomials of the
original parent and the four remaining source coefficient identities;
(IW15) is not merely an identity at the contact.

The same finite bounded-array criterion as (IW7), with the displayed
mass definitions, gives
\[
 \frac{2f_c}{Tx^2}\sim1,\quad
 \frac{2q_c}{Tx}\sim1,\quad
 \frac{2r_c}{TjB^2}\sim1,\quad
 \frac{k_c}{T^2B^2j}\in\mathfrak m_c,\quad
 \frac{g_c}{f_c}\in\mathfrak m_c.                     \tag{IW16}
\]
Here $\sim1$ means residue one. In particular the three principal
masses $f_c,q_c,r_c$ are positive, and
$|k_c|/r_c$, $|k_c|/(q_cy)$ are infinitesimal.
All normalized mass denominators remain units after
$j=1/\sigma$, $B=\delta\sigma^2$; the companion finite-array
checks test their residues independently of $H$.

For the signed terms use the exact identities
\[
\begin{split}
 (1+U^2)\Phi&=R_*^2+B_*^2+U^2Q_*^2+U^2(1-UVZ_*)^2,\\
 (1+V^2)\Phi&=Q_*^2+R_*^2+V^2B_*^2+V^2(1-UVZ_*)^2,\\
 Z_*(V^2-1)&=-VQ_*-R_*.
\end{split}                                                     \tag{IW17}
\]
The $e_c,g_c$ combination in (IW15) is exactly
\[
 b_+p^2y^2V^2B_*^2-g_c\Phi
       +g_c[Q_*^2+R_*^2+V^2(1-UVZ_*)^2].
\]
A $k_c$ of either sign costs at most
$3|k_c|V^2Q_*^2+2|k_c|R_*^2$, including its occurrence in $d_c$.
The corresponding relative costs are
$3|k_c|v^2/(q_cy)$ and $2|k_c|/r_c$, both in $\mathfrak n$.
Every inequality here abbreviates a finite SOS identity, using
$2A^2+2B^2-(A+B)^2=(A-B)^2$ and nonnegative pure-source factors.
Consequently the scalar supplies strict fixed real fractions of
\[
                       q_cQ_*^2,\qquad r_cR_*^2,
                       \qquad f_c\Phi.                \tag{IW18}
\]
All other nonnegative packets remain available as whole squares.

\subsubsection*{Finite budgets for the complete normal rows}

Call a row free if any fixed positive real multiple of its square can
be subtracted at infinitesimal cost from (IW18), by a finite separated
SOS identity. This definition allows a bounded normal coupling to be
absorbed by a sufficiently large fixed real helper pivot.
Put
\[
 q=vz-Tu,\quad b=v-uz/j,\quad R_z=z-uv/B,\quad
 R_u=v(u^2-x),\quad R_v=u(v^2-y),
\]
\[
 J_u=u(u^2-x),\qquad J_v=v(v^2-y),\qquad
 \widehat q=q/\sqrt T.
\]
Since $Q_*=-q/(T\sqrt x)$, the first mass in (IW18) has a positive
unit coefficient on $\widehat q^2$.

First, $R_u=\sqrt y\,x(-B_*-UR_*)$. Applying the first identity
of (IW17) to $B_*^2$, and the actual $r_c$ pivot to $U^2R_*^2$,
shows that $R_u$ is free. The three relative costs are bounded by
constant multiples of
\[
                         B,\qquad u^2/j,\qquad u^2,
                                                               \tag{IW19}
\]
all infinitesimal. Second,
$R_v=\sqrt x\,y(-Q_*-VR_*)$, so $\sqrt T R_v$ is free;
its two costs are $O(y^2)$ and $O(Tv^2)$.

The same-angle rows have exact decompositions
\[
\begin{split}
 J_v&=\sqrt y\,v(v+\sqrt{y/x}\,u)(V-U)+(y/x)R_u,\\
 J_u&=\sqrt x\,u(u+\sqrt{x/y}\,v)(U-V)+(x/y)R_v.
\end{split}                                                     \tag{IW20}
\]
If $\rho=O(\delta)$ is a nonnegative source factor, then
$\rho J_v$ is free. Indeed $\rho\sqrt{y/f_c}=O(\sqrt\delta)$,
$\sqrt{y/x}=\sqrt{T/j}$ is bounded, and
$\rho y/x$ is infinitesimal. Thus every coefficient after division
by the strict $\sqrt{f_c}(V-U)$ pivot is infinitesimal.
Also $\sqrt T\,TxJ_u$ is free: the two coefficients of its
first summand relative to $\sqrt{f_c}(U-V)$ are bounded multiples of
\[
                  T\sqrt x\,u^2,\qquad\sqrt{T\delta}\,uv,
\]
and its second summand is exactly
$\delta(\sqrt T R_v)$, because $Tx^2/y=j^2B=\delta$.
This last use of $\delta\in\mathfrak m_c$ is the additional
source fact specific to the inverse-$b$ sector.

Finally,
\[
 j\alpha b=-(\alpha/B)R_u-\alpha uR_z.                \tag{IW21}
\]
Here $\alpha/B=O(x)$, and the square of $\alpha uR_z$ costs
$O(x^2u^2)$ relative to $r_cR_*^2$. Hence the entire left side
of (IW21) is free. These identities involve no bound on $z$.

\paragraph{The unrestricted first normal.}
Define the bounded source quantities
\[
 K_z=H(\alpha+\tfrac12)-\tfrac12,\qquad
 K_R=-\mu-\alpha(1-p)-\tfrac12+b_+p+TK_z,
\]
\[
 \rho=b_+p+\beta px\geq0.
\]
Direct differentiation of the entire (IW11), followed by the actual
source contact identity, gives
\[
 L_N=-j\alpha b+(K_R+\beta pv^2)R_u
                   +\rho J_v+K_zuvq.                 \tag{IW22}
\]
For completeness, before imposing the source, the remainder is
$v S_N$, with
\[
 S_N=j\alpha-\mu+b_+p+K_Rx+b_+py+\beta pxy=0
\]
identically under (IW6). Thus no constant or first-normal term is
silently dropped. Equations (IW7) imply
$K_R\sim-\tfrac12$, $\rho=O(\delta)$ and
$\beta p=O(\delta)$. The first three terms in (IW22) are free by
(IW19)--(IW21); the last has coefficient
$K_z\sqrt T\,uv\in\mathfrak n$ on $\widehat q$.
Dividing by $\sqrt{h_0}$ preserves all these bounds because
$h_0\geq\tfrac12$.

\paragraph{The unrestricted second normal.}
Use the exact $K_d$ from (IW8), and put
\[
\begin{split}
 C_q={}&K_dD_u+AH+\alpha+\tfrac12
        +v^2[CHD_u+AH+(\alpha+\tfrac12)H],\\
 c_2={}&T[CH+AH+(\alpha+\tfrac12)H]
                      -A-(\alpha+\tfrac12)(1-p).
\end{split}
\]
The complete second normal has the exact division
\[
 L_M=C_qq+(c_2+TCHu^2)R_v+T(K_d+CHy)J_u.             \tag{IW23}
\]
Before substituting the actual source its only remainder is $uS_M$,
where
\[
\begin{split}
 c_0&=T[K_d+AH+\alpha+\tfrac12]-A,\\
 S_M&=c_0+TK_dx+yc_2+TCHxy=0.
\end{split}
\]
The last equality is an exact identity from (IW6). Furthermore
$c_2$ is bounded (in fact infinitesimal), $TCH$ is a unit, and
$(K_d+CHy)/x$ is bounded with residue $-\tfrac12$.
It follows from the two free rows in (IW20) that
\[
 \sqrt{TH(1-H)}[(c_2+TCHu^2)R_v+T(K_d+CHy)J_u]
\]
is free with bounded coefficients. For the other term,
\[
 \sqrt{TH(1-H)}C_qq
     =T\sqrt{H(1-H)}C_q\widehat q.
\]
Its coefficient is in $\mathfrak n$: multiplying $C_q$ by $T$
clears the only $CH$ pole, and every term then contains either
$T$ or $v^2$. This statement is made on the full expression in
(IW23), before any relative source/testing rate is considered.

\subsubsection*{Strict full Gram and state restoration}

Normalize the hard rows by pure-source factors,
\[
 \widehat N=\sqrt{h_0}\,N,\qquad
 M_0=\sqrt{TH(1-H)}\,\widehat M.
\]
Their full hard matrix is bounded: $b_+/h_0\leq1$,
$\beta/h_0\leq1$, $\alpha/h_0=O(\alpha)$, and the two
source poles of $H_M$ are exactly cleared by $TH(1-H)$.
Equations (IW7), (IW13) give its real residue
\[
                           \operatorname{diag}(1,\tfrac12).
                                                               \tag{IW24}
\]
Its off-diagonal coefficient has the factor
$uv\sqrt{TH(1-H)}/\sqrt{h_0}$ and is infinitesimal.
The $\widehat q$ couplings in (IW22)--(IW23) are infinitesimal.
All remaining complete-normal couplings are bounded coefficients
of free rows. Buy sufficiently large fixed real pivots on the
finitely many free rows, use the elementary real Young square
completion, and apply the strict-residue Gram lemma to the
remaining finite matrix. The strict scalar pivots (IW18) lose only
infinitesimal amounts. This proves $G\in\Sigma\mathcal T^2$.
No state normal has been set to zero in this conclusion.
Restoring $mpD_u^2N^2$, the pure-testing factor $D_u$, and the
positive source scale $d$ proves the assertion for (IW9).

The proof has used $u,v>0$ and $w=1,z>0$, with arbitrary $z$.
The exact sign/state lemma SJ1--SJ10 applies to the same original
core: sign gauges make $u,v$ nonnegative while preserving every
packet; a negative core at positive angles forces the rotated
$w,z$ to have the same nonzero sign. Dividing all state coordinates
by the nonzero pure-testing $w$ gives precisely this chart.
When the core is nonnegative it is itself a pure-testing square,
and every event packet is already positive. Zero angles have the
direct core SOS exits. The zero state is immediate, and a complex
state splits into its real and imaginary quadratic evaluations.
Thus (IW1) is proved for all states at the double-infinitesimal
angles. The existing two-current outer-angle identities give the
separate all-tests join if needed; they are not required as a premise
of the local conductor proved here.



\section{\texorpdfstring{The complete uniform infinitesimal C1 source, at all tests}{The complete uniform infinitesimal C1 source, at all tests}}
\label{component:p845b3cb8b114}
% Fresh complete join of actual proof bodies, 2026-10-22.

\label{p845b3cb8b114:sec:uniform-c1-complete-current-join}

\paragraph{Theorem and unchanged source scope.}
Use the original physical canonical C1 stationary source with
$\lambda=n=0$, $d>0$, every original nonnegative event coefficient,
and every original strict source gap. Suppose
\[
 k,f,g,e=kf/g>0\text{ are infinitesimal},\qquad
 r=1-af^2,\quad R=1+k^2r,\quad T_0=1+k^2-tk^2g^2,
\]
\[
                         a,t>0,\qquad J=T_0-a>0,\quad M=R-t>0.
                                                               \tag{UJ1}
\]
Then the entire original parent is a finite tensor sum of squares
for every independent projective testing pair and every real or
complex testing state. There is no restriction on $b=g/k$, $g/f$,
or any relative infinitesimal rate. In particular the original
$\beta=0$ face is retained.

This is a mathematical join of the actual bounded-$b$ endpoint
product-repair proof, the fresh whole $Z$-chart proof EZ1--EZ19,
the fresh inverse-$b$ local proof IW1--IW24, and the two full-parent
outer-current identities below. No unavailable historical EZ or IW
coefficient witness is used as a premise.

\subsubsection*{Exact original-source bounds for the two outer identities}

Normalize $d=1$ by a positive pure-source scale. The retained
Section~\ref{component:p7d4e5164e84a} prints seven original coquotients.
Their respective residues on (UJ1) are
\[
 (c_2,c_\beta,c_3,c_5,c_4,c_\alpha,c_d)
                    \sim(2,2,2,2,2,2,4).
\]
Their arguments are bounded, so they are positive source units.
The actual weights satisfy
\[
 C=\frac{c_3}{af^2c_d},\qquad
 D=\frac{B_c}{2Jf^2c_d},\qquad A=\frac{f^2c_4}{c_d},
 \qquad B_c=2c_5-Jf^2c_d\sim4.
\]
Consequently
\[
 C^{-1},D^{-1}\in\mathfrak m_c,\qquad
 g_B=\frac{CD}{C+D}>0,\qquad
 \gamma=A g_B=
 \frac{c_4c_3B_c}{c_d^2(aB_c+2Jc_3)}\sim\frac14.
                                                               \tag{UJ2}
\]
This is uniform in both openings $a,J$: the last denominator has
residue $16\cdot4(\bar a+\bar J)=64$, since $a+J=T_0\sim1$.
No positive lower bound on either opening is assumed.
These bounds use the actual original $C,D,A$, not selected
substitute weights. Any other positive event is restored whole.

\subsubsection*{The two exact outer-current identities}

Here are the retained identities with the $d=1$ normalization made
explicit. Write $D_u=1+u^2$, $D_v=1+v^2$, and let the original core be
\[
 \mathcal H=
 2(uW+uvX-Y-vZ)^2+2D_u(vY-Z)^2
 +2D_v(X+uZ)^2-2D_uD_vZ^2.
\]
Set
\[
 R_t=uW-Y,\quad C_t=v(uX-Z),\quad B_t=X+uZ,
 \quad h_B=\frac D{C+D},\quad N_t=R_t+h_BC_t,
\]
\[
 c_u=D_u^{-1/2},\quad t_u=uD_u^{-1/2}.
\]
All angle normalizations are pure testing.

If $1/v$ is bounded, put $\eta_v=\sqrt{D_v}/v$ and
\[
 K_0=2\sqrt{g_B}\,C_t,\qquad
 B_0=\sqrt2\sqrt{D_v/D_u}\,B_t.
\]
Then the entire selected parent has the exact identity
\[
\begin{split}
 4CD_vR_t^2+4D(R_t+C_t)^2+\mathcal H
 ={}&4(C+D)N_t^2+4Cv^2R_t^2+2(R_t+C_t)^2\\
 &+2D_u(vY-Z)^2+(K_0,B_0)Q(K_0,B_0)^T,\\
 Q={}&
 \begin{pmatrix}
 1-\dfrac{c_u^2\eta_v^2}{2g_B}&
          \sqrt{\dfrac1{2g_B}}t_u\eta_v\\
 \sqrt{\dfrac1{2g_B}}t_u\eta_v&1
 \end{pmatrix}.
\end{split}                                                     \tag{UJ3}
\]
Its bounded Gram has residue $I_2$, by (UJ2), since
$c_u,t_u,\eta_v$ are bounded. This proves the selected parent
tensor-SOS for every other angle and every state. Restoring the
other nonnegative original events proves the whole-parent result.

It remains outside the small-angle corner that $v$ is infinitesimal
and $1/u$ is bounded. Put $\epsilon=1/g_B$, $q=\sqrt\epsilon$, and
\[
\begin{gathered}
 K_1=2\sqrt{g_B}\,C_t+\sqrt\epsilon\,Y,\quad
 B_1=B_t-uvY,\quad V_0=vY,\quad
 Y_0=\sqrt\epsilon\sqrt{D_u}\,Y,\\
 L_0=\frac{qc_u}{2}(K_1-c_uY_0)-t_uvB_1+c_uvV_0.
\end{gathered}
\]
Direct multiplication gives
$L_0=\sqrt{D_u}\,v(vY-Z)$ and the exact selected-parent identity
\[
\begin{split}
 &4CD_vR_t^2+4D(R_t+C_t)^2+4AD_uD_vY^2+\mathcal H\\
 &\quad=4(C+D)N_t^2+4Cv^2R_t^2+2(R_t+C_t)^2+\mathcal G,\\
 \mathcal G={}&K_1^2+2D_v(B_1^2+V_0^2)
 +(4\gamma D_v-c_u^2)Y_0^2\\
 &\qquad+2qv(K_1-c_uY_0)V_0-2L_0^2.
\end{split}                                                     \tag{UJ4}
\]
The Gram residue on $(K_1,B_1,V_0,Y_0)$ is
\[
        \operatorname{diag}(1,2,2,4\bar\gamma-\bar c_u^2)
        =\operatorname{diag}(1,2,2,\bar t_u^2)\succ0.
\]
Here $\bar t_u^2>0$ follows from bounded $1/u$. Every correction
contains $q$ or $v$ and has zero residue. Again restore all other
positive original events.

Both identities are whole-parent identities, not state
restrictions. They use bounded normalized homogeneous angle
coordinates, and $\eta_v^2=1+(1/v)^2$; clearing the original
pure-testing homogeneous factors gives the same certificates
at literal Alice or Bob infinity. The exact sign gauges of
SJ1--SJ10 preserve all event packets and reduce finite angle
signs to the nonnegative ones used here. Literal zero angles
have the direct core SOS exits. Thus (UJ3)--(UJ4) cover every
testing chart except $u,v>0$ both infinitesimal.

\subsubsection*{Every remaining source and state chart}

At $u,v$ infinitesimal normalize a nonzero testing state by a
coordinate of largest absolute value. This is a pure-testing
operation; all four coordinates are bounded and at least one
is a unit. If $X$ or $Y$ is a unit, the original pure-testing
core has positive real residue $2(X^2+Y^2)$. It is a positive
testing square, and every original event remains a positive
separated square. We can therefore assume $X,Y$ infinitesimal.

If $Z$ is a unit, normalize $Z=1$. The remaining $W$ is arbitrary
bounded and may have either sign. EZ1--EZ19 applies in each of
its two actual-source coordinate rows:
\[
\begin{array}{c|c|c}
 \text{source}&\text{exact coordinates}&\text{actual proof}\\ \hline
 b=g/k\text{ bounded}& f=eb,\quad g=kb&
                     \mathrm{EZ1\!-\!EZ19,\ bounded\ }b\\
 b=g/k\text{ unbounded}&
        (k,f,g),\quad k/g\in\mathfrak m_c&
                     \mathrm{EZ1\!-\!EZ19,\ inverse\ }b.
\end{array}                                                     \tag{UJ5}
\]
The extra inverse-$b$ source inequality used by EZ is derived from
the actual $\beta\geq0$: $tb^2\leq2c_\beta/c_d$. Thus it adds no
source hypothesis to (UJ1). Its full proof retains both state
normals and the testing-angle normal, including unbounded scaled
normal coordinates through their exact $sN_y,sN_v$ coefficients.

If $Z$ is infinitesimal, $W$ must be a unit. Normalize $W=1$;
then $u,v,X,Y,Z$ are all infinitesimal.
For bounded $b$, use the actual retained
Section~\ref{component:p05180564ed05} proof. Its source variables are exactly
$(e,k,b,a,t)$ with $f=eb$, $g=kb$; its gaps
$R_c=1+k^2(1-ae^2b^2)$ and
$T_0=1+k^2-tk^4b^2$ coincide with (UJ1).
It allows every bounded positive $b$, including infinitesimal $b$,
both strict-gap openings, and all original coefficient faces.
Its whole seven-packet polynomial, complete first normals,
higher-normal support and nine-root product identity are printed
with Section~\ref{component:p7d4e5164e84a} and
Section~\ref{component:p935c9d6318ea}. Its final original-parent
restoration is an exact sum of complete positive packets.
There is no remaining condition on $k^3/e$ or $k/e$ in that proof.

For unbounded $b$, use IW1--IW24. Its exact source hypotheses
$k,f,g,k/g\in\mathfrak m_c$, physical $\beta\geq0$, and the
original gaps follow directly from (UJ1) and this ratio class.
It proves the complete original parent with both unrestricted
state normals for every positive rotated state ratio, so it
includes the present $W$-unit chart. The retained $\beta$ packet,
including its zero face, and both source openings are part of
that fresh proof. The state and sign reduction is the actual
SJ1--SJ10 identity, not a historical chart label.

These cases exhaust every source ratio and every state.
The zero state is immediate. A complex state splits into its
real and imaginary evaluations because the parent matrix is real.
Undo all pure-testing normalizations and the positive source
scale $d$. This proves the theorem on (UJ1) at all independent
projective tests.


